% Generated mechanically from frozen input inputs/p10/ja/formal-spaces.tex.
% Do not edit this generated file; edit the bound locale source instead.

% OK, start here.
%

\setcounter{chapter}{86}
\chapter{形式代数空間}



\phantomsection
\label{formal-spaces-section-phantom}


\section{序論}
\label{formal-spaces-section-introduction}

\noindent
形式スキームは \cite{EGA} で導入された。形式スキームのより一般的な
一つの版は \cite{McQuillan} で，別の版は \cite{Yasuda} で導入された。
形式代数空間は \cite{Kn} で導入された。関連する内容と，それをはるかに
越える多くの内容は \cite{Abbes} および \cite{Fujiwara-Kato} に見いだせる。
本章では，これから用いる形式代数空間の概念を導入する。本章の定義は，
文献に現れる形式スキームおよび形式空間のほとんどのクラスを充満部分圏
として含められるほど一般的である。

\medskip\noindent
これらの代替理論のいくつかと本章の理論との比較も論じるが，理論の論理的な
展開に不要な場合には，必ずしも詳細をすべて与えない。

\medskip\noindent
形式代数空間を導入するほか，いくつかのごく基本的な性質を証明し，
数種類の射についても論じる。








\section{EGA 流の形式スキーム}
\label{formal-spaces-section-formal-schemes-EGA}

\noindent
この節では，\cite{EGA} における形式スキームの構成を概観する。この概念は
代数幾何学において非常に有用であるが，考察すべき正しい概念とは限らない。
むしろ，この理論の枠組みではいくつかの選択がなされており，目的が異なれば
別の選択のほうが適切な場合もある，というべきであろう。実際，文献には，
著者が扱おうとする問題に適合させた，互いに密接に関連する多くの異なる理論が
見られる。それでも，ここで概説する理論の大きな利点の一つは，明確に定まった
幾何学的対象を扱えることである。

\medskip\noindent
始める前に，位相環の層，より一般には位相空間の層に対する層条件に関する
問題を指摘しておく。すなわち，次の大きな圏
\begin{enumerate}
\item 位相空間の圏，
\item 位相群の圏，
\item 位相環の圏，
\item 与えられた位相環上の位相加群の圏
\end{enumerate}
は，集合の圏 $\textit{Sets}$ への自然な忘却関手を備えているが，『層』第
\ref{sheaves-section-algebraic-structures} 節で定義した代数構造の型の例では
ない。したがって，同章で展開した機構を無造作に適用することはできない。
一方，上に挙げた各圏は極限と等化子をもち，集合，群，環，加群への忘却関手は
それらと可換する（『位相』補題 \ref{topology-lemma-limits}，
\ref{topology-lemma-topological-group-limits}，
\ref{topology-lemma-topological-ring-limits}，および
\ref{topology-lemma-topological-module-limits} を参照せよ）。
したがって，『層』定義
\ref{sheaves-definition-sheaf-values-in-category} と同様に層の概念を
定義でき，その基礎にある集合，群，環，または加群の前層は層になる。
重要な違いは，開被覆
$U = \bigcup_{i \in I} U_i$ に対して図式
$$
\xymatrix{
\mathcal{F}(U) \ar[r]
&
\prod\nolimits_{i\in I}
\mathcal{F}(U_i)
\ar@<1ex>[r] \ar@<-1ex>[r]
&
\prod\nolimits_{(i_0, i_1) \in I \times I}
\mathcal{F}(U_{i_0} \cap U_{i_1})
}
$$
が位相空間，位相群，位相環，または位相加群の圏における等化子図式でなければ
ならないことである。言い換えると，最初の写像は
$\mathcal{F}(U)$ を，積位相を備えた
$\prod_{i \in I} \mathcal{F}(U_i)$ の部分空間と同一視しなければならない。

\medskip\noindent
位相空間，位相群，位相環，または位相加群の層 $\mathcal{F}$ の点
$x \in X$ における茎 $\mathcal{F}_x$ は，開近傍全体にわたる対応する圏での
余極限
$$
\mathcal{F}_x = \colim_{x\in U} \mathcal{F}(U)
$$
として定義する。これは集合，群，環，または加群の水準で余極限を取ることと
同じである（『位相』補題 \ref{topology-lemma-colimits}，
\ref{topology-lemma-topological-group-colimits}，
\ref{topology-lemma-topological-ring-colimits}，および
\ref{topology-lemma-topological-module-colimits} を参照せよ）が，さらに
位相を備えている。注意すべきことに，得られる位相はどの圏で作業するかに
依存する。『例』第 \ref{examples-section-colimit-topology} 節を参照せよ。
位相空間，位相群，位相環，または位相加群の前層は層化でき，茎を取る操作は
この操作と可換する。評注
\ref{formal-spaces-remark-sheafification-of-presheaves-in-top} を参照せよ。

\medskip\noindent
$f : X \to Y$ を位相空間の連続写像とする。位相空間，位相群，位相環，
位相加群の層の圏から，$Y$ 上の対応する層の圏への関手 $f_*$ があり，
通常どおり $f_*\mathcal{F}(V) = \mathcal{F}(f^{-1}V)$ と定義される。
（この状況での引戻しについては後で論じる。）
$Y$ 上の位相空間の層 $\mathcal{G}$ と $X$ 上の位相空間の層
$\mathcal{F}$ の間の $f$-写像 $\xi : \mathcal{G} \to \mathcal{F}$ の
概念を，『層』定義 \ref{sheaves-definition-f-map} とまったく同じ方法で
定義する。ただし，すべての開集合 $V \subset Y$ に対して
$\xi_V : \mathcal{G}(V) \to \mathcal{F}(f^{-1}V)$ が連続であるという条件を
付け加える。『層』補題 \ref{sheaves-lemma-f-map} と同様に
$$
\{\mathcal{G}\text{ から }\mathcal{F}\text{ への }f\text{-写像}\} =
\Mor_{\Sh(Y, \textit{Top})}(\mathcal{G}, f_*\mathcal{F})
$$
を得る。位相群，位相環，位相加群の層についても同様である。最後に，
$\xi : \mathcal{G} \to \mathcal{F}$ を上のような $f$-写像とする。
$x \in X$ の像を $y = f(x)$ とすれば，茎の連続写像
$$
\xi_x : \mathcal{G}_y \longrightarrow \mathcal{F}_x
$$
が，『層』定義 \ref{sheaves-definition-composition-f-maps} に続く議論と
まったく同じ方法で定義される。

\medskip\noindent
以上の議論を用いて，「局所位相環付き空間」の圏 $LTRS$ を定義できる。
その対象は，位相空間 $X$ と，茎 $\mathcal{O}_{X,x}$ が（位相を忘れれば）
局所環である位相環の層 $\mathcal{O}_X$ とからなる対
$(X,\mathcal{O}_X)$ である。$LTRS$ の射
$(X, \mathcal{O}_X) \to (Y, \mathcal{O}_Y)$ は対 $(f,f^\sharp)$ であり，
$f : X \to Y$ は位相空間の連続写像，
$f^\sharp : \mathcal{O}_Y \to \mathcal{O}_X$ は $f$-写像であって，
すべての $x \in X$ に対して誘導される写像
$$
f^\sharp_x : \mathcal{O}_{Y, f(x)} \longrightarrow \mathcal{O}_{X, x}
$$
が（位相を忘れた）局所環の局所準同型となるものである。合成は局所環付き空間
の射の合成とまったく同じように定義される。

\medskip\noindent
位相空間 $X$ が準コンパクト開集合からなる基をもつと仮定する。
集合，群，環，またはある環上の加群の層 $\mathcal{F}$ が与えられたとき，
$\mathcal{F}$ には，それぞれ
位相空間，位相群，位相環，位相加群の層の構造を入れることができる。
実際，$U \subset X$ が準コンパクト開集合なら，
$\mathcal{F}(U)$ に離散位相を入れる。$U \subset X$ が任意なら，
準コンパクト開集合による開被覆 $U=\bigcup_{i\in I}U_i$ を選び，
$\mathcal{F}(U)$ に $\prod_{i\in I}\mathcal{F}(U_i)$ から誘導される
位相を入れる（上の議論によれば，そうしなければならない）。
この方法で位相空間，位相群，位相環，または位相加群の層が得られることの
確認は読者に委ねる。位相空間，位相群，位相環，または位相加群の層が
{\it 擬離散}であるとは，すべての準コンパクト開集合 $U \subset X$ に
ついて $\mathcal{F}(U)$ の位相が離散であることをいう。このとき上の構成は
忘却関手の随伴であり，集合の層の圏と位相空間の擬離散層の圏との間の同値を
誘導する（群，環，加群についても同様である）。

\medskip\noindent
Grothendieck と Dieudonn\'e は，まずアフィン形式スキームを定義する。
これは許容位相環 $A$ に対応する。『代数詳論』定義
\ref{more-algebra-definition-topological-ring} を参照せよ。
すなわち，$A$ が与えられたとき，環 $A$ の定義イデアルからなる基本系
$I_\lambda$ を考える。（任意の許容位相環では，すべての定義イデアルの族が
基本系をなす。）各 $\lambda$ に対してスキーム
$\Spec(A/I_\lambda)$ を考えることができる。
$I_\lambda \subset I_\mu$ なら，誘導される射
$$
\Spec(A/I_\mu) \to \Spec(A/I_\lambda)
$$
は，ある $n$ に対して $I_\mu^n\subset I_\lambda$ となるので肥厚である。
これを見る別の方法は，各写像
$$
\Spec(A/I_\lambda) \to \Spec(A)
$$
の像が $A$ の開素イデアルの集合への同相写像であることに注意することである。
これにより，次の定義が動機づけられる：
$$
\text{Spf}(A) = \{\text{開素イデアル }\mathfrak p \subset A\}
$$
ここには $\Spec(A)$ から来る位相を入れる。各 $\lambda$ に対して，
構造層 $\mathcal{O}_{\Spec(A/I_\lambda)}$ を $\text{Spf}(A)$ 上の層と
みなせる。対応する位相環の擬離散層を $\mathcal{O}_\lambda$ とする。
そこで
$$
\mathcal{O}_{\text{Spf}(A)} = \lim \mathcal{O}_\lambda
$$
とおく。ここで極限は位相環の層の圏で取る。対
$(\text{Spf}(A),\mathcal{O}_{\text{Spf}(A)})$ を $A$ の
{\it 形式スペクトル}という。

\medskip\noindent
ここでいくつかの事実を確認すべきである。第一に，茎
$\mathcal{O}_{\text{Spf}(A),x}$ は（位相を忘れれば）局所環である。
第二に，$f\in A$ に対して，対応する開集合
$D(f)\cap\text{Spf}(A)$ について，位相環として
$$
\Gamma(D(f) \cap \text{Spf}(A), \mathcal{O}_{\text{Spf}(A)})
= A_{\{f\}} = \lim (A/I_\lambda)_f
$$
が成り立つ。ここで $I_\lambda$ は上のような定義イデアルの基本系である。
さらに，環 $A_{\{f\}}$ も許容であり，
% P10-E0290
\par\noindent
$(\text{Spf}(A_f), \mathcal{O}_{\text{Spf}(A_{\{f\}})})$ は
\par\noindent
$(D(f) \cap \text{Spf}(A),
\mathcal{O}_{\text{Spf}(A)}|_{D(f) \cap \text{Spf}(A)})$
と同型である。最後に，許容位相環の対 $A,B$ が与えられれば
\begin{equation}
\label{formal-spaces-equation-morphisms-affine-formal-schemes}
\Mor_{LTRS}((\text{Spf}(B), \mathcal{O}_{\text{Spf}(B)}),
(\text{Spf}(A), \mathcal{O}_{\text{Spf}(A)}))
:= \Hom_{cont}(A, B)
\end{equation}
を得る。ここで $LTRS$ は上で定義した「局所位相環付き空間」の圏である。

\medskip\noindent
以上を踏まえ，\cite{EGA} では{\it 形式スキーム}を対
$(\mathfrak X,\mathcal{O}_\mathfrak X)$ として定義する。ここで
$\mathfrak X$ は位相空間であり，$\mathcal{O}_\mathfrak X$ は位相環の
層であって，各点がアフィン形式スキームと（$LTRS$ において）同型な
開近傍をもつ。{\it 形式スキームの射}
$
f : (\mathfrak X, \mathcal{O}_\mathfrak X) \to
(\mathfrak Y, \mathcal{O}_\mathfrak Y)
$
とは，圏 $LTRS$ における射である。

\medskip\noindent
$A$ を離散位相を備えた環とする。このとき $A$ は許容であり，
形式スキーム $\text{Spf}(A)$ は $\Spec(A)$ に等しい。構造層
$\mathcal{O}_{\text{Spf}(A)}$ は $\mathcal{O}_{\Spec(A)}$ に付随する
位相環の擬離散層である。すなわち，その基礎にある環の層は
$\mathcal{O}_{\Spec(A)}$ に等しく，準コンパクト開集合 $U$ 上では
$\mathcal{O}_{\text{Spf}(A)}(U)=\mathcal{O}_{\Spec(A)}(U)$ は
離散位相をもつが，一般の開集合上では必ずしもそうではない。
したがって，任意のアフィンスキームにアフィン形式スキームを対応させられる。
まったく同様に，一般のスキーム $(X,\mathcal{O}_X)$ から，
基礎にある環の層が $\mathcal{O}_X$ である位相環の擬離散層
$\mathcal{O}'_X$ を用いて $(X,\mathcal{O}'_X)$ を構成できる。
この構成は射と両立し，関手
\begin{equation}
\label{formal-spaces-equation-compare-schemes-formal-schemes}
\textit{Schemes} \longrightarrow \textit{Formal Schemes}
\end{equation}
を定める。式
(\ref{formal-spaces-equation-morphisms-affine-formal-schemes}) から直ちに，
この関手が充満忠実であることが従う。

\medskip\noindent
$\mathfrak X$ を形式スキームとする。形式スキームの{\it 大きさ}を
\par\noindent
$\text{size}(\mathfrak X)=\max(\aleph_0,\kappa_1,\kappa_2)$
\par\noindent
で定義する。ここで $\kappa_1$ は $\mathfrak X$ のアフィン形式開集合の
濃度であり，$\kappa_2$ は，アフィン形式開集合
$\mathfrak U\subset\mathfrak X$ にわたる
$\mathcal{O}_\mathfrak X(\mathfrak U)$ の濃度の上限である。

\begin{lemma}
\label{formal-spaces-lemma-fully-faithful}
『集合』補題 \ref{sets-lemma-construct-category} にあるようなスキームの圏
$\Sch_\alpha$ を選ぶ。形式スキーム $\mathfrak X$ に対し，
$$
h_\mathfrak X : (\Sch_\alpha)^{opp} \longrightarrow \textit{Sets},\quad
h_\mathfrak X(S) = \Mor_{\textit{Formal Schemes}}(S, \mathfrak X)
$$
をその点関手とする。このとき，$\mathfrak X$ の大きさが大きすぎなければ
$$
\Mor_{\textit{Formal Schemes}}(\mathfrak X, \mathfrak Y) =
\Mor_{\textit{PSh}(\Sch_\alpha)}(h_\mathfrak X, h_\mathfrak Y)
$$
が成り立つ。
\end{lemma}

\begin{proof}
まず，任意の形式スキーム $\mathfrak X$ に対して
$h_\mathfrak X$ が Zariski 位相の層条件を満たすことに注意する
（『スキーム』定義
\ref{schemes-definition-representable-by-open-immersions} を参照せよ）。
これは，形式スキームの圏における射の局所性から従う。また，形式スキームの
開埋込み $\mathfrak V\to\mathfrak W$ に対し，対応する関手の変換
$h_\mathfrak V\to h_\mathfrak W$ は単射であり，開埋込みで表現可能である
（『スキーム』定義
\ref{schemes-definition-representable-by-open-immersions} を参照せよ）。
形式スキームのアフィン形式スキーム $\mathfrak U_i$ による開被覆
$\mathfrak X=\bigcup\mathfrak U_i$ を選ぶ。このとき関手の族
$h_{\mathfrak U_i}$ は $h_\mathfrak X$ を被覆する
（『スキーム』定義
\ref{schemes-definition-representable-by-open-immersions} を参照せよ）。
最後に
$$
h_{\mathfrak U_i}\times_{h_\mathfrak X}h_{\mathfrak U_j}
=h_{\mathfrak U_i\cap\mathfrak U_j}
$$
である。したがって，写像 $h_\mathfrak X\to h_\mathfrak Y$ を与えることは，
重なり上で一致する写像の族
$h_{\mathfrak U_i}\to h_\mathfrak Y$ を与えることと同値である。
よって補題の写像の全単射性（それぞれ単射性）は，対
$(\mathfrak U_i,\mathfrak Y)$ に対する全単射性（それぞれ単射性）と，
$(\mathfrak U_i\cap\mathfrak U_j,\mathfrak Y)$ に対する単射性
（それぞれ条件なし）へ帰着できる。このようにして，
$\mathfrak X$ がアフィン形式スキームである場合に帰着する。
$\mathfrak X=\text{Spf}(A)$ と書く。ここで用いる位相環は許容である。
さらに，$A$ の定義イデアルからなる基本系
$I_\lambda\subset A$ を選ぶ。

\medskip\noindent
$\mathfrak Y$ 上でも局所化できる。実際，
$\mathfrak V\subset\mathfrak Y$ を開形式部分スキームとし，
$\varphi:h_\mathfrak X\to h_\mathfrak Y$ とする。このとき
$$
h_\mathfrak V\times_{h_\mathfrak Y,\varphi}h_\mathfrak X
\to h_\mathfrak X
$$
は開埋込みで表現可能である。すべての $\lambda$ について
$\Spec(A/I_\lambda)$ へ引き戻すと，開部分スキーム
$U_\lambda\subset\Spec(A/I_\lambda)$ が得られる。しかし
$I_\lambda\subset I_\mu$ なら，射
% P10-E0291
$\Spec(A/I_\lambda)\to\Spec(A/I_\mu)$ は $U_\mu$ を
$U_\lambda$ へ引き戻す。したがって，これらは貼り合わさって
開形式部分スキーム $\mathfrak U\subset\mathfrak X$ を与える。
省略する直接的な議論により
$$
h_\mathfrak U=h_\mathfrak V\times_{h_\mathfrak Y}h_\mathfrak X
$$
が分かる。このように，開被覆
$\mathfrak Y=\bigcup\mathfrak V_j$ と関手の変換
$\varphi:h_\mathfrak X\to h_\mathfrak Y$ が与えられると，
$\mathfrak X$ の対応する開被覆が得られる。
$\mathfrak X$ はアフィンなので，この被覆をアフィン形式部分スキームによる
有限開被覆
$\mathfrak X=\mathfrak U_1\cup\ldots\cup\mathfrak U_n$
で細分できる。言い換えると，各 $i$ に対して，ある $j$ と写像
$\varphi_i:h_{\mathfrak U_i}\to h_{\mathfrak V_j}$ が存在し，図式
$$
\xymatrix{
h_{\mathfrak U_i} \ar[r]_{\varphi_i} \ar[d] & h_{\mathfrak V_j} \ar[d] \\
h_{\mathfrak X} \ar[r]^\varphi & h_\mathfrak Y
}
$$
が可換になる。さらにいくつかの議論を加えると（その詳細は省略する），
$\mathfrak X$ と $\mathfrak Y$ がともにアフィン形式スキームである場合に
補題の全単射性を証明すれば十分であることが従う。

\medskip\noindent
$\mathfrak X$ と $\mathfrak Y$ をアフィン形式スキームとする。
$\mathfrak X=\text{Spf}(A)$ および
$\mathfrak Y=\text{Spf}(B)$ と書く。
$\varphi:h_\mathfrak X\to h_\mathfrak Y$ を関手の変換とする。
$I_\lambda\subset A$ を定義イデアルの基本系とする。標準的な包含射
$i_\lambda:\Spec(A/I_\lambda)\to\mathfrak X$ は，射
$\varphi(i_\lambda):\Spec(A/I_\lambda)\to\mathfrak Y$ へ写される。
式 (\ref{formal-spaces-equation-morphisms-affine-formal-schemes}) により，これは連続写像
$\chi_\lambda:B\to A/I_\lambda$ に対応する。
$\varphi$ は関手の変換なので，$I_\lambda\subset I_\mu$ に対して合成
$B\to A/I_\lambda\to A/I_\mu$ は $\chi_\mu$ に等しい。
言い換えると，環準同型
$$
\chi=\lim\chi_\lambda:B\longrightarrow\lim A/I_\lambda=A
$$
を得る。これは，すべての $\lambda$ に対して $I_\lambda$ の逆像が開で
あるため連続である（$A/I_\lambda$ は離散位相をもち，
$\chi_\lambda$ は連続である）。したがって，式
(\ref{formal-spaces-equation-morphisms-affine-formal-schemes}) により射
$\text{Spf}(\chi):\mathfrak X\to\mathfrak Y$ を得る。
この場合に，この構成が補題の写像の逆であることの確認は省略する。

\medskip\noindent
集合論に関する注意を述べる。与えられたスキームの圏
$\Sch_\alpha$ 上でこの議論を成立させるには，上の証明で用いたすべての
スキームが $\Sch_\alpha$ の対象と同型であることを確かめればよい。
実際，注意深く調べると，上に現れるスキーム
$\Spec(A/I_\lambda)$ が $\Sch_\alpha$ の対象と同型であれば十分である。
このためには，$\mathfrak X$ の大きさが $\Sch_\alpha$ に含まれるある
スキームの大きさ以下であると仮定すれば確かに十分である。
\end{proof}

\begin{lemma}
\label{formal-spaces-lemma-formal-scheme-sheaf-fppf}
\begin{slogan}
形式スキームは fpqc 層である
\end{slogan}
$\mathfrak X$ を形式スキームとする。その点関手
$h_\mathfrak X$（補題 \ref{formal-spaces-lemma-fully-faithful} を参照せよ）は，
fpqc 被覆に対する層条件を満たす。
\end{lemma}

\begin{proof}
『位相』補題 \ref{topologies-lemma-sheaf-property-fpqc} により，
Zariski 被覆の場合と，$R\to S$ が忠実平坦である被覆
$\{\Spec(S)\to\Spec(R)\}$ の場合へ帰着する。
補題 \ref{formal-spaces-lemma-fully-faithful} の証明で，
$h_\mathfrak X$ が Zariski 被覆に対する層条件を満たすことを確認した。

\medskip\noindent
$R\to S$ を忠実平坦な環準同型とする。対応するスキームの射を
$\pi:\Spec(S)\to\Spec(R)$ と書く。これは全射かつ平坦である。
$f:\Spec(S)\to\mathfrak X$ を，
$\Spec(S\otimes_R S)\to\mathfrak X$ という写像として
$f\circ\text{pr}_1=f\circ\text{pr}_2$ を満たす射とする。
『下降』補題 \ref{descent-lemma-equiv-fibre-product} により，基礎集合上の
写像として $f$ は，ある（集合論的な）写像
$g:\Spec(R)\to\mathfrak X$ を用いて $f=g\circ\pi$ と書ける。
『射』補題 \ref{morphisms-lemma-fpqc-quotient-topology} と
$f$ の連続性から，$g$ が連続であることが分かる。

\medskip\noindent
$y\in\Spec(R)$ を取る。$g(y)$ を含むアフィン形式開部分スキーム
$\mathfrak U\subset\mathfrak X$ を選ぶ。
$\mathfrak U=\text{Spf}(A)$ と書く。ここで $A$ はある許容位相環である。
上の議論により，$y\in D(r)\subset g^{-1}(\mathfrak U)$ となる
$r\in R$ を選べる。$\pi^{-1}(D(r))$ 上での $f$ の $\mathfrak U$ への
制限は，式 (\ref{formal-spaces-equation-morphisms-affine-formal-schemes}) により
連続環準同型 $A\to S_r$ に対応する。
$f$ に関する仮定により，誘導される二つの環準同型
$A\to S_r\otimes_{R_r}S_r=(S\otimes_R S)_r$ は等しい。
$R_r\to S_r$ が忠実平坦であることに注意する。
『下降』補題 \ref{descent-lemma-ff-exact} により，二つの射
$S_r\to S_r\otimes_{R_r}S_r$ の等化子は $R_r$ である。
したがって，$A\to S_r$ は写像 $A\to R_r$ を通って一意に分解する。
この写像は $A\to S_r$ と同じ開核をもつので連続でもある。
この写像はさらに，式
(\ref{formal-spaces-equation-morphisms-affine-formal-schemes}) により射
$D(r)\to\mathfrak U$ を与える。

\medskip\noindent
ここまでに何を証明したのだろうか。任意の $y\in\Spec(R)$ に対して，
標準アフィン開集合
$y\in D(r)\subset\Spec(R)$ が存在し，射
$f|_{\pi^{-1}(D(r))}:\pi^{-1}(D(r))\to\mathfrak X$ は，ある射
$D(r)\to\mathfrak X$ を通って一意に分解することを示した。
これらの射が貼り合わさって所望の射
$\Spec(R)\to\mathfrak X$ を与えることの確認は省略する。
\end{proof}

\begin{remark}[McQuillan の変種]
\label{formal-spaces-remark-mcquillan}
McQuillan による形式スキームの構成の変種がある。
\cite{McQuillan} を参照せよ。彼は許容性の条件をわずかに弱めることを
提案している。すなわち，許容位相環とは，線形位相をもち，かつ定義イデアルが
存在する完備な（本章の規約では分離的でもある）位相環 $A$ であったことを
思い出そう。ここで定義イデアルとは，任意の $0$ の近傍が，ある
$n\geq 1$ に対して $I^n$ を含むような開イデアル $I$ のことである。
McQuillan は，ここで{\it 弱許容}位相環と呼ぶものを用いる。
弱許容位相環 $A$ とは，線形位相をもち，かつ{\it 弱定義イデアル}が存在する
完備な（本章の規約では分離的でもある）位相環である。弱定義イデアルとは，
すべての $f\in I$ に対して $n\to\infty$ のとき
$f^n\to 0$ となるような開イデアル $I$ である。許容の場合と同様に，
$I$ が弱定義イデアルで $J\subset A$ が開イデアルなら，
$I\cap J$ は弱定義イデアルである。したがって，弱定義イデアルは
$0$ の開近傍の基本系をなし，この概念に基づくより大きな形式スキームの圏を
定義するために，上とほぼ同じ道筋を進むことができる。
補題 \ref{formal-spaces-lemma-fully-faithful} および
\ref{formal-spaces-lemma-formal-scheme-sheaf-fppf} の類似物は，この設定でも
（同じ証明により）成り立つ。
\end{remark}


\begin{remark}[位相空間の前層の層化]
\label{formal-spaces-remark-sheafification-of-presheaves-in-top}
\begin{reference}
\cite{Gray}
\end{reference}
この注意では，位相空間の前層の層化について簡単に論じる。まったく同じ議論が，
位相アーベル群，位相環，および（与えられた位相環上の）位相加群の前層にも
通用する。これを適切な一般性で行うため，サイト $\mathcal{C}$ 上で考える。
位相空間 $X$ 上の（前）層の場合に関心がある読者は，$\mathcal{C}$ の対象を
$X$ の開集合，$\mathcal{C}$ の射を開集合の包含，$\mathcal{C}$ における
被覆を $X$ における被覆と考えればよい。『サイト』例
\ref{sites-example-site-topological} を参照せよ。
$\mathcal{C}$ 上の位相空間の層の圏を
$\Sh(\mathcal{C},\textit{Top})$ と書き，$\mathcal{C}$ 上の位相空間の
前層の圏を $\textit{PSh}(\mathcal{C},\textit{Top})$ と書く。
$\mathcal{F}$ を $\mathcal{C}$ 上の位相空間の前層とする。層化
$\mathcal{F}^\#$ は，$\Sh(\mathcal{C},\textit{Top})$ の
$\mathcal{G}$ に関して関手的に，式
$$
\Mor_{\textit{PSh}(\mathcal{C}, \textit{Top})}(\mathcal{F}, \mathcal{G})
=
\Mor_{\Sh(\mathcal{C}, \textit{Top})}(\mathcal{F}^\#, \mathcal{G})
$$
を満たすべきである。言い換えると，包含関手
$\Sh(\mathcal{C},\textit{Top})\to
\textit{PSh}(\mathcal{C},\textit{Top})$ の左随伴を構成しようとしている。
まず，$\Sh(\mathcal{C},\textit{Top})$ は極限をもち，包含関手はそれらと
可換であると主張する。実際，圏 $\mathcal{I}$ と
$\Sh(\mathcal{C},\textit{Top})$ への関手
$i\mapsto\mathcal{G}_i$ が与えられたとき，単に
$$
(\lim \mathcal{G}_i)(U) = \lim \mathcal{G}_i(U)
$$
と定義する。ここで極限は位相空間の圏で取る（『位相』補題
\ref{topology-lemma-limits}）。極限どうしは可換するので（『圏』補題
\ref{categories-lemma-colimits-commute}），とりわけ層の公理で用いる演算である
積および等化子とも可換し，これは層を定義する。最後に，前層の射
$\mathcal{F}\to\lim\mathcal{G}_i$ を与えることは，明らかに両立する射の系
$\mathcal{F}\to\mathcal{G}_i$ を与えることと同じである。言い換えると，対象
$\lim\mathcal{G}_i$ は位相空間の前層の圏における極限であり，したがって特に
位相空間の層の圏における極限でもある。
第二の主張は，$\mathcal{G}$ が $\Sh(\mathcal{C},\textit{Top})$ の対象で
あるような前層の任意の射 $\mathcal{F}\to\mathcal{G}$ が，大きさの有界な
部分層 $\mathcal{G}'\subset\mathcal{G}$ を通って分解するということである。
ここで，位相空間の層 $\mathcal{H}$ の{\it 大きさ} $|\mathcal{H}|$ を基数
$\sup_{U\in\Ob(\mathcal{C})}|\mathcal{H}(U)|$ と定義する。
この主張を証明するため，
$$
\mathcal{G}'(U) =
\left\{
\quad
s \in \mathcal{G}(U)
\quad \middle| \quad
\begin{matrix}
\text{次の被覆が存在する： }\{U_i \to U\}_{i \in I} \\
\text{しかも }
s|_{U_i} \in \Im(\mathcal{F}(U_i) \to \mathcal{G}(U_i))
\end{matrix}
\quad
\right\}
$$
とおく。$\mathcal{G}'(U)$ には誘導位相を入れる。このとき
$\mathcal{G}'$ は位相空間の層であり（詳細は省略する），
$\mathcal{G}'\to\mathcal{G}$ は，与えられた写像
$\mathcal{F}\to\mathcal{G}$ がそれを通って分解する射である。さらに，
$\mathcal{G}'$ の大きさは，$\mathcal{C}$ と前層 $\mathcal{F}$ のみに依存する
ある基数 $\kappa$ で抑えられる（ヒント：本章の規約により，
$\mathcal{C}$ における被覆が集合をなすことを用いよ）。以上を合わせると，
『圏』定理 \ref{categories-theorem-adjoint-functor} の仮定が満たされ，
層から前層への包含関手の左随伴として層化が得られる。
最後に，サイト $\mathcal{C}$ の点 $p$ が関手
$u:\mathcal{C}\to\textit{Sets}$ によって与えられているとする。
『サイト』定義 \ref{sites-definition-point} を参照せよ。位相空間 $M$ に対し，
規則
$$
\begin{aligned}
U &\mapsto \text{Map}(u(U), M) \\
  &= \prod\nolimits_{x \in u(U)} M
\end{aligned}
$$
で定まり積位相を備えた前層は，位相空間の層である。したがって，\par\noindent
『サイト』補題 \ref{sites-lemma-point-pushforward-sheaf} の証明と同じ議論から
$\mathcal{F}_p=\mathcal{F}^\#_p$ が分かる。
言い換えると，層化は点における茎を取ることと可換する。
\end{remark}




\section{規約と記法}
\label{formal-spaces-section-conventions}

\noindent
以下の規約は，『空間の性質』節
\ref{spaces-properties-section-conventions} の規約と同様である。
したがって，以下ではすべてのスキームが大 fppf サイト
$\Sch_{fppf}$ に含まれることを常に仮定する。また，考えるすべての環 $A$ は，
$\Spec(A)$ がこの大サイトの対象と（同型で）あるという性質をもつものとする。
位相環 $A$ については，すべての離散商がこの性質をもつことだけを仮定する
（ただし通常はより多くを仮定する。注意 \ref{formal-spaces-remark-set-theoretic} と比較せよ）。

\medskip\noindent
$S$ をスキームとし，$X$ を $S$ 上の「空間」，すなわち
$(\Sch/S)_{fppf}$ 上の層とする。本章では，$X\times X$ ではなく，
$(\Sch/S)_{fppf}$ 上の層の圏における $X$ 自身との積を
$X\times_S X$ と書く。さらに，$X$ と $Y$ が「空間」であるとき，
「$f:X\to Y$ を射とする」とは，$f$ が関手の自然変換，すなわち
$(\Sch/S)_{fppf}$ 上の層の写像であることを意味する。同様に，
$U$ が $S$ 上のスキームで $X$ が $S$ 上の「空間」であるとき，
「$f:U\to X$ を射とする」または「$g:X\to U$ を射とする」とは，
$f$ または $g$ が層の写像 $h_U\to X$ または $X\to h_U$ であることを
意味する。ここで $h_U$ は『圏』例
\ref{categories-example-hom-functor} におけるものとする。






\section{位相環と加群}
\label{formal-spaces-section-topological-module}

\noindent
本節は『代数詳論』節
\ref{more-algebra-section-topological-ring} の続きである。
$R$ を位相環，$M$ を線形位相を備えた $R$-加群とする。
「{\it $M_\lambda$ を開部分加群の基本系とする}」というとき，
各 $M_\lambda$ が開部分加群であり，$0$ の任意の近傍が
$M_\lambda$ のいずれかを含むことを意味する。言い換えると，
$M_\lambda$ は部分加群からなる $M$ における $0$ の近傍の基本系である。
同様に，$R$ が線形位相環であるとき，
「{\it $I_\lambda$ を開イデアルの基本系とする}」とは，
$I_\lambda$ がイデアルからなる $R$ における $0$ の近傍の基本系であることを
意味する。

\begin{example}
\label{formal-spaces-example-what-does-it-mean}
$R$ を線形位相環，$M$ を線形位相を備えた $R$-加群とする。
$I_\lambda$ を $R$ の開イデアルの基本系，$M_\mu$ を $M$ の
開部分加群の基本系とする。$+:M\times M\to M$ の連続性は自動的であり，
$R\times M\to M$ の連続性は
$$
\forall f, x, \mu\ \exists \lambda, \nu,\ (f + I_\lambda)(x + M_\nu)
\subset fx + M_\mu
$$
を意味する。$M_\nu\subset M_\mu$ なら
$fM_\nu+I_\lambda M_\nu\subset M_\mu$ なので，この条件は
$$
\forall x, \mu\ \exists \lambda\ I_\lambda x \subset M_\mu
$$
と同値である。しかし，$\mu$ が与えられたとき
$I_\lambda M\subset M_\mu$ となる $\lambda$ が存在するとは限らない。
例えば，$t$-進位相を備えた $R=k[[t]]$ と，開部分加群の基本系
$$
M_m = \bigoplus\nolimits_{n \in \mathbf{N}} t^{nm}R
$$
を備えた $M=\bigoplus_{n\in\mathbf{N}}R$ を考える。
任意の $x\in M$ は非零成分を有限個しかもたないので，$m$ と $x$ が
与えられれば $t^kx\in M_m$ となる $k$ が存在する。したがって $M$ は
線形位相を備えた $R$-加群であるが，$m$ が与えられたとき
$t^kM\subset M_m$ となる $k$ が存在するとは限らない。一方，
$R\to S$ が線形位相環の連続写像なら，対応する主張は成り立つ。すなわち，
任意の開イデアル $J\subset S$ に対して $IS\subset J$ となる開イデアル
$I\subset R$ が存在する（写像 $R\to S$ の連続性から容易に従う）。
\end{example}

\begin{lemma}
\label{formal-spaces-lemma-closed}
$R$ を位相環とする。$M$ を線形位相を備えた $R$-加群とし，
$M_\lambda$（$\lambda\in\Lambda$）を開部分加群の基本系とする。
$N\subset M$ を部分加群とする。このとき $N$ の閉包は
$\bigcap_{\lambda\in\Lambda}(N+M_\lambda)$ である。
\end{lemma}

\begin{proof}
各 $N+M_\lambda$ は開なので閉でもある。したがって，その共通部分は閉である。
% P10-E0292
$x\in M$ が $N$ の閉包に属さないなら，ある $\lambda$ に対して
$(x+M_\lambda)\cap N=0$ である。したがって
$x\not\in N+M_\lambda$ である。これで補題が証明された。
\end{proof}

\noindent
特に断らない限り，部分加群と商加群には誘導位相を入れる。
$M$ を位相環 $R$ 上の線形位相加群とし，
$0\to N\to M\to Q\to 0$ を $R$-加群の短完全列とする。
$M_\lambda$ が $M$ の開部分加群の基本系なら，
$N\cap M_\lambda$ は $N$ の開部分加群の基本系である。
$\pi:M\to Q$ が商写像なら，$\pi(M_\lambda)$ は $Q$ の
開部分加群の基本系である。特に，これらの誘導位相は線形位相である。

\begin{lemma}
\label{formal-spaces-lemma-closure}
$R$ を位相環とする。$M$ を線形位相を備えた $R$-加群とし，
$N\subset M$ を部分加群とする。このとき
\begin{enumerate}
\item $0\to N^\wedge\to M^\wedge\to(M/N)^\wedge$ は完全であり，
\item $N^\wedge$ は $N\to M^\wedge$ の像の閉包である。
\end{enumerate}
\end{lemma}

\begin{proof}
$M_\lambda$（$\lambda\in\Lambda$）を開部分加群の基本系とする。
$N\cap M_\lambda$ は $N$ の開部分加群の基本系であり，
% P10-E0293
$M_\lambda+N/N$ は $M/N$ の開部分加群の基本系である。したがって，
(1) は完全列
$$
0 \to N/N \cap M_\lambda \to M/M_\lambda \to M/(M_\lambda + N) \to 0
$$
の完全性と，極限を取る操作が極限と可換することから従う。
第二の主張は，これと，$N\to N^\wedge$ の像が稠密であること，および
$M^\wedge\to(M/N)^\wedge$ の核が閉であることから従う。
\end{proof}

\begin{lemma}
\label{formal-spaces-lemma-quotient-by-closed}
$R$ を位相環とする。$M$ を完備な線形位相 $R$-加群とし，
$N\subset M$ を閉部分加群とする。$M$ が可算な $0$ の近傍の基本系を
もつなら，$M/N$ は完備であり，写像 $M\to M/N$ は開である。
\end{lemma}

\begin{proof}
$M_n$（$n\in\mathbf{N}$）を $M$ の開部分加群の基本系とする。
すべての $n$ に対して $M_{n+1}\subset M_n$ と仮定してよい。
系 $(M_n+N)/N$ は $M/N$ の基本系である。したがって
$M/N=\lim M/(M_n+N)$ を示せばよい。短完全列
$$
0 \to N/N \cap M_n \to M/M_n \to M/(M_n + N) \to 0
$$
を考える。系 $\{N/N\cap M_n\}$ の遷移写像は全射なので，
『代数』補題 \ref{algebra-lemma-ML-exact-sequence} により，
$M$ の完備性から得られる $M=\lim M/M_n$ は
$\lim M/(M_n+N)$ へ全射する。$N$ は閉なので，
$M\to\lim M/(M_n+N)$ の核は $N$ である
（補題 \ref{formal-spaces-lemma-closed} を参照せよ）。最後に，商位相の定義により
$M\to M/N$ は開である。
\end{proof}

\begin{lemma}
\label{formal-spaces-lemma-ses}
\begin{reference}
\cite[Theorem 8.1]{Ma}
\end{reference}
$R$ を位相環とする。$M$ を線形位相を備えた $R$-加群とし，
$N\subset M$ を部分加群とする。$M$ が可算な $0$ の近傍の基本系を
もつと仮定する。このとき
\begin{enumerate}
\item $0\to N^\wedge\to M^\wedge\to(M/N)^\wedge\to0$ は完全であり，
\item $N^\wedge$ は $N\to M^\wedge$ の像の閉包であり，
\item $M^\wedge\to(M/N)^\wedge$ は開である。
\end{enumerate}
\end{lemma}

\begin{proof}
補題 \ref{formal-spaces-lemma-closure} により，
$0\to N^\wedge\to M^\wedge\to(M/N)^\wedge$ は完全であり，
(2) も成り立つ。これにより標準写像
$c:M^\wedge/N^\wedge\to(M/N)^\wedge$ が得られる。
補題 \ref{formal-spaces-lemma-quotient-by-closed} により，加群
$M^\wedge/N^\wedge$ は完備であり，
$M^\wedge\to M^\wedge/N^\wedge$ は開である。完備化の普遍性により
標準写像 $b:(M/N)^\wedge\to M^\wedge/N^\wedge$ を得る。
$b$ と $c$ は稠密部分集合上で互いに逆なので，互いに逆である。
\end{proof}

\begin{lemma}
\label{formal-spaces-lemma-completion-adic-star}
$R$ を位相環とし，$M$ を位相 $R$-加群とする。
$I\subset R$ を有限生成イデアルとする。$M$ が，その位相が $I$-進である
開部分加群をもつと仮定する。このとき $M^\wedge$ は，その位相が $I$-進である
開部分加群をもち，すべての $n\geq1$ に対して
$M^\wedge/I^nM^\wedge=M/I^nM$ である。
\end{lemma}

\begin{proof}
$M'\subset M$ を，その位相が $I$-進である開部分加群とする。このとき
$\{I^nM'\}_{n\geq1}$ は $M$ の開部分加群の基本系である。したがって
$M^\wedge=\lim M/I^nM'$ は
$(M')^\wedge=\lim M'/I^nM'$ を開部分加群として含み，
『代数』補題 \ref{algebra-lemma-hathat-finitely-generated} により
$(M')^\wedge$ の位相は $I$-進である。
$I$ は有限生成なので $I^n$ も有限生成であり，その生成元を
$f_1,\ldots,f_r$ とする。全射
$(f_1,\ldots,f_r):M^{\oplus r}\to I^nM$ は，上で述べた $M$ の
位相の記述により連続かつ開であることに注意する。この全射と短完全列
$0\to I^nM\to M\to M/I^nM\to0$ に補題 \ref{formal-spaces-lemma-ses} を適用すると，
$$
(f_1, \ldots, f_r) :
(M^\wedge)^{\oplus r} \longrightarrow M^\wedge
$$
は全射 $M^\wedge\to M/I^nM$ の核へ全射する。
$f_1,\ldots,f_r$ は $I^n$ を生成するので結論を得る。
\end{proof}

\begin{definition}
\label{formal-spaces-definition-toplogy-tensor-product}
\begin{reference}
\cite[0, Section 7.7]{EGA} と比較せよ。
\end{reference}
$R$ を位相環とし，$M$ と $N$ を線形位相を備えた $R$-加群とする。
$M$ と $N$ の{\it テンソル積}とは，（通常の）テンソル積
$M\otimes_RN$ に，
$$
\Im(M_\mu \otimes_R N + M \otimes_R N_\nu \longrightarrow M \otimes_R N)
$$
を開部分加群の基本系と宣言して定まる線形位相を入れたものである。ここで
$M_\mu\subset M$ と $N_\nu\subset N$ は，それぞれ $M$ と $N$ の
開部分加群の基本系を動く。{\it 完備化テンソル積}
$$
M \widehat{\otimes}_R N =
\lim M \otimes_R N/(M_\mu \otimes_R N + M \otimes_R N_\nu) =
\lim M/M_\mu \otimes_R N/N_\nu
$$
とは，このテンソル積の完備化である。
\end{definition}

\noindent
テンソル積または完備化テンソル積の構成において，$R$ の位相は関与しないことに
注意する。$R\to A$ と $R\to B$ が線形位相環の連続写像なら，上の構成により
テンソル積 $A\otimes_RB$ と完備化テンソル積
$A\widehat{\otimes}_RB$ が得られる。

\medskip\noindent
注意 \ref{formal-spaces-remark-mcquillan} で導入した概念をここに記録しておく。

\begin{definition}
\label{formal-spaces-definition-weakly-admissible}
$A$ を線形位相環とする。
\begin{enumerate}
\item 元 $f\in A$ が{\it 位相的冪零}であるとは，$n\to\infty$ のとき
$f^n\to0$ となることをいう。
\item $A$ の{\it 弱定義イデアル}とは，位相的冪零元のみからなる
開イデアル $I\subset A$ のことである。
\item $A$ が弱定義イデアルをもつとき，$A$ は{\it 弱前許容}であるという。
\item $A$ が弱前許容かつ完備であるとき，$A$ は{\it 弱許容}であるという
\footnote{本章の規約では，これは分離的であることも含む。}。
\end{enumerate}
\end{definition}

\noindent
線形位相環 $A$ の弱定義イデアル $I$ と開イデアル $J$ が与えられると，
$I\cap J$ は弱定義イデアルである。したがって，一つでも弱定義イデアルが
存在すれば，弱定義イデアルからなる開イデアルの基本系が存在する。特に，
$A$ が弱許容位相環なら，弱定義イデアルの基本系 $\{I_\lambda\}$ を用いて
$A=\lim A/I_\lambda$ である。

\begin{lemma}
\label{formal-spaces-lemma-weakly-admissible-henselian}
$A$ を弱許容位相環とし，$I\subset A$ を弱定義イデアルとする。
このとき $(A,I)$ は Hensel 対である。
\end{lemma}

\begin{proof}
$A\to A'$ を \'etale 環準同型とし，$\sigma:A'\to A/I$ を
$A$-代数準同型とする。『代数詳論』補題
\ref{more-algebra-lemma-characterize-henselian-pair} により，
$\sigma$ を $A$-代数準同型 $A'\to A$ へ持ち上げれば十分である。
これを行うには，$A$ が完備なので，すべての開イデアル $J\subset I$ に対して
$\sigma$ を持ち上げる一意な $A$-代数準同型 $A'\to A/J$ を見つければ十分である。
$I$ は弱定義イデアルなので，イデアル $I/J$ は局所冪零である。
したがって『代数詳論』補題
\ref{more-algebra-lemma-locally-nilpotent-henselian} により結論を得る。
\end{proof}

\begin{lemma}
\label{formal-spaces-lemma-topologically-nilpotent}
$B$ を線形位相環とする。$B$ の位相的冪零元全体は，$B$ の閉根基イデアルである。
$\varphi:A\to B$ を線形位相環の連続写像とする。
\begin{enumerate}
\item $f\in A$ が位相的冪零なら，$\varphi(f)$ は位相的冪零である。
\item $I\subset A$ が位相的冪零元のみからなるなら，
$\varphi(I)B$ の閉包は位相的冪零元のみからなる。
\end{enumerate}
\end{lemma}

\begin{proof}
$\mathfrak b\subset B$ を位相的冪零元全体とする。
$\mathfrak b$ が根基イデアルであることの証明は省略する
（定義に関するよい演習である）。$g$ を $\mathfrak b$ の閉包の元とする。
$g$ が位相的冪零であることを示そう。$J\subset B$ を開イデアルとする。
ある $e\geq1$ に対して $g^e\in J$ となることを示せばよい。
補題 \ref{formal-spaces-lemma-closed} により $g\in\mathfrak b+J$ である。したがって，
ある $f\in\mathfrak b$ と $h\in J$ に対して $g=f+h$ と書ける。
$f^m\in J$ となる $m\geq1$ を選ぶ。このとき所望のとおり
$g^{m+1}\in J$ である。

\medskip\noindent
$\varphi:A\to B$ を補題の主張におけるものとする。
主張 (1) は明らかであり，主張 (2) はこれと $\mathfrak b$ が
閉イデアルであることから従う。
\end{proof}

\begin{lemma}
\label{formal-spaces-lemma-closure-image-ideal}
$A\to B$ を線形位相環の連続写像とし，$I\subset A$ をイデアルとする。
$IB$ の閉包は $B\to B\widehat{\otimes}_A A/I$ の核である。
\end{lemma}

\begin{proof}
% P10-E0294
$J_\mu$ を $B$ の開イデアルの基本系とする。補題 \ref{formal-spaces-lemma-closed} により，
$IB$ の閉包は $\bigcap(IB+J_\lambda)$ である。
$I_\mu$ を $A$ の開イデアルの基本系とする。このとき
$$
B \widehat{\otimes}_A A/I = \lim (B/J_\lambda \otimes_A A/(I_\mu + I)) =
\lim B/(J_\lambda + I_\mu B + I B)
$$
である。$A\to B$ は連続なので，すべての $\lambda$ に対して
$I_\mu B\subset J_\lambda$ となる $\mu$ が存在する。例
\ref{formal-spaces-example-what-does-it-mean} の議論を参照せよ。したがって極限は
$\lim B/(J_\lambda+IB)$ と書け，結論は明らかである。
\end{proof}

\begin{lemma}
\label{formal-spaces-lemma-completed-tensor-product}
$B\to A$ と $B\to C$ を線形位相環の連続準同型とする。
\begin{enumerate}
\item $A$ と $C$ が弱前許容なら，$A\widehat{\otimes}_BC$ は弱許容である。
\item $A$ と $C$ が前許容なら，$A\widehat{\otimes}_BC$ は許容である。
\item $A$ と $C$ が可算な開イデアルの基本系をもつなら，
$A\widehat{\otimes}_BC$ も可算な開イデアルの基本系をもつ。
\item $A$ と $C$ が pre-adic で有限生成の定義イデアルをもつなら，
$A\widehat{\otimes}_BC$ は adic で，有限生成の定義イデアルをもつ。
\item $A$ と $C$ が pre-adic Noether 環で，$\mathfrak a\subset A$ と
$\mathfrak b\subset B$ が位相的冪零元のイデアルであり，
$B/\mathfrak b\to A/\mathfrak a$ が有限型なら，
$A\widehat{\otimes}_BC$ は adic Noether 環である。
\end{enumerate}
\end{lemma}

\begin{proof}
$I_\lambda\subset A$（$\lambda\in\Lambda$）と
$J_\mu\subset C$（$\mu\in M$）を開イデアルの基本系とすると，定義により
$$
A \widehat{\otimes}_B C =
\lim_{\lambda, \mu} A/I_\lambda \otimes_B C/J_\mu
$$
であり，これには極限位相を入れる。したがって，写像
$A\widehat{\otimes}_BC\to A/I_\lambda\otimes_BC/J_\mu$ の核
$K_{\lambda,\mu}$ が開イデアルの基本系をなす。
$K_{\lambda,\mu}$ はイデアル
$I_\lambda(A\widehat{\otimes}_BC)+J_\mu(A\widehat{\otimes}_BC)$ の
閉包であることに注意する。最後に，稠密な像をもつ環準同型
$\tau:A\otimes_BC\to A\widehat{\otimes}_BC$ がある。

\medskip\noindent
(1) の証明。$I_\lambda$ と $J_\mu$ が位相的冪零元のみからなるなら，
補題 \ref{formal-spaces-lemma-topologically-nilpotent} により
$K_{\lambda,\mu}$ もそうである。したがって定義により
$A\widehat{\otimes}_BC$ は弱許容である。

\medskip\noindent
(2) の証明。ある $\lambda_0$ と $\mu_0$ に対して，イデアル
$I=I_{\lambda_0}\subset A$ と $J_{\mu_0}\subset C$ が
定義イデアルであると仮定する。すると，すべての $\lambda$ に対して
$I^n\subset I_\lambda$ となる $n$ が存在し，すべての $\mu$ に対して
$J^m\subset J_\mu$ となる $m$ が存在する。このとき
$$
\left(I(A \widehat{\otimes}_B C) + J(A \widehat{\otimes}_B C)\right)^{n + m}
\subset
I_\lambda(A \widehat{\otimes}_B C) + J_\mu(A \widehat{\otimes}_B C)
$$
である。したがって，開イデアル $K=K_{\lambda_0,\mu_0}$ は
$K^{n+m}\subset K_{\lambda,\mu}$ を満たす。よって $K$ は
$A\widehat{\otimes}_BC$ の定義イデアルであり，定義により
$A\widehat{\otimes}_BC$ は許容である。

\medskip\noindent
(3) の証明。$\Lambda$ と $M$ が可算なら，$\Lambda\times M$ も可算である。

\medskip\noindent
(4) の証明。$\Lambda=\mathbf{N}$，$M=\mathbf{N}$ と仮定し，
$I_n=I^n$ および $J_n=J^n$ となる有限生成イデアル
$I\subset A$ と $J\subset C$ があるとする。このとき
$$
I(A \widehat{\otimes}_B C) + J(A \widehat{\otimes}_B C)
$$
は有限生成イデアルであり，$A\widehat{\otimes}_BC$ がこのイデアルに関する
$A\otimes_BC$ の完備化であることは容易に分かる。したがって (4) は
『代数』補題 \ref{algebra-lemma-hathat-finitely-generated} から従う。

\medskip\noindent
(5) の証明。$\mathfrak c\subset C$ を位相的冪零元のイデアルとする。
$A$ と $C$ は adic Noether なので，$\mathfrak a$ と $\mathfrak c$ は定義イデアルである
（詳細は省略する）。(4) から，$A\widehat{\otimes}_BC$ は adic であり，
$\mathfrak a(A\widehat{\otimes}_BC)+
\mathfrak c(A\widehat{\otimes}_BC)$ が有限生成の定義イデアルであることを
既に知っている。さらに
$$
A \widehat{\otimes}_B C /
\left(\mathfrak a(A \widehat{\otimes}_B C) +
\mathfrak c(A \widehat{\otimes}_B C)\right)
=
A/\mathfrak a \otimes_{B/\mathfrak b} C/\mathfrak c
$$
は Noether 環 $C/\mathfrak c$ 上の有限型代数として Noether であるから，
『代数』補題 \ref{algebra-lemma-completion-Noetherian} により結論を得る。
\end{proof}









\section{taut な環写像}
\label{formal-spaces-section-taut-ring-maps}

\noindent
線形位相環の間の連続写像がもつ次の性質に名前を付けておくと便利である。

\begin{definition}
\label{formal-spaces-definition-taut}
$\varphi:A\to B$ を線形位相環の連続写像とする。
すべての開イデアル $I\subset A$ に対して，イデアル $\varphi(I)B$ の閉包が
開であり，これらの閉包が開イデアルの基本系をなすとき，
$\varphi$ は{\it taut}であるという
\footnote{これは標準的な用語ではない。この定義は加群にも一般化される。
すなわち，線形位相を備えた $A$-加群 $M$ が $A$-taut であるとは，
すべての開イデアル $I\subset A$ に対して $M$ における $IM$ の閉包が開で，
これらの閉包が $M$ における $0$ の近傍の基本系をなすことをいう。}。
\end{definition}

\noindent
$\varphi:A\to B$ が線形位相環の連続写像で，$I_\lambda$ が $A$ の
開イデアルの基本系なら，$\varphi$ が taut であることと，
% P10-E0295
$I_\lambda B$ の閉包が開で，$A$ の開イデアルの基本系をなすこととは同値である。

\begin{lemma}
\label{formal-spaces-lemma-taut-weakly-admissible}
$\varphi:A\to B$ を弱許容位相環の連続写像とする。次は同値である。
\begin{enumerate}
\item $\varphi$ は taut である。
\item すべての弱定義イデアル $I\subset A$ に対して，
$\varphi(I)B$ の閉包は $B$ の弱定義イデアルであり，これらは $B$ の
弱定義イデアルの基本系をなす。
\end{enumerate}
\end{lemma}

\begin{proof}
定義 \ref{formal-spaces-definition-taut} に続く注意から (2) は (1) を含意する。逆に，
$\varphi$ が taut であると仮定する。$I\subset A$ が弱定義イデアルなら，
$\varphi(I)B$ の閉包は taut 性の定義により開であり，補題
\ref{formal-spaces-lemma-topologically-nilpotent} により位相的冪零元のみからなる。
したがって $\varphi(I)B$ の閉包は弱定義イデアルである。さらに，
taut 性の定義により，これらのイデアルは開イデアルの基本系をなすので，
(2) が成り立つ。
\end{proof}

\begin{lemma}
\label{formal-spaces-lemma-completion-taut}
$A$ を線形位相環とする。$A$ からその完備化への写像
$A\to A^\wedge$ は taut である。
\end{lemma}

\begin{proof}
$I_\lambda$ を $A$ の開イデアルの基本系とする。
$A^\wedge=\lim A/I_\lambda$ には極限位相を入れたことを思い出そう。
これは，核 $J_\lambda=\Ker(A^\wedge\to A/I_\lambda)$ が
$A^\wedge$ の開イデアルの基本系をなすことを意味する。
$J_\lambda$ は $I_\lambda A^\wedge$ の閉包なので
（補題 \ref{formal-spaces-lemma-closure-image-ideal} と比較せよ），結論を得る。
\end{proof}

\begin{lemma}
\label{formal-spaces-lemma-composition-taut}
$A\to B$ と $B\to C$ を線形位相環の連続準同型とする。
$A\to B$ と $B\to C$ が taut なら，$A\to C$ は taut である。
\end{lemma}

\begin{proof}
省略する。ヒント：$I\subset A$ がイデアルで，$J$ が $IB$ の閉包なら，
$JC$ の閉包は $IC$ の閉包に等しい。
\end{proof}

\begin{lemma}
\label{formal-spaces-lemma-permanence-taut}
$A\to B$ と $B\to C$ を線形位相環の連続準同型とする。
$A\to C$ が taut なら，$B\to C$ は taut である。
\end{lemma}

\begin{proof}
$J\subset B$ を開イデアルとし，その逆像を $I\subset A$ とする。
$JC$ の閉包は $IC$ の閉包を含む。$A\to C$ は taut なので，この閉包は
開である。$I_\lambda$ を $A$ の開イデアルの基本系とする。
$K_\lambda$ を $I_\lambda C$ の閉包とする。$A\to C$ は taut なので，
これらは $C$ の開イデアルの基本系をなす。$K_\lambda$ の逆像を
$J_\lambda\subset B$ と書く。このとき $J_\lambda C$ の閉包は
$K_\lambda$ である。したがって，$J$ が $B$ の開イデアルを動くとき，
イデアル $JC$ の閉包は $C$ の開イデアルの基本系をなす。
\end{proof}

\begin{lemma}
\label{formal-spaces-lemma-base-change-taut}
$A\to B$ と $A\to C$ を線形位相環の連続準同型とする。
$A\to B$ が taut なら，$C\to B\widehat{\otimes}_AC$ は taut である。
\end{lemma}

\begin{proof}
$K\subset C$ を開イデアルとする。
% P10-E0296
$C$ における像が $J$ に含まれる任意の開イデアル $I\subset A$ を選ぶ。
仮定により，$IB$ の閉包 $J$ は開である。$A\to B$ は taut なので，
$B\widehat{\otimes}_AC$ は，$K$ と $I$ のすべての選択にわたる環
$B/J\otimes_{A/I}C/K$ の極限である。すなわち，イデアル
$J(B\widehat{\otimes}_AC)+K(B\widehat{\otimes}_AC)$ が
開イデアルの基本系をなす。さて，$B\to B\widehat{\otimes}_AC$ は
連続なので，$J$ は $K(B\widehat{\otimes}_AC)$ の閉包へ写る
（$I$ は $K$ へ写る）。したがって，この閉包は
$J(B\widehat{\otimes}_AC)+K(B\widehat{\otimes}_AC)$ に等しく，
証明が完了する。
\end{proof}

\begin{lemma}
\label{formal-spaces-lemma-taut-ascent-countable}
$\varphi:A\to B$ を線形位相環の連続準同型とする。
$\varphi$ が taut で，$A$ が可算な開イデアルの基本系をもつなら，
$B$ も可算な開イデアルの基本系をもつ。
\end{lemma}

\begin{proof}
定義から直ちに従う。
\end{proof}

\begin{lemma}
\label{formal-spaces-lemma-taut-ascent-weakly-admissible}
$\varphi:A\to B$ を線形位相環の連続準同型とする。
$\varphi$ が taut で，$A$ が弱前許容なら，$B$ は弱前許容である。
\end{lemma}

\begin{proof}
$I\subset A$ を弱定義イデアルとする。補題
\ref{formal-spaces-lemma-topologically-nilpotent} により，$IB$ の閉包 $J$ は開で，
位相的冪零元のみからなる。したがって $J$ は $B$ の弱定義イデアルである。
\end{proof}

\begin{lemma}
\label{formal-spaces-lemma-taut-ascent-admissible}
$\varphi:A\to B$ を線形位相環の連続準同型とする。
$\varphi$ が taut で，$A$ が前許容なら，$B$ は前許容である。
\end{lemma}

\begin{proof}
$I\subset A$ を定義イデアルとする。$I_\lambda\subset A$ を
開イデアルの基本系とする。このとき，$IB$ の閉包 $J$ は開であり，
$I_\lambda B$ の閉包 $J_\lambda$ は開で，$B$ の開イデアルの基本系をなす。
すべての $\lambda$ に対して，$I^n\subset I_\lambda$ となる $n$ が存在する。
$J^n$ は $I^nB$ の閉包に含まれることに注意する。したがって
$J^n\subset J_\lambda$ であり，$J$ は定義イデアルであると結論する。
\end{proof}

\begin{lemma}
\label{formal-spaces-lemma-dense-image-surjective}
$\varphi:A\to B$ を線形位相環の連続準同型とし，次を仮定する。
\begin{enumerate}
\item $\varphi$ は taut で，像が稠密である。
\item $A$ は完備で，可算な開イデアルの基本系をもつ。
\item $B$ は分離的である。
\end{enumerate}
このとき $\varphi$ は全射かつ開であり，$B$ は完備で，ある閉イデアル
$K\subset A$ に対して $B=A/K$ である。
\end{lemma}

\begin{proof}
開写像補題（『代数詳論』補題
\ref{more-algebra-lemma-open-mapping}）と $\varphi$ の taut 性を合わせると，
写像 $\varphi$ は開である。$\varphi$ の像は稠密なので，$\varphi$ は全射である。
$\varphi$ の核 $K$ は，$\varphi$ が連続なので閉である。したがって
$B=A/K$ は完備である。例えば補題 \ref{formal-spaces-lemma-quotient-by-closed} を参照せよ。
\end{proof}







\section{adic 環写像}
\label{formal-spaces-section-adic-ring-maps}

\noindent
次の定義を置く。

\begin{definition}
\label{formal-spaces-definition-adic-homomorphism}
$A$ と $B$ を pre-adic 位相環とする。環準同型
$\varphi:A\to B$ が{\it adic}であるとは
\footnote{これは標準的でない用語かもしれない。}，定義イデアル
$I\subset A$ であって，$B$ の位相が $I$-進位相となるものが存在することをいう。
\end{definition}

\noindent
$\varphi:A\to B$ が pre-adic 環の adic 準同型なら，$\varphi$ は連続であり，
$A$ の任意の定義イデアル $I$ に対して $B$ の位相は $I$-進位相である。

\begin{lemma}
\label{formal-spaces-lemma-composition-adic}
$A\to B$ と $B\to C$ を pre-adic 環の連続準同型とする。
$A\to B$ と $B\to C$ が adic なら，$A\to C$ は adic である。
\end{lemma}

\begin{proof}
省略する。
\end{proof}

\begin{lemma}
\label{formal-spaces-lemma-permanence-adic}
$A\to B$ と $B\to C$ を pre-adic 環の連続準同型とする。
$A\to C$ が adic なら，$B\to C$ は adic である。
\end{lemma}

\begin{proof}
$A$ の定義イデアル $I$ を選ぶ。$A\to C$ は adic なので，
$IC$ は $C$ の定義イデアルである。$B\to C$ は連続なので，
$IC$ へ写る定義イデアル $J\subset B$ を見つけられる。
$A\to B$ は連続なので，$I$ における $J$ の逆像
$I'\subset I$ も $A$ の定義イデアルである。したがって
$I'C\subset JC\subset IC$ は二つの定義イデアルに挟まれているので，
それ自身が定義イデアルである。
\end{proof}

\begin{lemma}
\label{formal-spaces-lemma-adic-taut}
$\varphi:A\to B$ を pre-adic 位相環の間の連続準同型とする。
$\varphi$ が adic なら，$\varphi$ は taut である。
\end{lemma}

\begin{proof}
定義から直ちに従う。
\end{proof}

\noindent
次の補題は二つのことを述べる。
\begin{enumerate}
\item adic であるという性質は，完備な線形位相環の taut 写像に沿って上昇する。
\item $A$ が有限生成の定義イデアルをもつとき，adic 環の間の連続写像
$\varphi:A\to B$ に対して，「$\varphi$ は taut」と
「$\varphi$ は adic」は同値である。
\end{enumerate}
(2) により，「taut 性」は任意の線形位相環の間の連続環写像に対する
「adic 性」の一般化であるということができる。節 \ref{formal-spaces-section-adic} も参照せよ。

\begin{lemma}
\label{formal-spaces-lemma-taut-is-adic}
$\varphi:A\to B$ を線形位相環の連続写像とする。
$\varphi$ が taut で，$A$ が pre-adic かつ有限生成の定義イデアルをもち，
$B$ が完備なら，$B$ は adic で有限生成の定義イデアルをもち，環写像
$\varphi$ は adic である。
\end{lemma}

\begin{proof}
$A$ の有限生成の定義イデアル $I$ を選ぶ。$J_n$ を $B$ における
$\varphi(I^n)B$ の閉包とする。$B$ は完備なので $B=\lim B/J_n$ である。
$B'=\lim B/I^nB$ を $B$ の $I$-進完備化とする。『代数』補題
\ref{algebra-lemma-hathat-finitely-generated} により，$B'$ 上の
$I$-進位相は完備で，$B'/I^nB'=B/I^nB$ である。したがって環写像
$B'\to B$ は連続であり，すべての $n$ に対して
$B'\to B/I^nB\to B/J_n$ が全射なので，その像は稠密である。最後に，
$(I^nB')B=I^nB$ であり $A\to B$ は taut なので，写像
$B'\to B$ は taut である。補題 \ref{formal-spaces-lemma-dense-image-surjective} により，
$B'\to B$ は開かつ全射である。したがって $B$ の位相は $I$-進位相であり，
証明が完了する。
\end{proof}








\section{弱 adic 環}
\label{formal-spaces-section-weakly-adic}

\noindent
本節は読み飛ばしてもよい。以下は adic 環の自然な一般化である。

\begin{definition}
\label{formal-spaces-definition-weakly-adic}
\begin{reference}
\cite[Definition 8.3.8]{Gabber-Ramero}
\end{reference}
$A$ を線形位相環とする。
\begin{enumerate}
\item あるイデアル $I\subset A$ が存在して，すべての $n\geq0$ に対して
$I^n$ の閉包が開であり，これらの閉包が開イデアルの基本系をなすとき，
$A$ は{\it 弱 pre-adic}であるという
\footnote{\cite{Gabber-Ramero} では，著者らは $A$ を{\it $c$-adic}と呼ぶ。}。
\item $A$ が弱 pre-adic かつ完備であるとき，$A$ は{\it 弱 adic}であるという
\footnote{本章の規約では，これは分離的であることも含む。}。
\end{enumerate}
\end{definition}

\noindent
完備な線形位相環について，次の含意がある。
$$
\xymatrix{
\text{adic + Noether} \ar@{=>}[d] \\
\text{adic + 有限生成の定義イデアル} \ar@{=>}[d] \\
\text{adic} \ar@{=>}[d] \\
\text{弱 adic} \ar@{=>}[d] \\
\text{許容 + 第一可算} \ar@{=>}[d] \ar@{=>}[r] &
\text{許容} \ar@{=>}[d] \\
\text{弱許容 + 第一可算} \ar@{=>}[r] &
\text{弱許容}
}
$$
ここで「第一可算」とは，位相環が可算な開イデアルの基本系をもつことを意味する。
完備とは限らない線形位相環にも，同様の含意図式がある（すなわち，
pre-adic，弱 pre-adic，前許容，弱前許容という概念を用いる）。
pre-adic 環の場合とは異なり，次の弱 pre-adic 環の特徴づけが示すように，
弱 pre-adic 環の完備化は弱 adic である。

\begin{lemma}
\label{formal-spaces-lemma-weakly-adic}
$A$ を線形位相環とする。次は同値である。
\begin{enumerate}
\item $A$ は弱 pre-adic である。
\item $A'$ が pre-adic 位相環であるような taut 連続環写像
$A'\to A$ が存在する。
\item $A$ は前許容であり，すべての $n\geq1$ に対して $I^n$ の閉包が
開となる定義イデアル $I$ が存在する。
\item $A$ は前許容であり，すべての定義イデアル $I$ とすべての
$n\geq1$ に対して $I^n$ の閉包が開である。
\end{enumerate}
弱 pre-adic 環の完備化は弱 adic である。
$A$ が弱 adic なら，$A$ は許容で，可算な開イデアルの基本系をもつ。
\end{lemma}

\begin{proof}
(1) を仮定する。すべての $n$ に対して $I^n$ の閉包が開で，これらの閉包が
開イデアルの基本系をなすようなイデアル $I$ を選ぶ。
$A'=A$ に $I$-進位相を入れたものとする。このとき定義により
$A'\to A$ は taut であり，(2) が成り立つ。

\medskip\noindent
(2) を仮定する。$I'\subset A'$ を定義イデアルとする。
$I'A$ の閉包を $I$ と書く。$A'\to A$ の taut 性は，
$(I')^nA$ の閉包 $I_n$ が開で，開イデアルの基本系をなすことを意味する。
したがって $I=I_1$ は開であり，$I^n$ の閉包は $I_n$ に等しいので開で，
開イデアルの基本系をなす。したがって，$I$ は確かに定義イデアルであり，
すべての $n$ に対して $I^n$ の閉包は開である。よって (3) が成り立つ。

\medskip\noindent
$I\subset A$ が (3) におけるものなら，$I$ は定義
\ref{formal-spaces-definition-weakly-adic} におけるイデアルであり，(1) が成り立つ。
また，$I'\subset A$ が他の任意の定義イデアルなら，$I'$ は開である
（『代数詳論』定義 \ref{more-algebra-definition-topological-ring} を参照せよ）。
したがって，ある $n\geq1$ に対して，これは $I^n$ を含む。
よってすべての $m\geq1$ に対して $(I')^m$ は $I^{nm}$ を含み，
すべての $m$ に対して $(I')^m$ の閉包が開であると結論する。
このようにして (3) は (4) を含意する。(4) $\Rightarrow$ (3) は自明である。

\medskip\noindent
$A$ を弱 pre-adic とする。(2) における $A'\to A$ を選ぶ。補題
\ref{formal-spaces-lemma-completion-taut} および \ref{formal-spaces-lemma-composition-taut} により，
合成 $A'\to A^\wedge$ は taut である。したがって，(2) と (1) の同値性により
$A^\wedge$ は弱 pre-adic である。線形位相環 $A$ の完備化は完備なので
（『代数詳論』節 \ref{more-algebra-section-topological-ring}），
$A^\wedge$ は弱 adic である。

\medskip\noindent
% P10-E0297
$A$ を弱 adic とする。このとき $A$ は完備であり，
(1) $\Rightarrow$ (3) により前許容なので，$A$ は許容である。
もちろん定義により，$A$ は可算な開イデアルの基本系をもつ。
\end{proof}

\noindent
弱 adic 環が adic で有限生成の定義イデアルをもつことを保証する二つの判定法を
与える。

\begin{lemma}
\label{formal-spaces-lemma-weakly-adic-finite-generation-bis}
$A$ を完備な線形位相環とする。$I\subset A$ を有限生成イデアルとし，
すべての $n\geq0$ に対して $I^n$ の閉包が開で，これらの閉包が
開イデアルの基本系をなすと仮定する。このとき $A$ は adic で，
有限生成の定義イデアルをもつ。
\end{lemma}

\begin{proof}
% P10-E0638
$A'$ を，環 $A$ に $I$-進位相を入れたものとする。仮定により
$A'\to A$ は taut である。補題 \ref{formal-spaces-lemma-taut-is-adic} により結論を得る
（念のため付言すると，この補題は $I$ が定義イデアルであることも示す）。
\end{proof}

\begin{lemma}
\label{formal-spaces-lemma-weakly-adic-finite-generation}
$A$ を弱 adic 位相環とする。$I$ を定義イデアルとし，$I_2$ を $I^2$ の閉包と
する。$I/I_2$ が有限生成加群であると仮定する。このとき $A$ は adic で，
有限生成の定義イデアルをもつ。
\end{lemma}

\begin{proof}
補題 \ref{formal-spaces-lemma-weakly-adic} の特徴づけを以後断りなく用いる。
$I/I_2$ の生成元へ写る $f_1,\ldots,f_r\in I$ を選ぶ。
$I'=(f_1,\ldots,f_r)$ とおく。このとき $I'+I_2=I$ である。
$I_2$ は $I^2=(I'+I_2)^2\subset I'+I_3$ の閉包である。ここで
$I_3$ は $I^3$ の閉包である。したがって $I'+I_3=I$ である。
同様に続けると，$I_n$ を $I^n$ の閉包として，すべての $n\geq2$ に対して
$I'+I_n=I$ であることが分かる。言い換えると，$A$ における $I'$ の閉包は
$I$ である。したがって $(I')^n$ の閉包は $I_n$ である。よって
$(I')^n$ の閉包は $A$ の開イデアルの基本系をなす。補題
\ref{formal-spaces-lemma-weakly-adic-finite-generation-bis} により結論を得る。
\end{proof}

\noindent
% P10-E0639
「弱 pre-adic」という性質の重要な特徴は，線形位相環の taut 環準同型に沿って
上昇することである。

\begin{lemma}
\label{formal-spaces-lemma-taut-ascent-weakly-adic}
$\varphi:A\to B$ を線形位相環の連続準同型とする。
$\varphi$ が taut で，$A$ が弱 pre-adic なら，$B$ は弱 pre-adic である。
\end{lemma}

\begin{proof}
$I\subset A$ を，$I^n$ の閉包 $I_n$ が開で，これらの閉包が開イデアルの
基本系をなすようなイデアルとする。$I^nB$ の閉包は $I_nB$ の閉包に等しい。
$\varphi$ は taut なので，これらの閉包は開で，$B$ の開イデアルの基本系を
なす。したがって $B$ は弱 pre-adic である。
\end{proof}

\begin{lemma}
\label{formal-spaces-lemma-completed-tensor-product-weakly-adic}
$B\to A$ と $B\to C$ を線形位相環の連続準同型とする。
$A$ と $C$ が弱 pre-adic なら，$A\widehat{\otimes}_BC$ は弱 adic である。
\end{lemma}

\begin{proof}
補題 \ref{formal-spaces-lemma-weakly-adic} の特徴づけを以後断りなく用いる。
補題 \ref{formal-spaces-lemma-completed-tensor-product} により，
$A\widehat{\otimes}_BC$ が許容であることを知っている。さらに，その補題の証明は，
% P10-E0298
$I\subset A$ と $J\subset C$ が定義イデアルであるとき，
$I(A\widehat{\otimes}_BC)+J(A\widehat{\otimes}_BC)$ の閉包
$K\subset A\widehat{\otimes}_BC$ が定義イデアルであることを示す。
そこで，すべての $n\geq1$ に対して $K^n$ の閉包が開であることを示せばよい。
イデアル $K^n$ は
$I^n(A\widehat{\otimes}_BC)+J^n(A\widehat{\otimes}_BC)$ を含み，
$A$ における $I^n$ の閉包は開で，$C$ における $J^n$ の閉包も開なので，
$K^n$ の閉包が $A\widehat{\otimes}_BC$ において開であることが分かる。
\end{proof}






\section{性質の下降}
\label{formal-spaces-section-taut-descent}

\noindent
本節では，次の状況を考える。
\begin{enumerate}
\item $\varphi:A\to B$ は線形位相を備えた位相環の連続写像である。
\item $\varphi$ は taut である。
\item すべての開イデアル $I\subset A$ に対して，$J\subset B$ が
$IB$ の閉包を表すなら，写像 $A/I\to B/J$ は忠実平坦である。
\end{enumerate}
この状況では，$B$ の性質が $A$ に受け継がれることを示す。

\begin{lemma}
\label{formal-spaces-lemma-taut-descent-countable}
上の状況で，$B$ が可算な開イデアルの基本系をもつなら，
$A$ も可算な開イデアルの基本系をもつ。
\end{lemma}

\begin{proof}
開イデアルの基本系 $B\supset J_1\supset J_2\supset\ldots$ を選ぶ。
$\varphi$ の taut 性により，すべての $n$ に対して
$J_n\supset I_nB$ となる開イデアル $I_n$ を見つけられる。
$I_n$ が $A$ の開イデアルの基本系であると主張する。実際，
$I\subset A$ が開であると仮定する。$\varphi$ は taut なので，
$IB$ の閉包は開であり，したがって十分大きいある $n$ に対して $J_n$ を含む。
よって $I_nB\subset IB$ である。$B$ における $IB$ の閉包を $J$ とする。
$A/I\to B/J$ は忠実平坦なので単射である。
$I_nB\subset IB\subset J$ により $I_n\to A/I\to B/J$ は零だから，
$I_n\to A/I$ は零である。したがって $I_n\subset I$ となり，主張を得る。
\end{proof}

\begin{lemma}
\label{formal-spaces-lemma-taut-descent-weakly-admissible}
上の状況で，$B$ が弱前許容なら，$A$ も弱前許容である。
\end{lemma}

\begin{proof}
$J\subset B$ を弱定義イデアルとする。$IB\subset J$ となる開イデアル
$I\subset A$ を取る。$I$ が弱定義イデアルであることを示すには，
任意の $f\in I$ が位相的冪零であることを示せばよい。
$I'\subset A$ を開イデアルとする。$I'B$ の閉包を $J'\subset B$ と書く。
$A/I'\to B/J'$ は忠実平坦なので単射である。したがって
$f^n\in I'$ を示すには $\varphi(f)^n\in J'$ を示せば十分である。
$\varphi(f)\in J$ であり，イデアル $J$ は $B$ の弱定義イデアルで，
$J'$ は $B$ において開なので，これは $n\gg0$ に対して成り立つ。
\end{proof}

\begin{lemma}
\label{formal-spaces-lemma-taut-descent-admissible}
上の状況で，$B$ が前許容なら，$A$ も前許容である。
\end{lemma}

\begin{proof}
% P10-E0299
$J\subset B$ を弱定義イデアルとする。$IB\subset J$ となる開イデアル
$I\subset A$ を取る。$I'\subset A$ を開イデアルとする。
$I$ が定義イデアルであることを示すには，$n\gg0$ に対して
$I^n\subset I'$ となることを示せばよい。$I'B$ の閉包を $J'\subset B$ と
書く。$A/I'\to B/J'$ は忠実平坦なので単射である。したがって
$I^n\subset I'$ を示すには $\varphi(I)^n\subset J'$ を示せば十分である。
$\varphi(I)\subset J$ であり，イデアル $J$ は $B$ の定義イデアルで，
$J'$ は $B$ において開なので，これは $n\gg0$ に対して成り立つ。
\end{proof}

\begin{lemma}
\label{formal-spaces-lemma-taut-descent-weakly-adic}
上の状況で，$B$ が弱 pre-adic なら，$A$ も弱 pre-adic である。
\end{lemma}

\begin{proof}
補題 \ref{formal-spaces-lemma-weakly-adic} による弱 pre-adic 環の特徴づけを
以後断りなく用いる。補題 \ref{formal-spaces-lemma-taut-descent-admissible} により，
位相環 $A$ は前許容である。$I\subset A$ を定義イデアルとし，
$n\geq1$ を固定する。補題を証明するには，$I^n$ の閉包が開であることを
示さなければならない。$I_\lambda\subset A$ を開イデアルの基本系とする。
$IB$ の閉包を $J\subset B$ と書き，$I_\lambda B$ の閉包を
$J_\lambda\subset B$ と書く。$B$ は弱 pre-adic なので，$J^n$ の閉包は
開である。したがって，ある $\lambda$ に対して
$$
J_\lambda \subset \bigcap\nolimits_\mu (J^n + J_\mu)
$$
となる。右辺は補題 \ref{formal-spaces-lemma-closed} により $J^n$ の閉包だからである。
これは，$B/J_\mu$ における $J_\lambda$ の像が，$B/J_\mu$ における
$J^n$ の像に含まれることを意味する。$B/J_\mu$ における $J^n$ の像は，
$B/J_\mu$ における $I^nB$ の像に等しいことに
注意する（$J$ の各元は $J_\mu$ を法として $IB$ の元と合同である）。
$A/I_\mu\to B/J_\mu$ は忠実平坦で，$I_\lambda B\subset J_\lambda$ なので，
$A/I_\mu$ における $I_\lambda$ の像は $I^n$ の像に含まれると結論する。
したがって $I_\lambda$ は $I^n$ の閉包に含まれ，証明が完了する。
\end{proof}

\begin{lemma}
\label{formal-spaces-lemma-taut-descent-adic-star}
上の状況で，$B$ が adic で有限生成の定義イデアルをもち，$A$ が完備なら，
$A$ は adic で有限生成の定義イデアルをもつ。
\end{lemma}

\begin{proof}
補題 \ref{formal-spaces-lemma-taut-descent-weakly-adic} により，$A$ は既に弱 adic であり，
したがって特に許容である（adic 環が弱 adic であることには補題
\ref{formal-spaces-lemma-weakly-adic} も用いる）。$I\subset A$ を定義イデアルとする。
$J\subset B$ を有限生成の定義イデアルとする。$IB$ の閉包は開なので，
$J^n$ が $IB$ の閉包に含まれるような $n>0$ を見つけられる。
そこで $J$ を $J^n$ に置き換えることにより，$J$ が $IB$ の閉包に含まれる
有限生成の定義イデアルであると仮定してよい。補題 \ref{formal-spaces-lemma-closed} により，
これは確かに
$$
J \subset IB + J^2
$$
を含意する。有限生成 $A$-加群 $M=(J+IB)/IB$ を考える。上の式は
$JM=M$ を示す。例えば補題 \ref{formal-spaces-lemma-weakly-admissible-henselian} により，
$J$ は $B$ の Jacobson 根基に含まれることが分かる。したがって Nakayama の補題，
より正確には『代数』補題 \ref{algebra-lemma-NAK} の (2) により，
$M=0$ と結論する。よって $J\subset IB$ である。

\medskip\noindent
$J$ は有限生成なので，$J\subset I'B$ となる有限生成イデアル
$I'\subset I$ を見つけられる。$A\to B$ は連続で，$J\subset B$ は開で，
$I$ は定義イデアルなので，$I^nB\subset J$ となる $n>0$ を見つけられる。
$J_{n+1}\subset B$ を $I^{n+1}B$ の閉包とする。このとき
$$
I^n \cdot (B/J_{n + 1}) \subset J \cdot (B/J_{n + 1}) \subset
I' \cdot (B/J_{n + 1})
$$
である。$A/I^{n+1}\to B/J_{n+1}$ は忠実平坦なので，これは
$I^n\cdot(A/I^{n+1})\subset I'\cdot(A/I^{n+1})$ を含意し，さらに
$$
I^n \subset I' + I^{n + 1}
$$
を意味する。これはすべての $k\geq1$ に対して
$I^n\subset I'+I^{n+k}$ を含意し，さらにすべての $k,m\geq1$ に対して
$I^{nm}\subset(I')^m+I^{nm+k}$ を含意する。よって $(I')^m$ の閉包は
$I^{nm}$ を含む。$A$ は弱 adic なので $I^{nm}$ の閉包は開である。
% P10-E0300
したがって，すべての $m$ に対して $(I')^m$ の閉包は開である。
これらの閉包は $A$ の開イデアルの基本系をなす
（$I^n$ の閉包について同じことが成り立つからである）。
補題 \ref{formal-spaces-lemma-weakly-adic-finite-generation-bis} により結論を得る。
\end{proof}









\section{アフィン形式代数空間}
\label{formal-spaces-section-affine-formal-algebraic-spaces}

\noindent
この節ではアフィン形式代数空間を導入する。実際，これは \cite{BVGD} で
アフィン形式スキームと呼ばれているものと同じである。しかし，\cite{EGA} で
定義されるアフィン形式スキームの概念との混同を避けるため，ここではこれを
アフィン形式代数空間と呼ぶ。

\medskip\noindent
スキームの肥厚とは，基礎位相空間上の全射を誘導する閉埋め込みであったことを
想起しよう。『射についてさらに』定義
\ref{more-morphisms-definition-thickening} を参照せよ。

\begin{definition}
\label{formal-spaces-definition-affine-formal-algebraic-space}
$S$ をスキームとする。$(\Sch/S)_{fppf}$ 上の層 $X$ が
{\it アフィン形式代数空間}であるとは，次が存在することをいう：
\begin{enumerate}
\item 有向集合 $\Lambda$，
\item $(\Sch/S)_{fppf}$ における $\Lambda$ 上の系
$(X_\lambda, f_{\lambda \mu})$ であって，
\begin{enumerate}
\item 各 $X_\lambda$ はアフィンであり，
\item 各 $f_{\lambda \mu} : X_\lambda \to X_\mu$ は肥厚であるもの。
\end{enumerate}
\end{enumerate}
さらに，fppf 層として
$$
X \cong \colim_{\lambda \in \Lambda} X_\lambda
$$
であり，$X$ は集合論的条件を満たすものとする
（評注 \ref{formal-spaces-remark-set-theoretic} を参照）。$S$ 上の
{\it アフィン形式代数空間の射}とは，層の写像をいう。
\end{definition}

\noindent
系 $(X_\lambda, f_{\lambda \mu})$ はデータの一部ではないことに注意せよ。
$U$ を $S$ 上の準コンパクトなスキームとする。遷移写像はモノ射なので，
『サイト』補題 \ref{sites-lemma-directed-colimits-sections} により
$$
X(U) = \colim X_\lambda(U)
$$
である。したがって，定義の余極限に内在する fppf 層化は Zariski 層化であり，
準コンパクトなスキームに対しては何も変えない。

\begin{lemma}
\label{formal-spaces-lemma-diagonal-affine-formal-algebraic-space}
$S$ をスキームとする。$X$ が $S$ 上のアフィン形式代数空間ならば，
対角射 $\Delta : X \to X \times_S X$ は表現可能な閉埋め込みである。
\end{lemma}

\begin{proof}
$U,V$ を $S$ 上のスキームとし，$U \to X$ および $V \to X$ が与えられたとする。
$U \times_X V$ が表現可能であることを示そう。定義
\ref{formal-spaces-definition-affine-formal-algebraic-space} のように
$X = \colim X_\lambda$ と書く。上の議論から，$U$ と $V$ の上で Zariski 局所的に，
% P10-E0301
これらの射はある $X_\lambda$ を経由する。この場合
$U \times_X V = U \times_{X_\lambda} V$ はスキームである。したがって対角射は
表現可能である。『代数空間』補題
\ref{spaces-lemma-representable-diagonal} を参照せよ。
$W$ を $S$ 上のスキームとし，$(a,b):W\to X\times_SX$ が与えられたとき，写像
$X \times_{\Delta, X \times_S X, (a, b)} W \to W$ を考える。
先と同様に，$W$ 上局所的に射 $a$ と $b$ はある $\lambda$ に対するアフィン・スキーム
$X_\lambda$ に入るので，射
$X_\lambda
\times_{\Delta_\lambda, X_\lambda \times_S X_\lambda, (a, b)} W \to W$
を得る。これは
$\Delta_\lambda:X_\lambda\to X_\lambda\times_SX_\lambda$ の基底変換である。
$X_\lambda$ はアフィンなので $X_\lambda\to S$ は分離であり，この対角射は
閉埋め込みである。ゆえに $X\to X\times_SX$ は閉埋め込みである。
\end{proof}

\noindent
スキームの射 $X\to X'$ が肥厚であるとは，それが閉埋め込みで，点の基礎集合上の
全射を誘導することをいう（『射についてさらに』定義
\ref{more-morphisms-definition-thickening} を参照）。したがって，肥厚であるという
性質は任意の基底変換で保たれ，かつ標的上 fpqc 局所的である。『代数空間』節
\ref{spaces-section-lists} を参照せよ。よって『代数空間』定義
\ref{spaces-definition-relative-representable-property} を「肥厚」に適用でき，
$(\Sch/S)_{fppf}$ 上の前層の表現可能な変換 $F\to G$ が肥厚であるとは何かが
定まる。$F$ と $G$ が代数空間である場合，これは『代数空間の射についてさらに』定義
\ref{spaces-more-morphisms-definition-thickening} における肥厚の定義と矛盾しない。

\begin{lemma}
\label{formal-spaces-lemma-covering-by-thickenings}
定義 \ref{formal-spaces-definition-affine-formal-algebraic-space} のように
$X_\lambda, \lambda\in\Lambda$ と $X=\colim X_\lambda$ が与えられているとする。
このとき $X_\lambda\to X$ は表現可能な肥厚である。
\end{lemma}

\begin{proof}
この主張は『代数空間』節 \ref{spaces-section-representable} および
\ref{spaces-section-representable-properties} の議論により意味をもつ。
補題 \ref{formal-spaces-lemma-diagonal-affine-formal-algebraic-space} により，射
$X_\lambda\to X$ は表現可能である。$U$ をスキームとし $U\to X$ が与えられたとする。
定義 \ref{formal-spaces-definition-affine-formal-algebraic-space} の後の議論から，$U$ 上 Zariski
局所的に，この射は $\lambda\leq\mu$ であるある $X_\mu$ を経由する。
この場合 $U\times_XX_\lambda=U\times_{X_\mu}X_\lambda$ なので，
$U\times_XX_\lambda\to U$ は肥厚 $X_\lambda\to X_\mu$ の基底変換である。
\end{proof}

\begin{lemma}
\label{formal-spaces-lemma-factor-through-thickening}
定義 \ref{formal-spaces-definition-affine-formal-algebraic-space} のように
$X_\lambda, \lambda\in\Lambda$ と $X=\colim X_\lambda$ が与えられているとする。
$Y$ が $S$ 上の準コンパクトな代数空間ならば，任意の射 $Y\to X$ はある
$X_\lambda$ を経由する。
\end{lemma}

\begin{proof}
アフィン・スキーム $V$ と全射エタール射 $V\to Y$ を選ぶ。合成
$V\to Y\to X$ は，定義 \ref{formal-spaces-definition-affine-formal-algebraic-space} の後の議論に
より，ある $\lambda$ に対する $X_\lambda$ を経由する。$V\to Y$ は層の全射なので，
結論を得る。
\end{proof}

\begin{lemma}
\label{formal-spaces-lemma-characterize-affine-formal-algebraic-space}
$S$ をスキームとし，$X$ を $(\Sch/S)_{fppf}$ 上の層とする。このとき，$X$ が
アフィン形式代数空間であるための必要十分条件は，次が成り立つことである：
\begin{enumerate}
\item $U$ が $S$ 上のアフィン・スキームである任意の射 $U\to X$ は，
$S$ 上のアフィン・スキーム $T$ をもち，表現可能な肥厚である射 $T\to X$ を
経由する。
\item 評注 \ref{formal-spaces-remark-set-theoretic} の集合論的条件が成り立つ。
\end{enumerate}
\end{lemma}

\begin{proof}
補題 \ref{formal-spaces-lemma-covering-by-thickenings} と
\ref{formal-spaces-lemma-factor-through-thickening} から，アフィン形式代数空間が (1)，(2) を
満たすことが従う。逆を示すには，$X$ が空でないと仮定してよい。
$T$ が $S$ 上のアフィン・スキームであり，表現可能な肥厚であるような
射 $T\to X$ の圏を $\Lambda$ とする。この圏は有向である。$X$ は空でないので，
$\Lambda$ は少なくとも一つの対象を含む。$T\to X$ と $T'\to X$ が $\Lambda$ に
属するならば，$T\amalg T'\to X$ を $\Lambda$ のある $T''\to X$ を通して
分解できる。$\Lambda$ の任意の二対象の間には，唯一の射が存在するか，射が
存在しないかのいずれかである。したがって $\Lambda$ は有向集合であり，仮定から
$X = \colim_{T \to X\text{ が }\Lambda\text{ に属する}} T$ である。
証明を終えるには，$\Lambda$ の任意の射 $T\to T'$ が肥厚であることを
示せばよい。$T'\to X$ は層のモノ射なので，
$T=T\times_{T'}T'=T\times_XT'$ である。ゆえに射 $T\to T'$ は射影
$T\times_XT'\to T'$ に等しく，これは $T\to X$ が肥厚であることから肥厚である。
\end{proof}

\noindent
一般のアフィン形式代数空間 $X$ に対し，例えば $X$ が点を分離するのに十分な関数をもつとは
限らない。『例』節 \ref{examples-section-affine-formal-algebraic-space} を参照せよ。
それをもつものを特徴づけるため，次の補題を与える。

\begin{lemma}
\label{formal-spaces-lemma-mcquillan-affine-formal-algebraic-space}
$S$ をスキームとする。$X$ を $(\Sch/S)_{fppf}$ 上の fppf 層で，評注
\ref{formal-spaces-remark-set-theoretic} の集合論的条件を満たすものとする。次は同値である：
\begin{enumerate}
\item $S$ 上の弱許容位相環 $A$（評注 \ref{formal-spaces-remark-mcquillan} を参照）が存在して
$X = \colim_{I \subset A\text{ は弱定義イデアル}} \Spec(A/I)$ となる。
\item $X$ はアフィン形式代数空間であり，$S$-代数 $A$ と写像
$X\to\Spec(A)$ が存在して，アフィン・スキーム $T$ による閉埋め込み
$T\to X$ に対し，合成 $T\to\Spec(A)$ が閉埋め込みとなる。
\item $X$ はアフィン形式代数空間であり，$S$-代数 $A$ と写像
$X\to\Spec(A)$ が存在して，スキーム $T$ による閉埋め込み $T\to X$ に対し，
合成 $T\to\Spec(A)$ が閉埋め込みとなる。
\item $X$ はアフィン形式代数空間であり，定義
\ref{formal-spaces-definition-affine-formal-algebraic-space} のような表示
$X=\colim X_\lambda$ のある選択に対して，射影
$\lim\Gamma(X_\lambda,\mathcal{O}_{X_\lambda})
\to\Gamma(X_\lambda,\mathcal{O}_{X_\lambda})$ が全射である。
\item $X$ はアフィン形式代数空間であり，定義
\ref{formal-spaces-definition-affine-formal-algebraic-space} のような表示
$X=\colim X_\lambda$ の任意の選択に対して，射影
$\lim\Gamma(X_\lambda,\mathcal{O}_{X_\lambda})
\to\Gamma(X_\lambda,\mathcal{O}_{X_\lambda})$ が全射である。
\end{enumerate}
さらに，この弱許容位相環は極限位相を備えた
$A=\lim\Gamma(X_\lambda,\mathcal{O}_{X_\lambda})$ であり，弱定義イデアルは，
表現可能な肥厚である射 $T\to X$ をちょうど分類する。
\end{lemma}

\begin{proof}
(5) が (4) を含意することは明らかである。

\medskip\noindent
定義 \ref{formal-spaces-definition-affine-formal-algebraic-space} のような
$X=\colim X_\lambda$ に対して (4) を仮定する。
\par\noindent
$A=\lim\Gamma(X_\lambda,\mathcal{O}_{X_\lambda})$ とおく。
$T$ をスキームとし，$T\to X$ を閉埋め込みとする（補題
\ref{formal-spaces-lemma-diagonal-affine-formal-algebraic-space} により $T\to X$ は表現可能で
あることに注意せよ）。$X_\lambda\to X$ は肥厚なので，
$X_\lambda\times_XT\to T$ も肥厚である。一方，
$X_\lambda\times_XT\to X_\lambda$ は閉埋め込みなので，
$X_\lambda\times_XT$ はアフィンである。したがって『極限』命題
\ref{limits-proposition-affine} により $T$ はアフィンである。
次に補題 \ref{formal-spaces-lemma-factor-through-thickening} により，$T\to X$ はある
$\lambda$ に対する $X_\lambda$ を経由する。よって
$A\to\Gamma(X_\lambda,\mathcal{O})\to\Gamma(T,\mathcal{O})$ は全射である。
このようにして (3) が成り立つ。

\medskip\noindent
(3) が (2) を含意することは明らかである。

\medskip\noindent
$A$ と $X\to\Spec(A)$ に対して (2) を仮定する。定義
\ref{formal-spaces-definition-affine-formal-algebraic-space} のように
$X=\colim X_\lambda$ と書く。仮定 (2) により
$A_\lambda=\Gamma(X_\lambda,\mathcal{O})$ は $A$ の商である。したがって
$A^\wedge=\lim A_\lambda$ は完備位相環である。『代数についてさらに』節
\ref{more-algebra-section-topological-ring} の議論を参照せよ。
$A\to A_\lambda$ が全射なので，$A^\wedge\to A_\lambda$ も全射である。
任意の $\lambda$ に対し，$A^\wedge\to A_\lambda$ の核
$I_\lambda\subset A^\wedge$ は弱定義イデアルであると主張する。
実際，それは極限位相の定義により開である。$f\in I_\lambda$ とすると，
任意の $\mu\in\Lambda$ に対し，$A_\mu$ における $f$ の像は $X_\mu$ の
すべての点の剰余体で零である。したがって，それは $A_\mu$ の冪零元であり，
ある冪について $f^n\in I_\mu$ となる。
% P10-E0302
よって $n\to0$ のとき $f^n\to0$ である。したがって $A^\wedge$ は弱許容である。
最後に，$I\subset A^\wedge$ を弱定義イデアルとする。すると
$I\subset A^\wedge$ は開なので，$I\supset I_\lambda$ となるある $\lambda$ が
存在する。これにより射
$\Spec(A^\wedge/I)\to\Spec(A_\lambda)\to X$ を得る。ゆえに，今度はすべての
弱定義イデアルにわたって余極限を取ると
$X=\colim\Spec(A^\wedge/I)$ となる。したがって (1) が成り立つ。

\medskip\noindent
(1) を仮定する。この場合，$X$ がアフィン形式代数空間であることは明らかである。
定義 \ref{formal-spaces-definition-affine-formal-algebraic-space} のような任意の表示
$X=\colim X_\lambda$ をとる。各 $\lambda$ に対し，補題
\ref{formal-spaces-lemma-factor-through-thickening} により，$X_\lambda\to X$ が
$\Spec(A/I)\to X$ を経由するような弱定義イデアル $I\subset A$ を見つけられる。
このとき $I\subset I_\lambda$ をもつ
$X_\lambda=\Spec(A/I_\lambda)$ である。逆に，任意の弱定義イデアル
$I\subset A$ に対し，射 $\Spec(A/I)\to X$ はある $\lambda$ に対する
$X_\lambda$ を経由する。
すなわち $I_\lambda\subset I$ である。したがって，各 $I_\lambda$ は
弱定義イデアルであり，これらは弱定義イデアル全体の集合の余終部分集合をなす。
ゆえに $A=\lim A/I=\lim A/I_\lambda$ であり，(5) が成り立つこと，さらに
$A=\lim\Gamma(X_\lambda,\mathcal{O}_{X_\lambda})$ であることが分かる。
\end{proof}

\noindent
この補題を用いて，次の定義を行える。

\begin{definition}
\label{formal-spaces-definition-types-affine-formal-algebraic-space}
$S$ をスキームとし，$X$ を $S$ 上のアフィン形式代数空間とする。
$X$ が補題 \ref{formal-spaces-lemma-mcquillan-affine-formal-algebraic-space} の同値な条件を
満たすとき，$X$ を {\it McQuillan} と呼ぶ。$A$ を $X$ に付随する弱許容位相環とする。
次のようにいう：
\begin{enumerate}
\item $X$ が {\it classical} であるとは，$X$ が McQuillan で $A$ が許容である
ことをいう（『代数についてさらに』定義
\ref{more-algebra-definition-topological-ring}）。
\item $X$ が {\it 弱 adic} であるとは，$X$ が McQuillan で $A$ が弱 adic で
あることをいう（定義 \ref{formal-spaces-definition-weakly-adic}）。
\item $X$ が {\it adic} であるとは，$X$ が McQuillan で $A$ が adic である
ことをいう（『代数についてさらに』定義
\ref{more-algebra-definition-topological-ring}）。
\item $X$ が {\it adic*} であるとは，$X$ が McQuillan で，$A$ が adic かつ
$A$ が有限生成の定義イデアルをもつことをいう。
\item $X$ が {\it Noether} であるとは，$X$ が McQuillan で，$A$ が Noether
かつ adic であることをいう。
\end{enumerate}
\end{definition}

\noindent
\cite{Fujiwara-Kato} では，adic 位相環 $A$ が有限生成の定義イデアルを含むという
性質に「有限イデアル型」という用語を使っている。アフィン形式代数空間 $X$ に
ついて，この定義で導入した概念の間には次の含意がある：
$$
\xymatrix{
X\text{ Noether} \ar@{=>}[r] &
X\text{ adic*} \ar@{=>}[r] &
X\text{ adic} \ar@{=>}[lld] \\
X\text{ 弱 adic} \ar@{=>}[r] &
X\text{ classical} \ar@{=>}[r] &
X\text{ McQuillan}
}
$$
節 \ref{formal-spaces-section-weakly-adic} の議論を参照し，精確な主張については補題
\ref{formal-spaces-lemma-implications-between-types} を参照せよ。

\begin{remark}
\label{formal-spaces-remark-compare-with-affine-formal-schemes}
classical アフィン形式代数空間は，EGA（\cite{EGA}）で考察されるアフィン形式
スキームに対応する。これを説明するため，基礎スキームを $\Spec(\mathbf{Z})$ と
仮定する。$\mathfrak X=\text{Spf}(A)$ をアフィン形式スキームとする。
補題 \ref{formal-spaces-lemma-fully-faithful} におけるその点関手を $h_\mathfrak X$ とする。
このとき $h_\mathfrak X=\colim h_{\Spec(A/I)}$ であり，ここで余極限は許容位相環
$A$ の定義イデアル全体にわたる。アフィン・スキーム上で評価すれば，これは
(\ref{formal-spaces-equation-morphisms-affine-formal-schemes}) から従う。また両辺は fppf 層なので，
アフィン・スキーム上で確認すれば十分である。補題
\ref{formal-spaces-lemma-formal-scheme-sheaf-fppf} を参照せよ。したがって
$h_\mathfrak X$ はアフィン形式代数空間である。実際，定義
\ref{formal-spaces-definition-types-affine-formal-algebraic-space} により，それは classical
アフィン形式代数空間である。よって補題 \ref{formal-spaces-lemma-fully-faithful} から，
アフィン形式スキームの圏は classical アフィン形式代数空間の圏と同値である。
\end{remark}

\noindent
以上でアフィン形式スキームとの関係を与えたので，次の定義は自然であろう。

\begin{definition}
\label{formal-spaces-definition-affine-formal-spectrum}
$S$ をスキームとする。$A$ を $S$ 上の弱許容位相環とする。定義
\ref{formal-spaces-definition-weakly-admissible}\footnote{classical の場合については
『代数についてさらに』定義
\ref{more-algebra-definition-topological-ring} を，差異の議論については評注
\ref{formal-spaces-remark-mcquillan} を参照せよ。}を参照せよ。
$A$ の {\it 形式スペクトル}とは，アフィン形式代数空間
$$
\text{Spf}(A) = \colim \Spec(A/I)
$$
をいう。ここで余極限は $A$ の弱定義イデアル全体の集合にわたり，圏
$\Sh((\Sch/S)_{fppf})$ で取られる。
\end{definition}

\noindent
このような形式スペクトルは構成により McQuillan であり，逆に，任意の
McQuillan アフィン形式代数空間は形式スペクトルと同型である。ただし，この理論には，
どの弱許容位相環の形式スペクトルでもないアフィン形式代数空間も存在する。
\cite{Yasuda} に従い，$S$-代数の圏の pro-対象として $S$-pro-環を導入することも
できる（『圏』評注 \ref{categories-remark-pro-category} を参照）。その場合，
$S$ 上の任意のアフィン形式代数空間は，そのような $S$-pro-環の形式スペクトルに
なる。ここではそうせず，対応するアフィン形式代数空間を直接扱う。

\medskip\noindent
形式スペクトルの構成は関手的である。これを説明するため，$S$ 上の弱許容位相環の
% P10-E0303
連続写像 $\varphi:B\to A$ をとる。このとき
$$
\text{Spf}(\varphi) : \text{Spf}(B) \to \text{Spf}(A)
$$
は，$I\subset A$ と $J\subset B$ が $\varphi(I)\subset J$ を満たす任意の
弱定義イデアルであるとき，図式
$$
\xymatrix{
\Spec(B/J) \ar[d] \ar[r] & \Spec(A/I) \ar[d] \\
\text{Spf}(B) \ar[r] & \text{Spf}(A)
}
$$
を可換にする唯一のアフィン形式代数空間の射である。$\varphi$ の連続性は，
任意の弱定義イデアル $J\subset B$ に対して，要求される性質をもつ弱定義イデアル
$I\subset A$ が存在することを含意する。したがって，上の可換性は
$\text{Spf}(\varphi)$ を一意に定め，かつ定義する。

\begin{lemma}
\label{formal-spaces-lemma-morphism-between-formal-spectra}
$S$ をスキームとし，$A$，$B$ を $S$ 上の弱許容位相環とする。
$S$ 上のアフィン形式代数空間の任意の射
$f:\text{Spf}(B)\to\text{Spf}(A)$ は，一意な連続 $S$-代数写像
$f^\sharp:A\to B$ に対する $\text{Spf}(f^\sharp)$ に等しい。
\end{lemma}

\begin{proof}
$f:\text{Spf}(B)\to\text{Spf}(A)$ を補題のような射とする。
$J\subset B$ を弱定義イデアルとする。補題
\ref{formal-spaces-lemma-factor-through-thickening} により，合成
$\Spec(B/J)\to\text{Spf}(B)\to\text{Spf}(A)$ が
$\Spec(A/I)$ を経由するような弱定義イデアル $I\subset A$ が存在する。
『スキーム』補題 \ref{schemes-lemma-morphism-into-affine} により，$S$-代数写像
$A/I\to B/J$ を得る。これらの写像は $J$ の変化に関して両立し，写像
$f^\sharp:A\to B$ を定める。任意の弱定義イデアル $J\subset B$ に対して，
$f^\sharp(I)\subset J$ を満たす弱定義イデアル $I\subset A$ が存在するので，
この写像は連続である。等式 $f=\text{Spf}(f^\sharp)$ は，$f^\sharp$ を構成する
環写像 $A/I\to B/J$ の選び方から従う。
\end{proof}

\begin{lemma}
\label{formal-spaces-lemma-presentation-representable}
$S$ をスキームとし，$f:X\to Y$ を $(\Sch/S)_{fppf}$ 上の前層の写像とする。
$X$ がアフィン形式代数空間で，$f$ が代数空間により表現可能かつ局所準有限ならば，
$f$ は（スキームにより）表現可能である。
\end{lemma}

\begin{proof}
$T$ を $S$ 上のスキームとし，写像 $T\to Y$ をとる。代数空間
$X\times_YT$ がスキームであることを示さなければならない。定義
\ref{formal-spaces-definition-affine-formal-algebraic-space} のように
$X=\colim X_\lambda$ と書く。$W\subset X\times_YT$ を準コンパクトな開部分空間とする。
射影 $X\times_YT\to X$ の $W$ への制限は，ある $\lambda$ に対する
$X_\lambda$ を経由する。このとき
$$
W \to X_\lambda \times_S T
$$
はモノ射（したがって分離）であり，局所準有限でもある（$X\to Y$ に関する仮定と
『代数空間の射』補題
\ref{spaces-morphisms-lemma-permanence-quasi-finite} により，
$W\to X\times_YT\to T$ は局所準有限だからである）。したがって
『代数空間の射』命題
\ref{spaces-morphisms-proposition-locally-quasi-finite-separated-over-scheme}
により $W$ はスキームである。ゆえに『代数空間の性質』補題
\ref{spaces-properties-lemma-subscheme} により $X\times_YT$ はスキームである。
\end{proof}






\section{可算添字付きアフィン形式代数空間}
\label{formal-spaces-section-countably-indexed}

\noindent
これは次の補題に現れるアフィン形式代数空間である。

\begin{lemma}
\label{formal-spaces-lemma-countable-affine-formal-algebraic-space}
$S$ をスキームとし，$X$ を $S$ 上のアフィン形式代数空間とする。次は同値である：
\begin{enumerate}
\item $S$ 上のアフィン・スキームの肥厚からなる系
$X_1\to X_2\to X_3\to\ldots$ で，$X=\colim X_n$ を満たすものが存在する。
\item 定義 \ref{formal-spaces-definition-affine-formal-algebraic-space} のような表示
$X=\colim X_\lambda$ で，$\Lambda$ が可算となる選択が存在する。
\end{enumerate}
\end{lemma}

\begin{proof}
可算有向集合は $(\mathbf{N},\geq)$ と同型な余終部分集合をもつことから従う。
『代数』補題 \ref{algebra-lemma-ML-limit-nonempty} の証明を参照せよ。
\end{proof}

\begin{definition}
\label{formal-spaces-definition-countable}
$S$ をスキームとし，$X$ を $S$ 上のアフィン形式代数空間とする。
補題 \ref{formal-spaces-lemma-countable-affine-formal-algebraic-space} の同値な条件を満たすとき，
$X$ は {\it 可算添字付き}であるという。
\end{definition}

\noindent
\cite{BVGD} の用語では，これは $X$ が $\aleph_0$-ind-スキームであるという。

\begin{lemma}
\label{formal-spaces-lemma-implications-between-types}
$X$ をスキーム $S$ 上のアフィン形式代数空間とする。
\begin{enumerate}
\item $X$ が Noether ならば，$X$ は adic* である。
\item $X$ が adic* ならば，$X$ は adic である。
\item $X$ が adic ならば，$X$ は弱 adic である。
\item $X$ が弱 adic ならば，$X$ は classical である。
\item $X$ が弱 adic ならば，$X$ は可算添字付きである。
\item $X$ が可算添字付きならば，$X$ は McQuillan である。
\end{enumerate}
\end{lemma}

\begin{proof}
% P10-E0304
(1)，(2)，(3)，(4) は，$X=\text{Spf}(A)$ と書き，$A$ が弱許容な
（したがって完備な）線形位相環であることと，節 \ref{formal-spaces-section-weakly-adic} で
論じたそのような環の各種の型の間の含意を用いれば従う。

\medskip\noindent
(5) の証明。定義により，$X=\colim\Spec(A/I)$ となる弱 adic 位相環 $A$ が
存在する。ここで余極限は $A$ の定義イデアル全体にわたる。$A$ は弱 adic なので，
% P10-E0305
特に開イデアルの可算基本系 $I_\lambda$ が存在する。定義
\ref{formal-spaces-definition-weakly-adic} を参照せよ。
% P10-E0306
このとき，$\text{Spf}(A)$ の定義により
$X=\colim\Spec(A/I_n)$ である。したがって $X$ は可算添字付きである。

\medskip\noindent
(6) の証明。$S$ 上のアフィン・スキームの肥厚からなるある系
$X_1\to X_2\to X_3\to\ldots$ に対して $X=\colim X_n$ と書く。このとき
$$
A = \lim \Gamma(X_n, \mathcal{O}_{X_n})
$$
は各 $\Gamma(X_n,\mathcal{O}_{X_n})$ へ全射となる。実際，射
$X_n\to X_{n+1}$ は閉埋め込みなので，遷移写像は全射である。
したがって $X$ は McQuillan である。
\end{proof}

\begin{lemma}
\label{formal-spaces-lemma-countably-indexed}
$S$ をスキームとし，$X$ を $(\Sch/S)_{fppf}$ 上の前層とする。次は同値である：
\begin{enumerate}
\item $X$ は可算添字付きアフィン形式代数空間である。
\item $X=\text{Spf}(A)$ であり，$A$ は $0$ の近傍の可算基本系をもつ
弱許容位相 $S$-代数である。
\item $X=\text{Spf}(A)$ であり，$A$ は弱定義イデアルの基本系
$A\supset I_1\supset I_2\supset I_3\supset\ldots$ をもつ
弱許容位相 $S$-代数である。
\item $X=\text{Spf}(A)$ であり，$A$ は，$I_n/I_{n+1}$ が局所冪零であるような
イデアルの可算列 $A\supset I_1\supset I_2\supset I_3\supset\ldots$ により
与えられる $0$ の開近傍の基本系をもつ完備位相 $S$-代数である。
\item $X=\text{Spf}(A)$ であり，$A=\lim B/J_n$ は極限位相をもち，ここで
$B\supset J_1\supset J_2\supset J_3\supset\ldots$ は $S$-代数 $B$ の
イデアルの列で，各 $J_n/J_{n+1}$ は局所冪零である。
\end{enumerate}
\end{lemma}

\begin{proof}
(1) を仮定する。補題 \ref{formal-spaces-lemma-implications-between-types} により，
$X=\text{Spf}(A)$ と書ける。ここで $A$ は弱許容位相 $S$-代数である。
補題 \ref{formal-spaces-lemma-countable-affine-formal-algebraic-space} の (1) のような任意の表示
$X=\colim X_n$ に対し，$X_n=\Spec(A_n)$，$A_n=A/I_n$ で，
$I_n\subset A$ はある弱定義イデアルであり，$A=\lim A_n$ となる。
これは補題 \ref{formal-spaces-lemma-mcquillan-affine-formal-algebraic-space} の最後の主張から従い，
同じ主張はさらに，$\{I_n\}$ が $0$ の開近傍の基本系であることを含意する。
したがって，弱定義イデアルの列
$$
A \supset I_1 \supset I_2 \supset I_3 \supset \ldots
$$
で $A=\lim A/I_n$ を満たすものを得る。このようにして，条件 (1) は条件
(2)--(5) のそれぞれを含意する。

\medskip\noindent
(5) を仮定する。まず，$A=\lim B/J_n$ 上の極限位相は線形位相かつ完備である。
『代数についてさらに』節 \ref{more-algebra-section-topological-ring} を参照せよ。
$f\in A$ が $B/J_1$ で零に写るならば，その $J_1/J_2$ における像が冪零なので，
ある冪は $B/J_2$ で零に写り，さらにそのある冪は $J_2/J_3$ で零に写る，等々である。
このようにして，開イデアル $\Ker(A\to B/J_1)$ は弱定義イデアルである。
したがって $A$ は弱許容である。これにより (5) が (2) を含意することが分かる。

\medskip\noindent
$B=A$ と取れば (4) が (5) の特別な場合であることは明らかである。
(3) が (2) の特別な場合であることも明らかである。

\medskip\noindent
$A$ が (2) のようなものと仮定する。$E_n$ を $A$ における $0$ の近傍の
可算基本系とする。$A$ は弱許容位相環なので，開イデアル
$I_n\subset E_n$ を見つけられる。弱定義イデアル $J\subset A$ も選べる。
このとき $J\cap I_n$ は $A$ の弱定義イデアルの基本系であり，
$X=\text{Spf}(A)=\colim\Spec(A/(J\cap I_n))$ を得る。これは $X$ が
可算添字付きアフィン形式代数空間であることを示す。
\end{proof}

\begin{lemma}
\label{formal-spaces-lemma-characterize-noetherian-affine}
$S$ をスキームとし，$X$ をアフィン形式代数空間とする。次は同値である：
\begin{enumerate}
\item $X$ は Noether である。
\item $X$ は adic* であり，スキーム $T$ による任意の閉埋め込み $T\to X$ に対し，
$T$ は Noether である。
\item $X$ は adic* であり，定義 \ref{formal-spaces-definition-affine-formal-algebraic-space} の
ような表示 $X=\colim X_\lambda$ のある選択に対し，スキーム $X_\lambda$ は
Noether である。
\item $X$ は弱 adic であり，定義 \ref{formal-spaces-definition-affine-formal-algebraic-space} の
ような表示 $X=\colim X_\lambda$ のある選択に対し，スキーム $X_\lambda$ は
Noether である。
\end{enumerate}
\end{lemma}

\begin{proof}
$X$ が Noether であると仮定する。このとき $X=\text{Spf}(A)$ で，$A$ は
Noether adic 環である。$T$ をスキームとし，$T\to X$ を閉埋め込みとする。
補題 \ref{formal-spaces-lemma-mcquillan-affine-formal-algebraic-space} により，$T$ はアフィンで，
$T\to\Spec(A)$ は閉埋め込みである。$A$ は Noether なので $T$ も Noether である。
このようにして (1) $\Rightarrow$ (2) が分かる。

\medskip\noindent
(2) $\Rightarrow$ (3) と (2) $\Rightarrow$ (4) は直ちに従う
（補題 \ref{formal-spaces-lemma-implications-between-types} を参照）。

\medskip\noindent
(3) $\Rightarrow$ (1) を示すため，有限生成の定義イデアル $I$ をもつ
ある adic 環 $A$ に対して $X=\text{Spf}(A)$ と書く。また，開イデアルの
ある基本系 $I_\lambda$ に対して環 $A/I_\lambda$ が Noether であることが
与えられている。$I$ は開なので，$I_\lambda\subset I$ となる $\lambda$ を
見つけられる。すると $A/I$ は Noether であり，『代数』補題
\ref{algebra-lemma-completion-Noetherian} により $A$ は Noether である。

\medskip\noindent
(4) $\Rightarrow$ (3) を示すため，ある弱 adic 環 $A$ に対して
$X=\text{Spf}(A)$ と書く。このとき $A$ は許容で，定義イデアル $I$ をもち，
$I^2$ の閉包 $I_2$ は開である。補題 \ref{formal-spaces-lemma-weakly-adic} を参照せよ。
また，開イデアルのある基本系 $I_\lambda$ に対して，環 $A/I_\lambda$ が
Noether であることも与えられている。$I_\lambda\subset I_2$ となる
$\lambda$ を選ぶ。このとき $A/I_2$ は $A/I_\lambda$ の商として Noether である。
したがって $I/I_2$ は有限 $A$-加群である。ゆえに補題
\ref{formal-spaces-lemma-weakly-adic-finite-generation} により，$A$ は有限生成の定義イデアルを
もつ adic 環である。したがって $X$ は adic* であり，(3) が成り立つ。
\end{proof}





\section{形式代数空間}
\label{formal-spaces-section-formal-algebraic-spaces}

\noindent
新しい概念をまず環について，次にスキームについて，その後に代数空間について
導入するという従来の習慣からここでは離れ，形式スキームを先に導入せずに
形式代数空間を導入する。一般的な考え方は，形式代数空間とは fppf 位相の層であり，
エタール局所的に \cite{BVGD} の意味でアフィン形式スキームとなるものだという
ことである。関連する内容は \cite{Yasuda} にも見いだせる。

\medskip\noindent
形式代数空間の定義では，『ブートストラップ』節
\ref{bootstrap-section-morphism-representable-by-spaces} および
\ref{bootstrap-section-representable-by-spaces-properties} の用語を借用する。

\begin{definition}
\label{formal-spaces-definition-formal-algebraic-space}
$S$ をスキームとする。$(\Sch/S)_{fppf}$ 上の層 $X$ が
{\it 形式代数空間}であるとは，層の写像の族
$\{X_i\to X\}_{i\in I}$ であって次を満たすものが存在することをいう：
\begin{enumerate}
\item $X_i$ はアフィン形式代数空間である。
\item $X_i\to X$ は代数空間により表現可能かつエタールである。
\item 層の写像として $\coprod X_i\to X$ は全射である。
\end{enumerate}
さらに，$X$ は集合論的条件を満たすものとする
（評注 \ref{formal-spaces-remark-set-theoretic} を参照）。$S$ 上の
{\it 形式代数空間の射}とは，層の写像をいう。
\end{definition}

\noindent
議論。確認：アフィン形式代数空間は形式代数空間である。定義の状況では，射
$X_i\to X$ は（スキームにより）表現可能である。補題
\ref{formal-spaces-lemma-presentation-representable} を参照せよ。『ブートストラップ』補題
\ref{bootstrap-lemma-surjective-flat-locally-finite-presentation} により，
$\coprod X_i\to X$ が層の写像として全射であることを要求する代わりに，単に
全射であることを要求してもよい（これは表現可能なので意味をもつ）。

\medskip\noindent
ここでいう形式代数空間は {\bf 非常に一般的}である。実際，上で定義した
アフィン形式代数空間でさえ，非常に扱いにくい対象である。

\begin{lemma}
\label{formal-spaces-lemma-diagonal-formal-algebraic-space}
$S$ をスキームとする。$X$ が $S$ 上の形式代数空間ならば，対角射
$\Delta:X\to X\times_SX$ は表現可能で，モノ射，局所準有限，局所有限型，
かつ分離である。
\end{lemma}

\begin{proof}
$U,V$ を $S$ 上のスキームとし，$U\to X$ および $V\to X$ が与えられたとする。
このとき $U\times_XV$ は層である。定義 \ref{formal-spaces-definition-formal-algebraic-space} の
ような $\{X_i\to X\}$ を選ぶ。各 $i$ に対し，射
$$
(U \times_X X_i) \times_{X_i} (V \times_X X_i)
= (U \times_X V) \times_X X_i \to U \times_X V
$$
は $X_i\to X$ の基底変換として表現可能かつエタールであり，その始域はスキーム
である（補題 \ref{formal-spaces-lemma-diagonal-affine-formal-algebraic-space} および
\ref{formal-spaces-lemma-presentation-representable} を用いよ）。これらの写像は合わせて全射なので，
『ブートストラップ』定理 \ref{bootstrap-theorem-final-bootstrap} により
$U\times_XV$ は代数空間である。射 $U\times_XV\to U\times_SV$ はモノ射である。
さらにこれは局所準有限である。実際，上に表示した射を前合成すると，合成
$$
(U \times_X X_i) \times_{X_i} (V \times_X X_i)
\to (U \times_X X_i) \times_S (V \times_X X_i)
\to U \times_S V
$$
を得る。これは閉埋め込み（補題
\ref{formal-spaces-lemma-diagonal-affine-formal-algebraic-space}）とエタール射の合成として
局所準有限である。『代数空間上の降下』補題
\ref{spaces-descent-lemma-locally-quasi-finite-etale-local-source} を参照せよ。
したがって『代数空間の射』命題
\ref{spaces-morphisms-proposition-locally-quasi-finite-separated-over-scheme}
により $U\times_XV$ はスキームである。ゆえに『代数空間』補題
\ref{spaces-lemma-representable-diagonal} により $\Delta$ は表現可能である。

\medskip\noindent
実際，上でスキームの射 $U\times_XV\to U\times_SV$ が常にモノ射かつ
局所準有限であることを示したので，『代数空間』補題
\ref{spaces-lemma-transformation-diagonal-properties} により
$\Delta:X\to X\times_SX$ はモノ射かつ局所準有限である。次に『代数空間』補題
\ref{spaces-lemma-representable-transformations-property-implication} の原理を用いれば，
$\Delta$ が分離かつ局所有限型であることが分かる。すなわち，スキームのモノ射は
分離であり（『スキーム』補題 \ref{schemes-lemma-monomorphism-separated}），
スキームの局所準有限射は局所有限型である（『射』節
\ref{morphisms-section-quasi-finite} の定義から従う）。
\end{proof}

\begin{lemma}
\label{formal-spaces-lemma-space-to-formal-space}
$S$ をスキームとする。$f:X\to Y$ を，$S$ 上の代数空間から $S$ 上の
形式代数空間への射とする。このとき $f$ は代数空間により表現可能である。
\end{lemma}

\begin{proof}
$Z$ を $S$ 上のスキームとし，射 $Z\to Y$ をとる。$X\times_YZ$ が
代数空間であることを示さなければならない。スキーム $U$ と全射エタール射
$U\to X$ を選ぶ。このとき $U\times_YZ\to X\times_YZ$ は表現可能な
全射エタール射であり（『代数空間』補題
\ref{spaces-lemma-base-change-representable-transformations-property}），
補題 \ref{formal-spaces-lemma-diagonal-formal-algebraic-space} により $U\times_YZ$ はスキームである。
したがって『ブートストラップ』定理 \ref{bootstrap-theorem-final-bootstrap} により
結論を得る。
\end{proof}

\begin{remark}
\label{formal-spaces-remark-compare-with-formal-schemes}
集合論的問題を除けば，EGA 流の形式スキームの圏（節
\ref{formal-spaces-section-formal-schemes-EGA} を参照）は，形式代数空間の圏の充満部分圏と
同値である。これを説明するため，基礎スキームを $\Spec(\mathbf{Z})$ と仮定する。
補題 \ref{formal-spaces-lemma-formal-scheme-sheaf-fppf} により，形式スキーム $\mathfrak X$ に
付随する点関手 $h_\mathfrak X$ は fppf 位相の層である。補題
\ref{formal-spaces-lemma-fully-faithful} により，対応 $\mathfrak X\mapsto h_\mathfrak X$ は
形式スキームの圏から fppf 層の圏への充満忠実埋め込みである。形式スキーム
$\mathfrak X$ が与えられたとき，$\mathfrak X_i$ がアフィン形式スキームである
開被覆 $\mathfrak X=\bigcup\mathfrak X_i$ を選ぶ。このとき評注
\ref{formal-spaces-remark-compare-with-affine-formal-schemes} により $h_{\mathfrak X_i}$ は
アフィン形式代数空間である。射 $h_{\mathfrak X_i}\to h_\mathfrak X$ は
表現可能な開埋め込みである。したがって
$\{h_{\mathfrak X_i}\to h_\mathfrak X\}$ は定義
\ref{formal-spaces-definition-formal-algebraic-space} のような族であり，$h_\mathfrak X$ は
形式代数空間である。
\end{remark}

\begin{remark}
\label{formal-spaces-remark-set-theoretic}
$S$ をスキームとし，$(\Sch/S)_{fppf}$ を『位相』定義
\ref{topologies-definition-big-small-fppf} のような大 fppf サイトとする。
定義 \ref{formal-spaces-definition-affine-formal-algebraic-space} および
\ref{formal-spaces-definition-formal-algebraic-space} における $X$ の集合論的条件として，次を採る：
$(\Sch/S)_{fppf}$ の対象 $U,R$，fppf 層の全射である射 $U\to X$，および
fppf 層の全射である射 $R\to U\times_XU$ が存在する。言い換えると，この層が
表現可能層の間の二つの写像の余等化子であることを要求する。この概念が適切に
振る舞うことを含意する観察をいくつか挙げる：
\begin{enumerate}
\item $X=\colim_{\lambda\in\Lambda}X_\lambda$ であり，この系が定義
\ref{formal-spaces-definition-affine-formal-algebraic-space} の条件 (1)，(2) を満たすとする。
このとき $U=\coprod_{\lambda\in\Lambda}X_\lambda\to X$ は fppf 層の全射である。
さらに補題 \ref{formal-spaces-lemma-diagonal-affine-formal-algebraic-space} により
$U\times_XU$ は $U\times_SU$ の閉部分スキームである。したがって $U$ が
$(\Sch/S)_{fppf}$ の対象により表現可能なら，$U\times_SU$ もそうであり
（『集合』補題 \ref{sets-lemma-what-is-in-it} を参照），集合論的条件が満たされる。
$\Lambda$ が可算ならば常にそうである。『集合』補題
\ref{sets-lemma-what-is-in-it} を参照せよ。
\item 確認。$\{X_i\to X\}_{i\in I}$ を定義
\ref{formal-spaces-definition-formal-algebraic-space} のような族（集合論的条件は上の定式化による）
とし，各 $X_i$ が実際にアフィン・スキームであると仮定する。このとき $X$ は
代数空間である。実際，$U'=\coprod X_i$ と
$R'=\coprod X_i\times_XX_j$ がその対象により表現可能となる，より大きな大 fppf
サイト $(\Sch'/S)_{fppf}$ を選べば，$X'=U'/R'$ はこの選択に対する
代数空間の圏の対象になる。次に『代数空間』補題
\ref{spaces-lemma-fully-faithful} を適用すると，$X$ が
$(\Sch/S)_{fppf}$ に対する代数空間であることが分かる。
\item $\{X_i\to X\}_{i\in I}$ を，定義
\ref{formal-spaces-definition-formal-algebraic-space} の条件 (1)，(2)，(3) を満たす層の写像の
族とする。各 $i$ に対し，$U_i\in\Ob((\Sch/S)_{fppf})$ と層の全射
$U_i\to X_i$ を選べる。したがって $I$ が大きすぎなければ（例えば可算なら），
$U=\coprod U_i\to X$ は層の全射であり，$U$ は $(\Sch/S)_{fppf}$ の対象により
表現可能である。$U\times_XU$ へ全射となる
$R\in\Ob((\Sch/S)_{fppf})$ を得るには，対角射
$\Delta:X\to X\times_SX$ が複雑すぎないと仮定すれば十分である。例えば，
$X$ の対角射が準コンパクト，すなわち $X$ が準分離ならば，これは常に可能である。
\end{enumerate}
\end{remark}

\section{被約化}
\label{formal-spaces-section-reduction}

\noindent
次の補題が示すように，任意の形式代数空間には基礎にある被約代数空間が存在する。

\begin{lemma}
\label{formal-spaces-lemma-reduction-formal-algebraic-space}
$S$ をスキームとする。$X$ を $S$ 上の形式代数空間とする。
被約代数空間 $X_{red}$ と，肥厚である表現可能射
$X_{red}\to X$ が存在する。$U$ が被約代数空間である射 $U\to X$ は，
$X_{red}$ を通って一意的に分解する。
\end{lemma}

\begin{proof}
まず $X$ がアフィン形式代数空間であると仮定する。定義
\ref{formal-spaces-definition-affine-formal-algebraic-space} のように
$X=\colim X_\lambda$ と書く。遷移射が肥厚なので，アフィン・スキーム
$X_\lambda$ の被約化はすべて互いに同型であり，これを $X_{red}$ と書く。
射 $X_{red}\to X$ は，補題 \ref{formal-spaces-lemma-covering-by-thickenings} と
肥厚の合成が肥厚であることから，表現可能な肥厚である。普遍性の確認は省略する
（『スキーム』定義 \ref{schemes-definition-reduced-induced-scheme}，
『スキーム』補題 \ref{schemes-lemma-map-into-reduction}，
『代数空間の性質』定義
\ref{spaces-properties-definition-reduced-induced-space}，および
『代数空間の性質』補題
\ref{spaces-properties-lemma-map-into-reduction} を用いよ）。

\medskip\noindent
$X$ と $\{X_i\to X\}_{i\in I}$ を定義
\ref{formal-spaces-definition-formal-algebraic-space} のようなものとする。各 $i$ に対し，
$X_{i,red}\to X_i$ を上で構成した被約化とする。$i,j\in I$ に対し，射影
$X_{i,red}\times_XX_j\to X_{i,red}$ は，（仮定により）エタールであり，
（補題 \ref{formal-spaces-lemma-presentation-representable} により）スキームの射である。
したがって $X_{i,red}\times_XX_j$ は被約である（『降下』補題
\ref{descent-lemma-reduced-local-smooth} を参照）。ゆえに射影
$X_{i,red}\times_XX_j\to X_j$ は普遍性により $X_{j,red}$ を通って分解する。
そこで
$$
R_{ij} = X_{i, red} \times_X X_j = X_{i, red} \times_X X_{j, red} =
X_i \times_X X_{j, red}
$$
を得る。実際，射 $X_{i,red}\to X_i$ は層の単射だからである。
$U=\coprod X_{i,red}$，$R=\coprod R_{ij}$ とおき，二つの射影を
$s,t:R\to U$ と書く。層として $R=U\times_XU$ であり，$s$ と $t$ は
エタールである。
% P10-E0307
すると上の観察により $(t,s):R\to U$ はエタール同値関係を定める。
したがって『代数空間』定理 \ref{spaces-theorem-presentation} により
$X_{red}=U/R$ は代数空間である。構成により図式
$$
\xymatrix{
\coprod X_{i, red} \ar[r] \ar[d] & \coprod X_i \ar[d] \\
X_{red} \ar[r] & X
}
$$
はカルテシアンである。右の垂直射は全射エタールであり，上の水平射は
表現可能な肥厚なので，『ブートストラップ』補題
\ref{bootstrap-lemma-after-fppf-sep-lqf} により $X_{red}\to X$ は表現可能である
（補題の仮定を確認するには，全射エタール射が全射，平坦，かつ局所有限表示であり，
肥厚が分離かつ局所準有限であることを用いる）。次に『代数空間』補題
\ref{spaces-lemma-descent-representable-transformations-property}
を用いれば，$X_{red}\to X$ が肥厚であることが分かる
（肥厚であることは，全射閉埋め込みであることと同値であることを用いる）。

\medskip\noindent
最後に，$U$ が $S$ 上の被約代数空間である射 $U\to X$ を考える。
各 $X_i\times_XU$ は $U$ 上エタールであり，したがって被約である
（『代数空間の性質』節
\ref{spaces-properties-section-types-properties} における被約代数空間の定義による）。
よって $X_i\times_XU\to X_i$ は $X_{i,red}$ を通って分解する。
$\{X_i\times_XU\to U\}$ はエタール被覆なので，$U\to X$ は
$X_{red}$ を通って分解する。
\end{proof}

\begin{example}
\label{formal-spaces-example-reduction-affine-formal-spectrum}
$A$ を弱許容位相環とする。この場合
$$
\text{Spf}(A)_{red} = \Spec(A/\mathfrak a)
$$
である。ここで $\mathfrak a\subset A$ は位相的冪零元全体のイデアルである。
実際，$\mathfrak a$ は根基イデアルであり（補題
\ref{formal-spaces-lemma-topologically-nilpotent}），$A$ が弱許容であることから開である。
\end{example}

\begin{lemma}
\label{formal-spaces-lemma-reduction-smooth}
$S$ をスキームとする。$f:X\to Y$ を，代数空間により表現可能かつ滑らかな
$S$ 上の形式代数空間の射とする（例えばエタール射）。このとき
$X_{red}=X\times_YY_{red}$ である。
\end{lemma}

\begin{proof}
（エタールの場合は，補題
\ref{formal-spaces-lemma-reduction-formal-algebraic-space} の証明における基礎の被約代数空間の
構成から直ちに従う。）$f$ が滑らかであると仮定する。
$X\times_YY_{red}\to Y_{red}$ は代数空間の滑らかな射であることに注意せよ。
したがって『代数空間上の降下』補題
\ref{spaces-descent-lemma-reduced-local-smooth} により
$X\times_YY_{red}$ は被約代数空間である。すると被約化の普遍性から，標準射
$X_{red}\to X\times_YY_{red}$ は同型である。
\end{proof}

\begin{lemma}
\label{formal-spaces-lemma-reduction-surjective}
$S$ をスキームとする。$f:X\to Y$ を，代数空間により表現可能な
$S$ 上の形式代数空間の射とする。このとき $f$ が『ブートストラップ』定義
\ref{bootstrap-definition-property-transformation} の意味で全射であることと，
$f_{red}:X_{red}\to Y_{red}$ が代数空間の全射であることは同値である。
\end{lemma}

\begin{proof}
省略する。
\end{proof}

\section{肥厚に沿う代数空間の余極限}
\label{formal-spaces-section-global-colimits}

\noindent
形式代数空間の特別な型として，次の補題のように，肥厚に沿う代数空間の
余フィルター的余極限として大域的に書けるものがある。後で（節
\ref{formal-spaces-section-quasi-compact-quasi-separated} において），任意の準コンパクトかつ
準分離な形式代数空間がこのような大域的余極限であることを見る。

\begin{lemma}
\label{formal-spaces-lemma-colimit-is-formal}
$S$ をスキームとする。有向集合 $\Lambda$ と，$\Lambda$ 上の代数空間の系
$(X_\lambda,f_{\lambda\mu})$ が与えられ，各
$f_{\lambda\mu}:X_\lambda\to X_\mu$ が肥厚であると仮定する。このとき
$X=\colim_{\lambda\in\Lambda}X_\lambda$ は $S$ 上の形式代数空間である。
\end{lemma}

\begin{proof}
余極限を fppf 層の圏でとるので，$X$ は層である。
$\lambda\in\Lambda$ を一つ選んで固定する。
% P10-E0308
$X_i$ が $S$ 上のアフィン・スキームであるようなエタール被覆
$\{X_{i,\lambda}\to X_\lambda\}$ を選ぶ。『代数空間の性質』補題
\ref{spaces-properties-lemma-cover-by-union-affines} を参照せよ。
各 $\mu\geq\lambda$ に対し，エタールな垂直射をもつカルテシアン図式
$$
\xymatrix{
X_{i, \lambda} \ar[r] \ar[d] & X_{i, \mu} \ar[d] \\
X_\lambda \ar[r] & X_\mu
}
$$
が存在する。『代数空間の射についてさらに』定理
\ref{spaces-more-morphisms-theorem-topological-invariance} を参照せよ
（ここでは，肥厚が定理の条件を満たす全射閉埋め込みであることも用いる）。
さらに，これらの図式は一意な同型を除いて一意である。したがって
$\mu'\geq\mu$ に対し
$X_{i,\mu}=X_\mu\times_{X_{\mu'}}X_{i,\mu'}$ である。
% P10-E0309
射 $X_{i,\mu}\to X_{i,\mu'}$ は肥厚の基底変換なので肥厚である。
各 $X_{i,\mu}$ は，『代数空間の極限』命題
\ref{spaces-limits-proposition-affine} と $X_{i,\lambda}$ がアフィンであることにより，
アフィン・スキームである。
$X_i=\colim_{\mu\geq\lambda}X_{i,\mu}$ とおく。このとき $X_i$ は
アフィン形式代数空間である。射 $X_i\to X$ はエタールである。実際，
アフィン・スキーム $U$ が与えられれば，任意の $U\to X$ は，ある
$\mu\geq\lambda$ に対して $X_\mu$ を通って分解する（詳細は省略する）。
このようにして $X$ が形式代数空間であることが分かる。
\end{proof}

\noindent
$S$ をスキームとし，$X$ を $S$ 上の形式代数空間とする。
$X$ が補題 \ref{formal-spaces-lemma-colimit-is-formal} のような大域的余極限であることを，
どのように証明または確認すればよいだろうか。そのためには，$Z$ が $S$ 上の
代数空間であり，$i$ が全射かつ閉埋め込み，言い換えれば $i$ が肥厚であるような写像
$i:Z\to X$ を探す。$i$ は代数空間により表現可能なので（補題
\ref{formal-spaces-lemma-space-to-formal-space}），これは意味をもち，前と同様に
『ブートストラップ』定義
\ref{bootstrap-definition-property-transformation} を用いることができる。

\begin{example}
\label{formal-spaces-example-david-hansen}
$(A,\mathfrak m,\kappa)$ を付値環とし，ある非零元
$\pi\in\mathfrak m$ に関して $(\pi)$-進完備であるとする。
さらに $\mathfrak m$ は有限生成でないと仮定する。一例は
$A=\mathcal{O}_{\mathbf{C}_p}$，$\pi=p$ である。ここで
$\mathcal{O}_{\mathbf{C}_p}$ は $p$-進複素数体 $\mathbf{C}_p$ の整数環であり，
これは $\mathbf{Q}_p$ の代数閉包の完備化である。別の例は
$$
A =
\left\{
\sum\nolimits_{\alpha \in \mathbf{Q},\ \alpha \geq 0} a_\alpha t^\alpha
\middle|
\begin{matrix}
a_\alpha \in \kappa \text{ かつ各 }n\text{ に対し，} \\
\text{有限個の非零な }a_\alpha
\text{ のみが }\alpha \leq n
\end{matrix}
\right\}
$$
および $\pi=t$ である。このとき $X=\text{Spf}(A)$ はアフィン形式代数空間であり，
$\Spec(\kappa)\to X$ は肥厚である。これは弱定義イデアル
$\mathfrak m\subset A$ に対応するが，これは定義イデアルではない。
\end{example}

\begin{remark}[弱定義イデアル]
\label{formal-spaces-remark-weak-ideals-of-definition}
$\mathfrak X$ を McQuillan の意味での形式スキームとする。評注
\ref{formal-spaces-remark-mcquillan} を参照せよ。$\mathfrak X$ の{\it 弱定義イデアル}とは，
任意のアフィン形式開部分スキーム $\mathfrak U\subset\mathfrak X$ に対して，
イデアル
$\mathcal{I}(\mathfrak U)\subset\mathcal{O}_\mathfrak X(\mathfrak U)$ が
弱許容位相環 $\mathcal{O}_\mathfrak X(\mathfrak U)$ の弱定義イデアルとなるような
イデアル層 $\mathcal{I}\subset\mathcal{O}_\mathfrak X$ をいう。
この条件はアフィン開被覆の各要素上で確認すれば十分である。一対一対応
$$
\{\mathfrak X\text{ に対する弱定義イデアル}\}
\leftrightarrow
\{\text{肥厚 }i : Z \to h_\mathfrak X\text{（上記のもの）}\}
$$
が存在する。この対応は，$\mathcal{I}$ に，スキーム
$Z=(\mathfrak X,\mathcal{O}_\mathfrak X/\mathcal{I})$ と
$\mathfrak X$ への明らかな射を対応させる。
{\it 弱定義イデアルの基本系}とは，任意のアフィン形式開部分スキーム
$\mathfrak U\subset\mathfrak X$ 上で，イデアル
$$
I_\lambda = \mathcal{I}_\lambda(\mathfrak U) \subset
A = \Gamma(\mathfrak U, \mathcal{O}_\mathfrak X)
$$
が弱許容位相環 $A$ の弱定義イデアルの基本系をなすような，弱定義イデアルの族
$\mathcal{I}_\lambda$ をいう。これもアフィン開被覆の各要素上で確認すれば十分である。
したがって，McQuillan 形式スキーム $\mathfrak X$ に付随する形式代数空間
$h_\mathfrak X$ が補題 \ref{formal-spaces-lemma-colimit-is-formal} のようなスキームの余極限で
あることと，$\mathfrak X$ に弱定義イデアルの基本系が存在することは同値である。
\end{remark}

\begin{remark}[定義イデアル]
\label{formal-spaces-remark-ideals-of-definition}
$\mathfrak X$ を EGA 流の形式スキームとする。$\mathfrak X$ の
{\it 定義イデアル}とは，任意のアフィン形式開部分スキーム
$\mathfrak U\subset\mathfrak X$ に対して，イデアル
$\mathcal{I}(\mathfrak U)\subset\mathcal{O}_\mathfrak X(\mathfrak U)$ が
許容位相環 $\mathcal{O}_\mathfrak X(\mathfrak U)$ の定義イデアルとなるような
イデアル層 $\mathcal{I}\subset\mathcal{O}_\mathfrak X$ をいう。
この条件はアフィン開被覆の各要素上で確認すれば十分である。評注
\ref{formal-spaces-remark-weak-ideals-of-definition} におけるような，定義イデアルと肥厚
$Z\to h_\mathfrak X$ の間の同じ対応は{\bf 得られない}。例
\ref{formal-spaces-example-david-hansen} がその一例を与える。
{\it 定義イデアルの基本系}とは，任意のアフィン形式開部分スキーム
$\mathfrak U\subset\mathfrak X$ 上で，イデアル
$$
I_\lambda = \mathcal{I}_\lambda(\mathfrak U) \subset
A = \Gamma(\mathfrak U, \mathcal{O}_\mathfrak X)
$$
が許容位相環 $A$ の定義イデアルの基本系をなすような，定義イデアルの族
$\mathcal{I}_\lambda$ をいう。これもアフィン開被覆の各要素上で確認すれば十分である。
$\mathfrak X$ が準コンパクトで，$\{\mathcal{I}_\lambda\}_{\lambda\in\Lambda}$ が
弱定義イデアルの基本系であると仮定する。$A$ が許容位相環なら，十分小さい
任意の開イデアルは定義イデアルである（すなわち，定義イデアルに含まれる
任意の開イデアルは定義イデアルである）。したがって，有限個のアフィン開被覆の
要素上で確認すれば足りることから，$\lambda$ が十分大きいとき
$\mathcal{I}_\lambda$ は定義イデアルである。評注
\ref{formal-spaces-remark-weak-ideals-of-definition} の議論を用いると，EGA 流の準コンパクト
形式スキーム $\mathfrak X$ に付随する形式代数空間 $h_\mathfrak X$ が補題
\ref{formal-spaces-lemma-colimit-is-formal} のようなスキームの余極限であることと，
$\mathfrak X$ に定義イデアルの基本系が存在することは同値である。
\end{remark}

\section{閉部分集合に沿う完備化}
\label{formal-spaces-section-completion}

\noindent
ここで用いる形式代数空間の概念は，閉部分集合に沿う完備化をとることに
よく適合している。

\begin{lemma}
\label{formal-spaces-lemma-completion-affine-is-affine-formal-algebraic-space}
$S$ をスキームとする。$X$ を $S$ 上のアフィン・スキームとする。
$T\subset|X|$ を閉部分集合とする。このとき関手
$$
(\Sch/S)_{fppf} \longrightarrow \textit{Sets},\quad
U \longmapsto \{f : U \to X \mid f(|U|) \subset T\}
$$
は McQuillan アフィン形式代数空間である。
\end{lemma}

\begin{proof}
$X=\Spec(A)$ とし，$T$ は根基イデアル $I\subset A$ に対応するとする。
$U=\Spec(B)$ を $S$ 上のアフィン・スキームとし，
$f:U\to X$ を $F(U)$ の元とする。このとき $f$ は，$B$ のすべての素イデアルが
$\varphi(I)B$ を含むような環写像 $\varphi:A\to B$ に対応する。
したがって $\varphi(I)$ のすべての元は $B$ において冪零である。『代数』補題
\ref{algebra-lemma-Zariski-topology} を参照せよ。
% P10-E0310
$J=\Ker(\varphi)$ とおくと，$I/J$ は $A/J$ の局所冪零イデアルである。
同値な言い方をすれば，$V(J)=V(I)=T$ である。すなわち，補題の関手は
$V(J)=T$ を満たすイデアル $J$ 全体にわたる余極限
$\colim\Spec(A/J)$ に等しい。したがってこの関手はアフィン形式代数空間である。
さらに，写像 $A\to A/J$ は全射であり，従って
$A^\wedge=\lim A/J\to A/J$ も全射なので，この空間は McQuillan である
（定義 \ref{formal-spaces-definition-types-affine-formal-algebraic-space} および補題
\ref{formal-spaces-lemma-mcquillan-affine-formal-algebraic-space} を参照）。
\end{proof}

\begin{lemma}
\label{formal-spaces-lemma-completion-is-formal-algebraic-space}
$S$ をスキームとする。$X$ を $S$ 上の代数空間とする。
$T\subset|X|$ を閉部分集合とする。このとき関手
$$
(\Sch/S)_{fppf} \longrightarrow \textit{Sets},\quad
U \longmapsto \{f : U \to X \mid f(|U|) \subset T\}
$$
は形式代数空間である。
\end{lemma}

\begin{proof}
この関手を $F$ と書く。$\{U_i\to U\}$ を fppf 被覆とする。このとき
$\coprod|U_i|\to|U|$ は全射である。$X$ は fppf 層なので，$F$ も fppf 層である。

\medskip\noindent
すべての $i$ に対して $X_i$ がアフィンであるようなエタール被覆
$\{g_i:X_i\to X\}$ をとる。
『代数空間の性質』補題
\ref{spaces-properties-lemma-cover-by-union-affines} を参照せよ。
射 $F\times_XX_i\to F$ はエタールであり（『代数空間』補題
\ref{spaces-lemma-base-change-representable-transformations-property} を参照），
写像 $\coprod F\times_XX_i\to F$ は層の全射である。したがって，
$F\times_XX_i$ がアフィン形式代数空間であることを示せば十分である。
$F\times_XX_i$ の $U$-値点とは，その像が閉部分集合
$g_i^{-1}(T)\subset|X_i|$ に含まれる射 $U\to X_i$ である。
ゆえに補題 \ref{formal-spaces-lemma-completion-affine-is-affine-formal-algebraic-space} から従う。
\end{proof}

\begin{definition}
\label{formal-spaces-definition-completion}
$S$ をスキームとする。$X$ を $S$ 上の代数空間とする。
$T\subset|X|$ を閉部分集合とする。補題
\ref{formal-spaces-lemma-completion-is-formal-algebraic-space} の形式代数空間を
{\it $X$ の $T$ に沿う完備化}という。
\end{definition}

\noindent
\cite[Chapter I, Section 10.8]{EGA} では完備化を表すために記号 $X_{/T}$ が
用いられており，本書でも時折この記号を用いる。$f:X\to X'$ を，スキーム
$S$ 上の代数空間の射とする。$T\subset|X|$ と $T'\subset|X'|$ が閉部分集合で，
$|f|(T)\subset T'$ を満たすと仮定する。このとき明らかに，$f$ は完備化の間の
形式代数空間の射
$$
X_{/T} \longrightarrow X'_{/T'}
$$
を定める。

\begin{lemma}
\label{formal-spaces-lemma-map-completions-representable}
$S$ をスキームとする。$f:X'\to X$ を $S$ 上の代数空間の射とする。
$T\subset|X|$ を閉部分集合とし，
$T'=|f|^{-1}(T)\subset|X'|$ とおく。このとき
$$
\xymatrix{
X'_{/T'} \ar[r] \ar[d] & X' \ar[d]^f \\
X_{/T} \ar[r] & X
}
$$
は層のカルテシアン図式である。特に，射
$X'_{/T'}\to X_{/T}$ は代数空間により表現可能である。
\end{lemma}

\begin{proof}
実際，$Y\to X$ を，その基礎空間 $|Y|$ が $T$ に写るようなスキームから
$X$ への射とする。
% P10-E0311
このとき $Y\times_XX'\to X$ は代数空間の射であり，
$|Y\times_XX'|$ は $T'$ に写る。したがって関手
$Y\times_{X_{/T}}X'_{/T'}$ は $Y\times_XX'$ により表現され，補題が従う。
\end{proof}

\begin{lemma}
\label{formal-spaces-lemma-reduction-completion}
$S$ をスキームとする。$X$ を $S$ 上の代数空間とする。
$T\subset|X|$ を閉部分集合とする。$X$ の $T$ に沿う完備化 $X_{/T}$ の被約化
$(X_{/T})_{red}$ は，$T$ に対応する $X$ の被約誘導閉部分空間 $Z$ である。
\end{lemma}

\begin{proof}
補題 \ref{formal-spaces-lemma-reduction-formal-algebraic-space}，『代数空間の性質』定義
\ref{spaces-properties-definition-reduced-induced-space}
（これは $Z$ の構成に『代数空間の性質』補題
\ref{spaces-properties-lemma-reduced-closed-subspace} を用いる），および
$X_{/T}$ の定義から従う。
% P10-E0312
実際，$Z$ と $(X_{/T})_{red}$ は同じ写像性質で特徴づけられる被約代数空間である：
始域が被約代数空間である射 $g:Y\to X$ は，$|Y|$ が
$T\subset|X|$ に写るとき，かつそのときに限り，これらを通って分解する。
\end{proof}

\begin{lemma}
\label{formal-spaces-lemma-affine-formal-completion-types}
$S$ をスキームとする。$X=\Spec(A)$ を $S$ 上のアフィン・スキームとする。
$T\subset X$ を閉部分集合とする。$X_{/T}$ を $X$ の $T$ に沿う形式完備化とする。
\begin{enumerate}
\item $X\setminus T$ が準コンパクト，すなわち $T$ が構成可能ならば，
$X_{/T}$ は adic* である。
\item ある有限生成イデアル $I\subset A$ に対して $T=V(I)$ ならば，
$X_{/T}=\text{Spf}(A^\wedge)$ である。ここで $A^\wedge$ は $A$ の $I$-進完備化である。
\item $X$ が Noether ならば，$X_{/T}$ も Noether である。
\end{enumerate}
\end{lemma}

\begin{proof}
(1) が成り立つならば，『代数』補題 \ref{algebra-lemma-qc-open} により，
(2) のようなイデアル $I\subset A$ を見つけられる。(3) が成り立つ場合にも，
(2) のようなイデアル $I\subset A$ を見つけられる。さらに，Noether 環の完備化は
『代数』補題 \ref{algebra-lemma-completion-Noetherian-Noetherian} により Noether である。
以上より，(2) を証明すれば十分である。

\medskip\noindent
(2) の証明。$T$ を切り出すイデアルを
$I=(f_1,\ldots,f_r)\subset A$ とする。$Z=\Spec(B)$ がアフィン・スキームで，
$g:Z\to X$ が集合論的に $g(Z)\subset T$ を満たす射ならば，各 $i$ に対して
$g^\sharp(f_i)$ は $B$ で冪零である。したがって，ある $n$ に対して
$I^n$ は $B$ において零に写る。ゆえに
$X_{/T}=\colim\Spec(A/I^n)=\text{Spf}(A^\wedge)$ である。
\end{proof}

\noindent
次の補題は Ofer Gabber による。

\begin{lemma}
\label{formal-spaces-lemma-completion-countably-indexed}
\begin{reference}
Ofer Gabber による2014年9月11日付電子メール。
\end{reference}
$S$ をスキームとする。$X=\Spec(A)$ を $S$ 上のアフィン・スキームとする。
$T\subset X$ を閉部分スキームとする。
\begin{enumerate}
\item 形式完備化 $X_{/T}$ が可算添字付きであり，
$T=V(f_1,f_2,f_3,\ldots)$ を満たす可算個の元
$f_1,f_2,f_3,\ldots\in A$ が存在するならば，$X_{/T}$ は adic* である。
\item $T$ が $X$ において可算個の関数により切り出せるという仮定を除くと，
(1) の結論は誤りである。
\end{enumerate}
\end{lemma}

\begin{proof}
$X_{/T}$ が可算添字付きであるという仮定は，イデアルの列
$$
A \supset J_1 \supset J_2 \supset J_3 \supset \ldots
$$
が存在して，$V(J_n)=T$ であり，
% P10-E0313
$V(J)=T$ を満たすすべてのイデアル $J\subset A$ に対して，
$J\supset J_n$ となる $n$ が存在することを意味する。

\medskip\noindent
(2) の例を構成するため，$\omega_1$ を最初の非可算順序数とする。
$k$ を体とし，$A$ を，$x_\alpha$（$\alpha\in\omega_1$）と
$y_{\alpha\beta}$（$\alpha\in\beta\in\omega_1$）により生成され，関係式
$x_\alpha=y_{\alpha\beta}x_\beta$ を課した $k$-代数とする。
$T=V(x_\alpha)$，$J_n=(x_\alpha^n)$ とおく。
$J\subset A$ が $V(J)=T$ を満たすイデアルならば，ある $n_\alpha\geq1$ に対して
$x_\alpha^{n_\alpha}\in J$ である。集合
$\{\alpha\mid n_\alpha=n\}$ のうち一つは $\omega_1$ において非有界でなければならない。
すると関係式から $J_n\subset J$ が従う。

\medskip\noindent
(2) を示すには，$A^\wedge=\lim A/J_n$ が有限生成イデアルに関して完備な環で
ないことを示せば十分である。$\gamma\in\omega_1$ に対し，$A_\gamma$ を，
$x_\alpha$（$\alpha\in\gamma$）と $y_{\alpha\beta}$（$\alpha\in\gamma$）により
生成されるイデアルによる $A$ の商とする。$A/J_1$ は被約なので，
$\lim A/J_n$ の位相的冪零元 $f$ はすべて
$J_1^\wedge=\lim J_1/J_n$ に属する。これは，$f$ が可算個の生成元のみを含む
無限級数であることを意味する。したがって，ある $\gamma$ に対して
$A_\gamma^\wedge=\lim A_\gamma/J_nA_\gamma$ における $f$ は零になる。
補題 \ref{formal-spaces-lemma-ses} により $A^\wedge\to A_\gamma^\wedge$ は連続かつ開である。
% P10-E0314
$A^\wedge$ の位相が，ある有限生成イデアル $I\subset A^\wedge$ に関する
$I$-進位相であるとすれば，ある $A_\gamma^\wedge$ において $I$ は零に写る。
すると $A_\gamma^\wedge$ は離散となるが，これは成り立たない。実際，
$\gamma=\alpha$ または $\gamma\in\alpha$ に対して
$x_\alpha\mapsto t$，$y_{\alpha\beta}\mapsto1$ で与えられる全射連続開写像
$A_\gamma^\wedge\to k[[t]]$ が存在する（開であることは補題
\ref{formal-spaces-lemma-ses} による）。

\medskip\noindent
(1) を証明する前に，まず次を示す：$I\subset A^\wedge$ が有限生成イデアルで，
その閉包 $\bar I$ が開ならば，$I=\bar I$ である。
$V(J_n^2)=T$ なので，$J_n^2\supset J_m$ を満たす $m$ が存在する。
したがって，部分列をとることにより，すべての $n$ に対して
$J_n^2\supset J_{n+1}$ と仮定してよい。
$J_n^\wedge=\lim_{k\geq n}J_n/J_k\subset A^\wedge$ とおく。
閉包 $\bar I=\bigcap(I+J_n^\wedge)$ は開なので（補題
\ref{formal-spaces-lemma-closed}），すべての $n\geq m$ に対して
$I+J_n^\wedge\supset J_m^\wedge$ を満たす $m$ が存在する。そのような $m$ を固定する。
$n\geq m+1$ に対して
$$
J_{n - 1}^\wedge I + J_{n + 1}^\wedge \supset
J_{n - 1}^\wedge (I + J_{n + 1}^\wedge) \supset
J_{n - 1}^\wedge J_m^\wedge
$$
である。実際，最初の包含は自明であり，二番目は上で示した。
\par\noindent
$J_{n-1}J_m\supset J_{n-1}^2\supset J_n$ が成り立つ。したがって，これらの包含から，
$A^\wedge$ における $J_n$ の像はイデアル
$J_{n-1}^\wedge I+J_{n+1}^\wedge$ に含まれる。このイデアルは開なので
$$
J_{n - 1}^\wedge I + J_{n + 1}^\wedge \supset J_n^\wedge.
$$
を得る。$I=(g_1,\ldots,g_t)$ と書く。$f\in J_{m+1}^\wedge$ をとる。
すべての $n\geq m+1$ に対して有効な最後の表示包含を用いると，
$c\geq0$ に関する帰納法により
$$
f = \sum f_{i, c} g_i \mod J_{m + 1+ c}^\wedge
$$
と書ける。ここで $f_{i,c}\in J_m^\wedge$ かつ
$f_{i,c}\equiv f_{i,c-1}\bmod J_{m+c}^\wedge$ である。
したがって $IJ_m^\wedge\supset J_{m+1}^\wedge$ である。
$I+J_{m+1}^\wedge\supset J_m^\wedge$ と合わせて，$I$ が開であることを得る。

\medskip\noindent
(1) の証明。$T=V(f_1,f_2,f_3,\ldots)$ と仮定する。
$I_m\subset A^\wedge$ を $f_1,\ldots,f_m$ により生成されるイデアルとする。
二つの場合に分ける。

\medskip\noindent
場合 I：ある $m$ に対して $I_m$ の閉包が開である。
このとき前段落の結果により $I_m$ は開である。構成により，任意の $n$ に対して
$(J_n)^2\supset J_{n+1}$ である。したがって $(J_n^\wedge)^2$ の閉包は
$J_{n+1}^\wedge$ を含み，従って開である。$n$ を十分大きくとると，
$A^\wedge$ の任意の二つの開イデアルの積の閉包が開であることが従う。
$I_m^k$ が開であることを，$k\ge 1$ に関する $k$ の帰納法により示そう。
$k=1$ の場合は場合 I における $m$ の仮定である。$k>1$ とし，
$I_m^{k-1}$ が開であると仮定する。このとき
$I_m^k=I_m^{k-1}\cdot I_m$ は二つの開イデアルの積なので，その閉包は開である。
しかし $I_m^k$ は有限生成なので，前段落を $I=I_m^k$ に適用すれば
$I_m^k$ 自身が開である。従って $k$ に関する帰納法を続けられる。
$I_m$ の各元は位相的冪零なので，$I_m$ は定義イデアルである。
従って $A^\wedge$ は有限生成の定義イデアルをもつ adic 環，すなわち
$X_{/T}$ は adic* である。

\medskip\noindent
場合 II：すべての $m$ に対して $I_m$ の閉包 $\bar I_m$ は開でない。
このとき $A^\wedge/\bar I_m$ の位相は離散でない。これは
$$
\Im(J_{\phi(m)} \to A/(f_1, \ldots, f_m)) \not =
\Im(J_{\phi(m) + 1} \to A/(f_1, \ldots, f_m))
$$
を満たす $\phi(m)\geq m$ を選べることを意味する。ここでは
$A^\wedge/(\bar I_m+J_n^\wedge)=A/((f_1,\ldots,f_m)+J_n)$ を用いた。
$0<m<i$ に対して $f_i^{e_i}\in J_{\phi(m)+1}$ となる指数 $e_i>0$ を選ぶ。
$J=(f_1^{e_1},f_2^{e_2},f_3^{e_3},\ldots)$ とおく。このとき $V(J)=T$ である。
すべての $n$ に対して $J\not\supset J_n$ であると主張するが，これは矛盾であり，
場合 II が起こらないことを示す。実際，$A/(f_1,\ldots,f_m)$ における $J$ の像は
$J_{\phi(m)+1}$ の像に含まれ，これは $J_m$ の像に真に含まれる。
\end{proof}

\section{ファイバー積}
\label{formal-spaces-section-fibre-products}

\noindent
形式代数空間のファイバー積に関する恒例の節である。

\begin{lemma}
\label{formal-spaces-lemma-etale-covering-by-formal-algebraic-spaces}
$S$ をスキームとする。$\{X_i\to X\}_{i\in I}$ を
$(\Sch/S)_{fppf}$ 上の層の写像の族とする。次を仮定する：
(a) $X_i$ は $S$ 上の形式代数空間である；
(b) $X_i\to X$ は代数空間により表現可能かつエタールである；
(c) $\coprod X_i\to X$ は層の全射である。
このとき $X$ は $S$ 上の形式代数空間である。
\end{lemma}

\begin{proof}
各 $i$ に対し，定義 \ref{formal-spaces-definition-formal-algebraic-space} のような
$\{X_{ij}\to X_i\}_{j\in J_i}$ を選ぶ。このとき
$\{X_{ij}\to X\}_{i\in I,j\in J_i}$ は，$X$ に対する定義
\ref{formal-spaces-definition-formal-algebraic-space} のような族である。
\end{proof}

\begin{lemma}
\label{formal-spaces-lemma-fibre-products-general}
$S$ をスキームとする。$X,Y$ を $S$ 上の形式代数空間とし，
$Z$ を，その対角射が代数空間により表現可能な層とする。
$X\to Z$ と $Y\to Z$ を層の写像とする。このとき
$X\times_ZY$ は形式代数空間である。
\end{lemma}

\begin{proof}
定義 \ref{formal-spaces-definition-formal-algebraic-space} のような
$\{X_i\to X\}$ と $\{Y_j\to Y\}$ を選ぶ。このとき
$\{X_i\times_ZY_j\to X\times_ZY\}$ は，代数空間により表現可能かつ
エタールな写像の族である。したがって補題
\ref{formal-spaces-lemma-etale-covering-by-formal-algebraic-spaces} により，$X$ と $Y$ が
アフィン形式代数空間である場合に $X\times_ZY$ が形式代数空間であることを
示せば十分である。

\medskip\noindent
$X$ と $Y$ がアフィン形式代数空間であると仮定する。定義
\ref{formal-spaces-definition-affine-formal-algebraic-space} のように
$X=\colim X_\lambda$，$Y=\colim Y_\mu$ と書く。このとき
$X\times_ZY=\colim X_\lambda\times_ZY_\mu$ である。各
$X_\lambda\times_ZY_\mu$ は代数空間である。
$\lambda\leq\lambda'$ かつ $\mu\leq\mu'$ に対し，射
$$
X_\lambda \times_Z Y_\mu \to
X_\lambda \times_Z Y_{\mu'} \to
X_{\lambda'} \times_Z Y_{\mu'}
$$
は，肥厚の基底変換の合成なので肥厚である。従って補題
\ref{formal-spaces-lemma-colimit-is-formal} を適用して結論を得る。
\end{proof}

\begin{lemma}
\label{formal-spaces-lemma-fibre-products}
$S$ をスキームとする。$S$ 上の形式代数空間の圏はファイバー積をもつ。
\end{lemma}

\begin{proof}
形式代数空間は表現可能な対角射をもつので，補題
\ref{formal-spaces-lemma-fibre-products-general} の特別な場合である。補題
\ref{formal-spaces-lemma-diagonal-formal-algebraic-space} を参照せよ。
\end{proof}

\begin{lemma}
\label{formal-spaces-lemma-reduction-fibre-products}
$S$ をスキームとする。$X\to Z$ と $Y\to Z$ を
$S$ 上の形式代数空間の射とする。このとき
$(X\times_ZY)_{red}=(X_{red}\times_{Z_{red}}Y_{red})_{red}$ である。
\end{lemma}

\begin{proof}
補題 \ref{formal-spaces-lemma-reduction-formal-algebraic-space} における被約化の普遍性から従う。
\end{proof}

\noindent
次の補題はすでに証明した（その時点ではファイバー積の存在を知らなかった）。

\begin{lemma}
\label{formal-spaces-lemma-diagonal-morphism-formal-algebraic-spaces}
$S$ をスキームとする。$f:X\to Y$ を $S$ 上の形式代数空間の射とする。
対角射 $\Delta:X\to X\times_YX$ は（スキームにより）表現可能であり，
モノ射，局所準有限，局所有限型，かつ分離である。
\end{lemma}

\begin{proof}
$T$ をスキームとし，$T\to X\times_YX$ を射とする。このとき
$$
T \times_{(X \times_Y X)} X = T \times_{(X \times_S X)} X
$$
である。従って補題 \ref{formal-spaces-lemma-diagonal-formal-algebraic-space} から直ちに従う。
\end{proof}

\section{形式代数空間の分離公理}
\label{formal-spaces-section-separation}

\noindent
本節では，形式代数空間に対する「絶対的な」分離条件を扱う。
形式代数空間の射に対する分離条件は後で論じる。

\begin{lemma}
\label{formal-spaces-lemma-characterize-quasi-separated}
$S$ をスキームとする。$X$ を $S$ 上の形式代数空間とする。次は同値である。
\begin{enumerate}
\item $X$ の被約化（補題 \ref{formal-spaces-lemma-reduction-formal-algebraic-space}）が
準分離な代数空間である。
\item 射 $U\to X$，$V\to X$ があり，$U$，$V$ が準コンパクトなスキームならば，
ファイバー積 $U\times_XV$ は準コンパクトである。
\item 射 $U\to X$，$V\to X$ があり，$U$，$V$ がアフィンならば，
ファイバー積 $U\times_XV$ は準コンパクトである。
\end{enumerate}
\end{lemma}

\begin{proof}
補題 \ref{formal-spaces-lemma-diagonal-formal-algebraic-space} により
$U\times_XV$ はスキームであることに注意せよ。
$U,V,X$ の被約化を $U_{red},V_{red},X_{red}$ と書く。このとき
$$
U_{red} \times_{X_{red}} V_{red} = U_{red} \times_X V_{red} \to U \times_X V
$$
はスキームの肥厚である。これと代数空間に対する類似の補題を念頭に置けば，
(1) と (2) の同値性は明らかである。『代数空間の性質』補題
\ref{spaces-properties-lemma-characterize-quasi-separated} を参照せよ。
(2) と (3) の同値性の証明は省略する。
\end{proof}

\begin{lemma}
\label{formal-spaces-lemma-characterize-separated}
$S$ をスキームとする。$X$ を $S$ 上の形式代数空間とする。次は同値である。
\begin{enumerate}
\item $X$ の被約化（補題 \ref{formal-spaces-lemma-reduction-formal-algebraic-space}）が
分離代数空間である。
\item 射 $U\to X$，$V\to X$ があり，$U$，$V$ がアフィンならば，
ファイバー積 $U\times_XV$ はアフィンであり，写像
$$
\mathcal{O}(U) \otimes_\mathbf{Z} \mathcal{O}(V)
\longrightarrow
\mathcal{O}(U \times_X V)
$$
は全射である。
\end{enumerate}
\end{lemma}

\begin{proof}
(2) が成り立つと仮定する。
% P10-E0315
アフィンな $U,V$ からの射 $U\to X_{red}$，$V\to X_{red}$ に
『代数空間の性質』補題
\ref{spaces-properties-lemma-characterize-quasi-separated} を適用し，
$U\times_{X_{red}}V=U\times_XV$ を用いると，$X_{red}$ は分離代数空間である。

\medskip\noindent
(1) を仮定する。$U\to X$ と $V\to X$ を (2) のような射とする。
補題 \ref{formal-spaces-lemma-diagonal-formal-algebraic-space} により
$U\times_XV$ はスキームであることに注意せよ。
$U,V,X$ の被約化を $U_{red},V_{red},X_{red}$ と書く。このとき
$$
U_{red} \times_{X_{red}} V_{red} = U_{red} \times_X V_{red} \to U \times_X V
$$
はスキームの肥厚である。従って
$(U\times_XV)_{red}=(U_{red}\times_{X_{red}}V_{red})_{red}$ である。
特に，$(U\times_XV)_{red}$ はアフィン・スキームであり，
% P10-E0315
写像
$$
\mathcal{O}(U) \otimes_\mathbf{Z} \mathcal{O}(V)
\longrightarrow
\mathcal{O}((U \times_X V)_{red})
$$
は全射である。『代数空間の性質』補題
\ref{spaces-properties-lemma-characterize-quasi-separated} を参照せよ。
すると『代数空間の極限』命題 \ref{spaces-limits-proposition-affine} により
$U\times_XV$ はアフィンである。他方，アフィン・スキームの射
$U\times_XV\to U\times V$ は合成
$$
U \times_X V = X \times_{(X \times_S X)} (U \times_S V)
\to U \times_S V \to U \times V
$$
である。第一の射はモノ射かつ局所有限型である（補題
\ref{formal-spaces-lemma-diagonal-formal-algebraic-space}）。第二の射は埋め込みである
（『スキーム』補題 \ref{schemes-lemma-fibre-product-after-map}）。
従って合成は局所有限型なモノ射である。他方，基礎にある被約アフィン・スキーム上の
写像は上で見たように閉埋め込みであり，従って普遍閉なので，この合成は整である
（『射』補題 \ref{morphisms-lemma-integral-universally-closed} を用いよ）。
したがって環写像
$$
\mathcal{O}(U) \otimes_\mathbf{Z} \mathcal{O}(V)
\longrightarrow
\mathcal{O}(U \times_X V)
$$
は有限型な整エピ射であり，従って有限，従って全射である
（『射』補題 \ref{morphisms-lemma-finite-integral} および
『代数』補題 \ref{algebra-lemma-finite-epimorphism-surjective} を用いよ）。
\end{proof}

\begin{definition}
\label{formal-spaces-definition-separated}
$S$ をスキームとする。$X$ を $S$ 上の形式代数空間とする。
\begin{enumerate}
\item 補題 \ref{formal-spaces-lemma-characterize-quasi-separated} の同値な条件を満たすとき，
$X$ は{\it 準分離}であるという。
\item 補題 \ref{formal-spaces-lemma-characterize-separated} の同値な条件を満たすとき，
$X$ は{\it 分離}であるという。
\end{enumerate}
\end{definition}

\noindent
次の補題は特に，弱許容位相環の完備化テンソル積が弱許容位相環であることを含意する。

\begin{lemma}
\label{formal-spaces-lemma-fibre-product-affines-over-separated}
$S$ をスキームとする。$X\to Z$ と $Y\to Z$ を
$S$ 上の形式代数空間の射とする。$Z$ は分離であると仮定する。
\begin{enumerate}
\item $X$ と $Y$ がアフィン形式代数空間ならば，$X\times_ZY$ もそうである。
\item $X$ と $Y$ が McQuillan アフィン形式代数空間ならば，$X\times_ZY$ もそうである。
\item $X$，$Y$，$Z$ が，それぞれ弱許容位相 $S$-代数 $A$，$B$，$C$ に対応する
McQuillan アフィン形式代数空間ならば，$X\times_ZY$ は
$A\widehat{\otimes}_CB$ に対応する。
\end{enumerate}
\end{lemma}

\begin{proof}
定義 \ref{formal-spaces-definition-affine-formal-algebraic-space} のように
$X=\colim X_\lambda$，$Y=\colim Y_\mu$ と書く。このとき
$X\times_ZY=\colim X_\lambda\times_ZY_\mu$ である。$Z$ は分離なので
これらのファイバー積はアフィンであり，(1) が成り立つ。
% P10-E0316
$X$ と $Y$ が弱許容位相 $S$-代数 $A$ と $B$ に対応し，
$X_\lambda=\Spec(A/I_\lambda)$，$Y_\mu=\Spec(B/J_\mu)$ であると仮定する。このとき
$$
X_\lambda \times_Z Y_\mu \to
X_\lambda \times Y_\mu \to \Spec(A \otimes B)
$$
は閉埋め込みである。従って補題
\ref{formal-spaces-lemma-mcquillan-affine-formal-algebraic-space} の条件の一つが成り立ち，
$X\times_ZY$ は McQuillan であると結論できる。さらに $Z$ が $C$ に対応する
McQuillan 形式代数空間ならば
$$
X_\lambda \times_Z Y_\mu = \Spec(A/I_\lambda \otimes_C B/J_\mu)
$$
である。従って $X\times_ZY$ に対応する弱許容位相環は完備化テンソル積である
（定義 \ref{formal-spaces-definition-toplogy-tensor-product} を参照）。
\end{proof}

\begin{lemma}
\label{formal-spaces-lemma-separated-from-separated}
$S$ をスキームとする。$X$ を $S$ 上の形式代数空間とする。
$U\to X$ を射とし，$U$ は $S$ 上の分離代数空間であるとする。
このとき $U\to X$ は分離である。
\end{lemma}

\begin{proof}
$U\to X$ は代数空間により表現可能なので（補題
\ref{formal-spaces-lemma-space-to-formal-space}），この主張は意味をもつ。
$T$ をスキームとし，$T\to X$ を射とする。
$U\times_XT\to T$ が分離であることを示さなければならない。
$U\times_XT\to U\times_ST$ はモノ射なので，$U\times_ST\to T$ が
分離であることを示せば十分である。これは $U\to S$ の基底変換なので従う。
上の議論では『代数空間の射』補題
\ref{spaces-morphisms-lemma-base-change-separated}，
\ref{spaces-morphisms-lemma-composition-separated}，
\ref{spaces-morphisms-lemma-monomorphism-separated}，および
\ref{spaces-morphisms-lemma-separated-implies-morphism-separated} を用いた。
\end{proof}

\section{準コンパクト形式代数空間}
\label{formal-spaces-section-quasi-compact}

\noindent
準コンパクト形式代数空間の特徴づけを与える。

\begin{lemma}
\label{formal-spaces-lemma-characterize-quasi-compact}
$S$ をスキームとする。$X$ を $S$ 上の形式代数空間とする。次は同値である。
\begin{enumerate}
\item $X$ の被約化（補題 \ref{formal-spaces-lemma-reduction-formal-algebraic-space}）が
準コンパクトな代数空間である。
\item 定義 \ref{formal-spaces-definition-formal-algebraic-space} のような
$\{X_i\to X\}_{i\in I}$ で，$I$ が有限であるものを選べる。
\item 代数空間により表現可能で，エタールかつ全射であり，$Y$ が
アフィン形式代数空間であるような射 $Y\to X$ が存在する。
\end{enumerate}
\end{lemma}

\begin{proof}
省略する。
\end{proof}

\begin{definition}
\label{formal-spaces-definition-quasi-compact}
$S$ をスキームとする。$X$ を $S$ 上の形式代数空間とする。
補題 \ref{formal-spaces-lemma-characterize-quasi-compact} の同値な条件を満たすとき，
$X$ は{\it 準コンパクト}であるという。
\end{definition}

\begin{lemma}
\label{formal-spaces-lemma-characterize-quasi-compact-morphism}
$S$ をスキームとする。$f:X\to Y$ を $S$ 上の形式代数空間の射とする。
次は同値である。
\begin{enumerate}
\item 被約化の間の誘導射 $f_{red}:X_{red}\to Y_{red}$
（補題 \ref{formal-spaces-lemma-reduction-formal-algebraic-space}）が
代数空間の準コンパクト射である。
\item 任意の準コンパクト・スキーム $T$ と射 $T\to Y$ に対し，
ファイバー積 $X\times_YT$ は準コンパクト形式代数空間である。
\item 任意のアフィン・スキーム $T$ と射 $T\to Y$ に対し，
ファイバー積 $X\times_YT$ は準コンパクト形式代数空間である。
\item 定義 \ref{formal-spaces-definition-formal-algebraic-space} のような被覆
$\{Y_j\to Y\}$ で，各 $X\times_YY_j$ が準コンパクト形式代数空間となるものが
存在する。
\end{enumerate}
\end{lemma}

\begin{proof}
省略する。
\end{proof}

\begin{definition}
\label{formal-spaces-definition-quasi-compact-morphism}
$S$ をスキームとする。$f:X\to Y$ を $S$ 上の形式代数空間の射とする。
補題 \ref{formal-spaces-lemma-characterize-quasi-compact-morphism} の同値な条件を満たすとき，
$f$ は{\it 準コンパクト}であるという。
\end{definition}

\noindent
射が代数空間により表現可能である場合（特に表現可能である場合），
これは既存の概念と一致する。

\begin{lemma}
\label{formal-spaces-lemma-quasi-compact-representable}
$S$ をスキームとする。$f:X\to Y$ を，代数空間により表現可能な
$S$ 上の形式代数空間の射とする。このとき $f$ が定義
\ref{formal-spaces-definition-quasi-compact-morphism} の意味で準コンパクトであることと，
$f$ が『ブートストラップ』定義
\ref{bootstrap-definition-property-transformation} の意味で準コンパクトであることは
同値である。
\end{lemma}

\begin{proof}
定義と補題 \ref{formal-spaces-lemma-characterize-quasi-compact-morphism} から直ちに従う。
\end{proof}

\section{準コンパクトかつ準分離な形式代数空間}
\label{formal-spaces-section-quasi-compact-quasi-separated}

\noindent
次の結果は Yasuda による。\cite[Proposition 3.32]{Yasuda} を参照せよ。

\begin{lemma}
\label{formal-spaces-lemma-structure-quasi-compact-quasi-separated}
\begin{reference}
\cite[Proposition 3.32]{Yasuda}
\end{reference}
$S$ をスキームとする。$X$ を $S$ 上の準コンパクトかつ準分離な形式代数空間とする。
このとき，有向集合 $\Lambda$ 上の代数空間の系
$(X_\lambda,f_{\lambda\mu})$ で，各
$f_{\lambda\mu}:X_\lambda\to X_\mu$ が肥厚であるものに対し，
$X=\colim X_\lambda$ と書ける。
\end{lemma}

\begin{proof}
補題 \ref{formal-spaces-lemma-characterize-quasi-compact} により，アフィン形式代数空間 $Y$ と
表現可能な全射エタール射 $Y\to X$ を選べる。定義
\ref{formal-spaces-definition-affine-formal-algebraic-space} のように
$Y=\colim Y_\lambda$ と書く。

\medskip\noindent
$\lambda\in\Lambda$ を選ぶ。補題 \ref{formal-spaces-lemma-presentation-representable} により
$Y_\lambda\times_XY$ はスキームである。$Y_\lambda\times_XY$ の被約化
（補題 \ref{formal-spaces-lemma-reduction-formal-algebraic-space}）は
$Y_{red}\times_{X_{red}}Y_{red}$ の被約化に等しい。$X$ は準分離であり，
$Y_{red}$ はアフィンなので，後者は準コンパクトである。したがって
$Y_\lambda\times_XY$ は準コンパクト・スキームである。
ゆえに，ある $\mu\geq\lambda$ が存在して，
$\text{pr}_2:Y_\lambda\times_XY\to Y$ は $Y_\mu$ を通って分解する。
補題 \ref{formal-spaces-lemma-factor-through-thickening} を参照せよ。
$Z_\lambda$ を射 $\text{pr}_2:Y_\lambda\times_XY\to Y_\mu$ の
スキーム論的像とする。これは $\mu$ の選択に依存せず，
$Z_\lambda\subset Y$ を射 $\text{pr}_2:Y_\lambda\times_XY\to Y$ の
スキーム論的像とみなすことができ，またそうみなす。
さらに $Z_\lambda$ は射 $\text{pr}_1:Y\times_XY_\lambda\to Y$ の
スキーム論的像にも等しい。実際，この射は $Z_\lambda$ の定義に用いた射と同型である。
$Y\times_XY$ の部分関手として
$Z_\lambda\times_XY=Y\times_XZ_\lambda$ であると主張する。
実際，$Y\to X$ はエタールなので，『代数空間の射』補題
\ref{spaces-morphisms-lemma-quasi-compact-scheme-theoretic-image} により，
$Z_\lambda\times_XY$ は射
$$
\text{pr}_{13} = \text{pr}_1 \times \text{id}_Y :
Y \times_X Y_\lambda \times_X Y \longrightarrow Y \times_X Y
$$
のスキーム論的像である。同様に，$Y\times_XZ_\lambda$ は射
$$
\text{pr}_{13} = \text{id}_Y \times \text{pr}_2 : 
Y \times_X Y_\lambda \times_X Y \longrightarrow Y \times_X Y
$$
のスキーム論的像である。主張が従う。すると
$R_\lambda=Z_\lambda\times_XY=Y\times_XZ_\lambda$ と射
$R_\lambda\to Z_\lambda\times_SZ_\lambda$ はエタール同値関係を定める。
このようにして代数空間 $X_\lambda=Z_\lambda/R_\lambda$ を得る。
構成により図式
$$
\xymatrix{
Z_\lambda \ar[r] \ar[d] & Y \ar[d] \\
X_\lambda \ar[r] & X
}
$$
はカルテシアンである（$X$ は二つの射影 $R=Y\times_XY\to Y$ の余等化子であり，
$Z_\lambda\subset Y$ は $R$-不変であり，$R_\lambda$ は $R$ の
$Z_\lambda$ への制限だからである）。従って $X_\lambda\to X$ は表現可能な
閉埋め込みである。『代数空間』補題
\ref{spaces-lemma-morphism-sheaves-with-P-effective-descent-etale} を参照せよ。
他方，$Y_\lambda\subset Z_\lambda$ なので，
$(X_\lambda)_{red}=X_{red}$ である。言い換えると，$X_\lambda\to X$ は肥厚である。
最後に
$$
X = \colim X_\lambda
$$
と主張する。$Y\times_XX_\lambda=Z_\lambda\supset Y_\lambda$ である。
$T$ が $S$ 上のスキームである任意の射 $T\to X$ は，エタール局所的に $Y$ への射へ
持ち上がり，さらにエタール局所的にある $Y_\lambda$ への射へ持ち上がる。
従って $T\to X$ は $T$ 上エタール局所的に $X_\lambda$ への射へ持ち上がる。
これで証明が完了する。
\end{proof}


\begin{remark}
\label{formal-spaces-remark-structure-quasi-compact-quasi-separated}
この評注では，補題 \ref{formal-spaces-lemma-structure-quasi-compact-quasi-separated} の主張と
証明を EGA 流の形式スキームの言葉に翻訳する。評注
\ref{formal-spaces-remark-ideals-of-definition} を見ると，補題は次のように翻訳できる。
\begin{itemize}
\item[(*)] 任意の準コンパクトかつ準分離な形式スキームは，
定義イデアルの基本系をもつ。
\end{itemize}
これを証明するには，まず『スキームのコホモロジー』補題
\ref{coherent-lemma-induction-principle} の帰納原理
（準コンパクトかつ準分離な形式スキームに対して言い換えたもの）を用いて，
次の状況に帰着する：
$\mathfrak X=\mathfrak U\cup\mathfrak V$ であり，$\mathfrak U$，$\mathfrak V$ は
開形式部分スキーム，$\mathfrak V$ はアフィンで，結論が
$\mathfrak U$，$\mathfrak V$，および $\mathfrak U\cap\mathfrak V$ に対して
成り立つとする。任意の定義イデアル
$\mathcal{I}\subset\mathcal{O}_\mathfrak U$ と
$\mathcal{J}\subset\mathcal{O}_\mathfrak V$ を選ぶ。
$\mathfrak U$ と $\mathfrak V$ 上に定義イデアルの基本系があるという仮定と，
$\mathfrak U\cap\mathfrak V$ が準コンパクトであることから，定義イデアル
$\mathcal{I}'\subset\mathcal{I}$ と $\mathcal{J}'\subset\mathcal{J}$ で
$$
\mathcal{I}'|_{\mathfrak U \cap \mathfrak V} \subset
\mathcal{J}|_{\mathfrak U \cap \mathfrak V}
\quad\text{および}\quad
\mathcal{J}'|_{\mathfrak U \cap \mathfrak V} \subset
\mathcal{I}|_{\mathfrak U \cap \mathfrak V}
$$
を満たすものを見つけられる。定義イデアル
$\mathcal{I}'\subset\mathcal{I}\subset\mathcal{O}_\mathfrak U$ と
$\mathcal{J}'\subset\mathcal{J}\subset\mathcal{O}_\mathfrak V$ により定まる閉埋め込みを
$U\to U'\to\mathfrak U$ および $V\to V'\to\mathfrak V$ とする。
$\mathfrak U\cap V$ を，基礎位相空間が $\mathfrak U\cap\mathfrak V$ のそれに
等しい $V$ の開部分スキームとする。$\mathcal{I}'$ の選び方により，分解
$\mathfrak U\cap V\to U'$ が存在する。同様に，$V'$ を通って分解する
$U\cap\mathfrak V$ を定義する。次に
$$
Z_U = \text{次のスキーム論的像：}
U \amalg (\mathfrak U \cap V) \longrightarrow U'
$$
および
$$
Z_V = \text{次のスキーム論的像：}
(U \cap \mathfrak V) \amalg V \longrightarrow V'
$$
を考える。準コンパクト射のスキーム論的像をとる操作は開集合への制限と可換するので
（『射』補題 \ref{morphisms-lemma-quasi-compact-scheme-theoretic-image}），
$Z_U\cap\mathfrak V=\mathfrak U\cap Z_V$ である。したがって $Z_U$ と $Z_V$ は
スキーム $Z$ に貼り合わさり，これは射 $Z\to\mathfrak X$ を備える。
評注 \ref{formal-spaces-remark-weak-ideals-of-definition} の議論と同様に，$Z$ は弱定義イデアル
$\mathcal{I}_Z\subset\mathcal{O}_\mathfrak X$ に対応する。
$Z_U\subset U'$ かつ $Z_V\subset V'$ であることに注意せよ。
従ってこの方法で構成される $\mathcal{I}_Z$ 全体の族は弱定義イデアルの基本系をなす。
ゆえに，ある部分族が定義イデアルの基本系を与える。評注
\ref{formal-spaces-remark-ideals-of-definition} を参照せよ。
\end{remark}

\begin{lemma}
\label{formal-spaces-lemma-characterize-affine}
$S$ をスキームとする。$X$ を $S$ 上の形式代数空間とする。このとき
$X$ がアフィン形式代数空間であることと，その被約化 $X_{red}$
（補題 \ref{formal-spaces-lemma-reduction-formal-algebraic-space}）がアフィンであることは同値である。
\end{lemma}

\begin{proof}
補題 \ref{formal-spaces-lemma-characterize-quasi-separated}，
\ref{formal-spaces-lemma-characterize-quasi-compact} および定義
\ref{formal-spaces-definition-separated}，\ref{formal-spaces-definition-quasi-compact} により，
$X$ は準コンパクトかつ準分離である。Yasuda の補題（補題
\ref{formal-spaces-lemma-structure-quasi-compact-quasi-separated}）により，代数空間の
肥厚のフィルター余極限 $X=\colim X_\lambda$ と書ける。しかし，各
$X_\lambda$ はアフィンである。実際，『代数空間の極限』補題
\ref{spaces-limits-lemma-reduction-scheme} と
$(X_\lambda)_{red}=X_{red}$ から従う。従って定義により $X$ は
アフィン形式代数空間である。
\end{proof}

\section{代数空間により表現可能な射}
\label{formal-spaces-section-representable}

\noindent
$f:X\to Y$ を，代数空間により表現可能な形式代数空間の射とする。
この種の射については，代数空間に関する諸章でなされた研究のおかげで，
利用できる理論が豊富にある。

\begin{lemma}
\label{formal-spaces-lemma-composition-representable}
代数空間により表現可能な射の合成は，代数空間により表現可能である。
（スキームにより）表現可能な射についても同じことが成り立つ。
\end{lemma}

\begin{proof}
『ブートストラップ』補題
\ref{bootstrap-lemma-composition-transformation} を参照せよ。
\end{proof}

\begin{lemma}
\label{formal-spaces-lemma-base-change-representable}
代数空間により表現可能な射の基底変換は，代数空間により表現可能である。
（スキームにより）表現可能な射についても同じことが成り立つ。
\end{lemma}

\begin{proof}
『ブートストラップ』補題
\ref{bootstrap-lemma-base-change-transformation} を参照せよ。
\end{proof}

\begin{lemma}
\label{formal-spaces-lemma-permanence-representable}
$S$ をスキームとする。$f:X\to Y$ および $g:Y\to Z$ を
$S$ 上の形式代数空間の射とする。$g\circ f:X\to Z$ が代数空間により
表現可能ならば，$f:X\to Y$ も代数空間により表現可能である。
\end{lemma}

\begin{proof}
$Y\to Z$ の対角射は補題
\ref{formal-spaces-lemma-diagonal-morphism-formal-algebraic-spaces} により表現可能である。
したがって『ブートストラップ』補題
\ref{bootstrap-lemma-representable-by-spaces-permanence} により，
$X\to Y$ は代数空間により表現可能である。
\end{proof}

\noindent
代数空間により表現可能であるという性質は，始域と終域の双方に関して局所的である。

\begin{lemma}
\label{formal-spaces-lemma-representable-by-algebraic-spaces-local}
$S$ をスキームとする。$f:X\to Y$ を $S$ 上の形式代数空間の射とする。
次の条件は同値である。
\begin{enumerate}
\item 射 $f$ は代数空間により表現可能である；
\item 可換図式
$$
\xymatrix{
U \ar[d] \ar[r] & V \ar[d] \\
X \ar[r] & Y
}
$$
が存在し，ここで $U$，$V$ は形式代数空間，縦の射は代数空間により
表現可能，$U\to X$ は全射エタール，かつ $U\to V$ は
代数空間により表現可能である；
\item 任意の可換図式
$$
\xymatrix{
U \ar[d] \ar[r] & V \ar[d] \\
X \ar[r] & Y
}
$$
であって，$U$，$V$ が形式代数空間で縦の射が代数空間により
表現可能なものに対して，射 $U\to V$ は代数空間により表現可能である；
\item 定義 \ref{formal-spaces-definition-formal-algebraic-space} におけるような
被覆 $\{Y_j\to Y\}$ が存在し，各 $j$ に対して定義
\ref{formal-spaces-definition-formal-algebraic-space} におけるような被覆
$\{X_{ji}\to Y_j\times_YX\}$ が存在して，各 $j$ および $i$ に対する
$X_{ji}\to Y_j$ は代数空間により表現可能である；
\item 定義 \ref{formal-spaces-definition-formal-algebraic-space} におけるような
被覆 $\{X_i\to X\}$ が存在し，各 $i$ に対して分解
$X_i\to Y_i\to Y$ が存在して，$Y_i$ はアフィン形式代数空間，
$Y_i\to Y$ は代数空間により表現可能，かつ $X_i\to Y_i$ は
代数空間により表現可能である；そして
% P10-E0318
\item ここにさらに追加する。
\end{enumerate}
\end{lemma}

\begin{proof}
(1) ならば $U=X$ および $V=Y$ と取れるので (2) であることは明らかである。
逆に，$U\to X\to Y$ に『ブートストラップ』補題
\ref{bootstrap-lemma-representable-by-spaces-cover} を適用すれば，
(2) から (1) が従う。

\medskip\noindent
(1) を仮定し，(3) におけるような図式を考える。このとき $U\to Y$ は
代数空間により表現可能であり（これは合成 $U\to X\to Y$ だからである；
『ブートストラップ』補題
\ref{bootstrap-lemma-composition-transformation} を参照せよ），
しかも $V$ を経由する。したがって補題
\ref{formal-spaces-lemma-permanence-representable} により $U\to V$ は
代数空間により表現可能である。

\medskip\noindent
(3) から (2) が従うことは明らかである。従って今や (1)--(3) は同値である。

\medskip\noindent
(4) の条件は意味をもつことに注意する。実際，補題
\ref{formal-spaces-lemma-fibre-products} によりファイバー積
$Y_j\times_YX$ は形式代数空間である。(4) から (5) が従うことは明らかである。

\medskip\noindent
(5) におけるような $X_i\to Y_i\to Y$ を仮定する。このとき
$V=\coprod Y_i$ および $U=\coprod X_i$ と置けば，(5) から (2) が従う。

\medskip\noindent
最後に (1)--(3) を仮定する。定義
\ref{formal-spaces-definition-formal-algebraic-space} におけるような任意の被覆
$\{Y_j\to Y\}$ を選び，各 $j$ に対して定義
\ref{formal-spaces-definition-formal-algebraic-space} におけるような任意の被覆
$\{X_{ji}\to Y_j\times_YX\}$ を選ぶことができる。
% P10-E0317
すると (3) により $X_{ij}\to Y_j$ は代数空間により表現可能であり，
(4) が成り立つことが分かる。これで証明が完了する。
\end{proof}

\begin{lemma}
\label{formal-spaces-lemma-algebraic-space-over-affine-formal}
$S$ をスキームとする。$Y$ を $S$ 上のアフィン形式代数空間とする。
$f:X\to Y$ を，代数空間により表現可能な
$(\Sch/S)_{fppf}$ 上の層の写像とする。このとき $X$ は形式代数空間である。
\end{lemma}

\begin{proof}
定義 \ref{formal-spaces-definition-affine-formal-algebraic-space} におけるように
$Y=\colim Y_\lambda$ と書く。各 $\lambda$ に対し，ファイバー積
$X\times_YY_\lambda$ は代数空間である。したがって補題
\ref{formal-spaces-lemma-colimit-is-formal} により
$X=\colim X\times_YY_\lambda$ は形式代数空間である。
\end{proof}

\begin{lemma}
\label{formal-spaces-lemma-representable-by-algebraic-spaces}
$S$ をスキームとする。$Y$ を $S$ 上の形式代数空間とする。
$f:X\to Y$ を，代数空間により表現可能な
$(\Sch/S)_{fppf}$ 上の層の写像とする。このとき $X$ は形式代数空間である。
\end{lemma}

\begin{proof}
$\{Y_i\to Y\}$ を定義
\ref{formal-spaces-definition-formal-algebraic-space} におけるものとする。
すると $X\times_YY_i\to X$ は，代数空間により表現可能でエタールかつ
合わせて全射な射の族である。従って補題
\ref{formal-spaces-lemma-etale-covering-by-formal-algebraic-spaces} により，
$X\times_YY_i$ が形式代数空間であることを示せば十分である。
これは補題 \ref{formal-spaces-lemma-algebraic-space-over-affine-formal} から従う。
\end{proof}

\begin{lemma}
\label{formal-spaces-lemma-affine-representable-by-algebraic-spaces}
$S$ をスキームとする。$f:X\to Y$ を，代数空間により表現可能な
アフィン形式代数空間の射とする。このとき $f$ は
（スキームにより）表現可能かつアフィンである。
\end{lemma}

\begin{proof}
$f$ がアフィンであることを示す。そうすれば『空間の射』補題
\ref{spaces-morphisms-lemma-affine-local} により，$f$ が表現可能かつ
アフィンであることが従う。定義
\ref{formal-spaces-definition-affine-formal-algebraic-space} におけるように
$Y=\colim Y_\mu$ および $X=\colim X_\lambda$ と書く。
$T$ を $S$ 上のスキームとし，射 $T\to Y$ を取る。
『ブートストラップ』定義
\ref{bootstrap-definition-property-transformation} に従って，
$X\times_YT\to T$ がアフィンであることを示さなければならない。
このためには $T$ をアフィンと仮定してよく，$X\times_YT$ が
アフィンであることを証明すればよい。この場合，補題
\ref{formal-spaces-lemma-factor-through-thickening} により $T\to Y$ はある $\mu$ に対する
$Y_\mu\to Y$ を経由する。$f$ は準コンパクトなので，補題
\ref{formal-spaces-lemma-characterize-quasi-compact-morphism} により
$X\times_YT$ は準コンパクトである。従って
$X\times_YT\to X$ はある $\lambda$ に対する $X_\lambda$ を経由する。
同様に，$\mu$ を大きくすれば $X_\lambda\to Y$ は $Y_\mu$ を経由する。
このとき
$X\times_YT=X_\lambda\times_{Y_\mu}T$ である。
アフィンスキームのファイバー積はアフィンなので結論を得る。
\end{proof}

\begin{lemma}
\label{formal-spaces-lemma-representable-affine}
$S$ をスキームとする。$\varphi:A\to B$ を，$S$ 上の弱許容位相環の
連続写像とする。次の条件は同値である。
\begin{enumerate}
\item $\text{Spf}(\varphi):\text{Spf}(B)\to\text{Spf}(A)$ は
代数空間により表現可能である；
\item $\text{Spf}(\varphi):\text{Spf}(B)\to\text{Spf}(A)$ は
（スキームにより）表現可能である；
\item $\varphi$ は taut である。定義 \ref{formal-spaces-definition-taut} を参照せよ。
\end{enumerate}
\end{lemma}

\begin{proof}
(1) と (2) は補題
\ref{formal-spaces-lemma-affine-representable-by-algebraic-spaces} により同値である。

\medskip\noindent
同値な条件 (1) と (2) が成り立つと仮定する。
$I\subset A$ を弱定義イデアルとすると，
$\Spec(A/I)\to\text{Spf}(A)$ は表現可能な肥厚である
（これは形式スペクトルの構成から明らかだが，補題
\ref{formal-spaces-lemma-mcquillan-affine-formal-algebraic-space} からも従う）。
するとその基底変換
$\Spec(A/I)\times_{\text{Spf}(A)}\text{Spf}(B)\to\text{Spf}(B)$
も表現可能な肥厚である。したがって補題
\ref{formal-spaces-lemma-mcquillan-affine-formal-algebraic-space} により，
$\text{Spf}(B)$ の部分関手として
$\Spec(A/I)\times_{\text{Spf}(A)}\text{Spf}(B)=\Spec(B/J(I))$
を満たす弱定義イデアル $J(I)\subset B$ が存在する。
Cartesian 図式
$$
\xymatrix{
\Spec(B/J(I)) \ar[d] \ar[r] & \Spec(A/I) \ar[d] \\
\text{Spf}(B) \ar[r] & \text{Spf}(A)
}
$$
を得る。補題 \ref{formal-spaces-lemma-fibre-product-affines-over-separated} により
$B/J(I)=B\widehat{\otimes}_AA/I$ であることが分かる。
従って補題 \ref{formal-spaces-lemma-closure-image-ideal} により，
$J(I)$ はイデアル $\varphi(I)B$ の閉包である。
上のような $I$ について
$\text{Spf}(A)=\colim\Spec(A/I)$ なので，
$\text{Spf}(B)=\colim\Spec(B/J(I))$ を得る。
従ってイデアル $J(I)$ は弱定義イデアルの基本系をなす
（補題 \ref{formal-spaces-lemma-mcquillan-affine-formal-algebraic-space} を参照せよ）。
ゆえに (3) が成り立つ。

\medskip\noindent
(3) を仮定する。本質的には直前の段落の議論を逆にたどるだけである。
$I\subset A$ を弱定義イデアルとする。補題
\ref{formal-spaces-lemma-fibre-product-affines-over-separated} により
Cartesian 図式
$$
\xymatrix{
\text{Spf}(B \widehat{\otimes}_A A/I) \ar[d] \ar[r] & \Spec(A/I) \ar[d] \\
\text{Spf}(B) \ar[r] & \text{Spf}(A)
}
$$
を得る。$J(I)$ を $IB$ の閉包とすると，$\varphi$ の taut 性により
$J(I)$ は $B$ において開である。
% P10-E0319
従って $J$ が $B$ において開で $J\subset J(B)$ ならば，
補題 \ref{formal-spaces-lemma-closed} により
$J(I)=\bigcap_{J \subset B\text{ は開}}(IB+J)$ であるから，
$B/J\otimes_AA/I=B/(IB+J)=B/J(I)$ である。
従って完備化テンソル積を定義する極限は停留し，
$B\widehat{\otimes}_AA/I=B/J(I)$ を与える。
ゆえに
$\text{Spf}(B\widehat{\otimes}_AA/I)=\Spec(B/J(I))$ である。
これは各弱定義イデアル $I\subset A$ に対して
$\text{Spf}(B)\times_{\text{Spf}(A)}\Spec(A/I)$ が表現可能であることを示す。
準コンパクトな $T$ をもつ任意の射 $T\to\text{Spf}(A)$ は，
ある弱定義イデアル $I$ に対する $\Spec(A/I)$ を経由するので
（補題 \ref{formal-spaces-lemma-factor-through-thickening}），
$\text{Spf}(\varphi)$ は表現可能，すなわち (2) が成り立つ。
これで証明が完了する。
\end{proof}

\begin{lemma}
\label{formal-spaces-lemma-property-goes-up-affine-morphism}
$S$ をスキームとする。$Y$ をアフィン形式代数空間とする。
$f:X\to Y$ を，表現可能かつアフィンな
$(\Sch/S)_{fppf}$ 上の層の写像とする。このとき次が成り立つ。
\begin{enumerate}
\item $X$ はアフィン形式代数空間である；
\item $Y$ が可算添字付きならば，$X$ も可算添字付きである；
\item $Y$ が可算添字付きかつ classical ならば，$X$ も可算添字付きかつ
classical である；
\item $Y$ が弱 adic ならば，$X$ も弱 adic である；
\item $Y$ が adic* ならば，$X$ も adic* である；そして
\item $Y$ が Noether で $f$ が（局所）有限型ならば，$X$ は Noether である。
\end{enumerate}
\end{lemma}

\begin{proof}
(1) の証明。定義 \ref{formal-spaces-definition-affine-formal-algebraic-space} におけるように
$Y=\colim_{\lambda\in\Lambda}Y_\lambda$ と書く。$f$ は表現可能かつ
アフィンなので，ファイバー積
$X_\lambda=Y_\lambda\times_YX$ はアフィンである。
% P10-E0320
さらに $X=\colim Y_\lambda\times_YX$ である。
従って $X$ はアフィン形式代数空間である。

\medskip\noindent
(2) の証明。$Y$ が可算添字付きならば，上の議論において
$\Lambda$ は可算であると仮定できる。このとき $X$ も可算添字付きであることが
直ちに分かる。

\medskip\noindent
(3)，(4)，(5) の証明。いずれの場合も仮定から $Y$ は可算添字付き
アフィン形式代数空間であり（補題
\ref{formal-spaces-lemma-implications-between-types}），従って (2) により $X$ もそうである。
ゆえに補題 \ref{formal-spaces-lemma-countably-indexed} により，ある弱許容位相
$S$-代数 $A$ および $B$ に対して $X=\text{Spf}(A)$ および
$Y=\text{Spf}(B)$ と書ける。補題
\ref{formal-spaces-lemma-morphism-between-formal-spectra} により，射 $f$ は連続な
$S$-代数準同型 $\varphi:B\to A$ に対応する。補題
\ref{formal-spaces-lemma-representable-affine} から $\varphi$ は taut であることが分かる。
従って (3) は補題 \ref{formal-spaces-lemma-taut-ascent-admissible} から，
(4) は補題 \ref{formal-spaces-lemma-taut-ascent-weakly-adic} から，(5) は補題
\ref{formal-spaces-lemma-taut-is-adic} から従う。

\medskip\noindent
(6) の証明。(3) と補題 \ref{formal-spaces-lemma-implications-between-types} を合わせると
$X$ は adic* であることが分かる。従って補題
\ref{formal-spaces-lemma-characterize-noetherian-affine} の判定法を使える。
まず，この判定法からアフィンスキーム $Y_\lambda$ は Noether である。
次に $X_\lambda\to Y_\lambda$ は有限型なので，$X_\lambda$ も Noether である
（『射』補題 \ref{morphisms-lemma-finite-type-noetherian}）。
再びこの判定法により $X$ は Noether である。これで証明が完了する。
\end{proof}

\begin{lemma}
\label{formal-spaces-lemma-property-goes-up-affine}
$S$ をスキームとする。$f:X\to Y$ を，代数空間により表現可能な
アフィン形式代数空間の射とする。このとき次が成り立つ。
\begin{enumerate}
\item $Y$ が可算添字付きならば，$X$ も可算添字付きである；
\item $Y$ が可算添字付きかつ classical ならば，$X$ も可算添字付きかつ
classical である；
\item $Y$ が弱 adic ならば，$X$ も弱 adic である；
\item $Y$ が adic* ならば，$X$ も adic* である；そして
\item $Y$ が Noether で $f$ が（局所）有限型ならば，$X$ は Noether である。
\end{enumerate}
\end{lemma}

\begin{proof}
補題 \ref{formal-spaces-lemma-affine-representable-by-algebraic-spaces} と
\ref{formal-spaces-lemma-property-goes-up-affine-morphism} を合わせればよい。
\end{proof}

\begin{example}
\label{formal-spaces-example-representable-morphism-from-completion}
$B$ を弱許容位相環とし，$B\to A$ を環の写像（位相は考えない）とする。
このとき
$$
A^\wedge = \lim A/JA
$$
を考えることができる。ここで極限は $B$ の弱定義イデアル $J$ 全体にわたる。
$A^\wedge$ は（極限位相を備えて）完備線形位相環となる。
全射 $A^\wedge\to A/JA$ の（開）核 $I$ は $JA^\wedge$ の閉包である
（補題 \ref{formal-spaces-lemma-closed} を参照せよ）。補題
\ref{formal-spaces-lemma-topologically-nilpotent} により，$I$ は位相的冪零元のみからなる。
従って $I$ は $A^\wedge$ の弱定義イデアルであり，$A^\wedge$ は
弱許容位相環であることが分かる。
% P10-E0321
従って $\varphi:B\to A^\wedge$ は弱許容位相環の taut な写像であり，
$$
\text{Spf}(A^\wedge) \longrightarrow \text{Spf}(B)
$$
は補題 \ref{formal-spaces-lemma-representable-affine} で研究した現象の特別な場合である。
\end{example}

\begin{remark}[警告]
\label{formal-spaces-remark-warning}
補題 \ref{formal-spaces-lemma-representable-affine}，
\ref{formal-spaces-lemma-property-goes-up-affine-morphism}，および
\ref{formal-spaces-lemma-property-goes-up-affine} の議論は，次の二つの意味で鋭い。
\begin{enumerate}
\item $A$ と $B$ が弱許容環で $\varphi:A\to B$ が連続写像ならば，
$\text{Spf}(\varphi):\text{Spf}(B)\to\text{Spf}(A)$ は一般には
表現可能でない。
\item $f:Y\to X$ がアフィン形式代数空間の表現可能な射で，
$X=\text{Spf}(A)$ が McQuillan であっても，$Y$ が McQuillan であるとは限らない。
\end{enumerate}
(1) の例として，$A=k$ を離散位相をもつ体とし，
$B=k[[t]]$ に $t$-進位相を入れればよい。
(2) の例は『例』の節
\ref{examples-section-affine-formal-algebraic-space} に与えられている。
\end{remark}

\noindent
上の警告にもかかわらず，次の結果は成り立つ。

\begin{lemma}
\label{formal-spaces-lemma-etale}
$S$ をスキームとする。$Y$ を $S$ 上の McQuillan アフィン形式代数空間，
すなわちある弱許容位相 $S$-代数 $B$ に対する
$Y=\text{Spf}(B)$ とする。このとき次の二つの圏は同値である。
\begin{enumerate}
\item 代数空間により表現可能かつエタールな，アフィン形式代数空間の射
$f:X\to Y$ の圏；
\item $A$ がエタール $B$-代数であり，$J\subset B$ が $B$ の弱定義イデアル
全体を動くとき $A^\wedge=\lim A/JA$ で与えられる形の位相 $B$-代数
$A^\wedge$ の圏。
\end{enumerate}
この同値は $A^\wedge$ を $X=\text{Spf}(A^\wedge)$ に送ることにより与えられる。
特に，(1) における任意の $X$ は McQuillan である。
\end{lemma}

\begin{proof}
$A$ をエタール $B$-代数とする。任意の開イデアル $J\subset B$ に対して
$B/J\to A/JA$ はエタールである。従って射
$\text{Spf}(A^\wedge)\to Y$ は表現可能かつエタールである。
関手 $\text{Spf}$ は補題
\ref{formal-spaces-lemma-morphism-between-formal-spectra} により充満忠実である。
証明を完了するため，次の段落で (1) における任意の $X\to Y$ が
本質像に属することを示す。

\medskip\noindent
弱定義イデアル $J_0\subset B$ を選ぶ。
$Y_0=\Spec(B/J_0)$ および $X_0=Y_0\times_YX$ と置く。
すると $X_0\to Y_0$ はアフィンスキームのエタール射である
（補題 \ref{formal-spaces-lemma-affine-representable-by-algebraic-spaces} を参照せよ）。
$X_0=\Spec(A_0)$ と書く。『代数』補題
\ref{algebra-lemma-lift-etale} により，$A_0\cong A/J_0A$ を満たす
エタール代数写像 $B\to A$ を見つけられる。
% P10-E0322
定義イデアル $J\subset J_0$ を考える。上と同様に，あるエタール環写像
$B/J\to\bar A$ に対して
$\Spec(B/J)\times_YX=\Spec(\bar A)$ と書ける。
すると $B/J\to\bar A$ と $B/J\to A/JA$ はともに，
エタール環写像 $B/J_0\to A_0$ の持ち上げである。
『代数詳論』補題
\ref{more-algebra-lemma-locally-nilpotent-henselian} により，
$J_0$ を法とする同一視を持ち上げる一意な $B/J$-代数同型
$\varphi_J:A/JA\to\bar A$ が存在する。写像 $\varphi_J$ は一意なので，
$J$ を変化させても両立する。従って
% P10-E0323
$$
X = \colim \Spec(B/J) \times_Y X = \colim \Spec(A/JA) = \text{Spf}(A)
$$
であり，補題が成り立つことが分かる。
\end{proof}

\begin{lemma}
\label{formal-spaces-lemma-etale-surjective}
補題 \ref{formal-spaces-lemma-etale} の記法と仮定の下で，
$f:X\to Y$ が $B\to A^\wedge$ に対応するとする。次の条件は同値である。
\begin{enumerate}
\item $f:X\to Y$ は全射である；
\item $B\to A$ は忠実平坦である；
\item 各弱定義イデアル $J\subset B$ に対して，環写像
$B/J\to A/JA$ は忠実平坦である；そして
\item ある弱定義イデアル $J\subset B$ に対して，環写像
$B/J\to A/JA$ は忠実平坦である。
\end{enumerate}
\end{lemma}

\begin{proof}
$J\subset B$ を弱定義イデアルとする。$J$ の各元は位相的冪零なので，
$1+J$ の各元は単元である。従って $J$ は $B$ の Jacobson 根基に含まれる
（『代数』補題 \ref{algebra-lemma-contained-in-radical}）。
ゆえに，平坦環写像 $B\to A$ が忠実平坦であることと
$B/J\to A/JA$ が忠実平坦であることは同値である
（『代数』補題 \ref{algebra-lemma-ff-rings}）。
このようにして (2)--(4) が同値であることが分かる。
(1) が成り立つならば，各弱定義イデアル $J\subset B$ に対して射
$\Spec(A/JA)=\Spec(B/J)\times_YX\to\Spec(B/J)$ は全射であり，
(3) が従う。逆に (3) を仮定する。準コンパクトな $T$ をもつ射
$T\to Y$ は，
% P10-E0324
$B$ のある定義イデアル $J$ に対する $\Spec(B/J)$ を経由する
（補題 \ref{formal-spaces-lemma-factor-through-thickening}）。
従って
$X\times_YT=\Spec(A/JA)\times_{\Spec(B/J)}T\to T$
は，全射 $\Spec(A/JA)\to\Spec(B/J)$ の基底変換として全射である。
ゆえに (1) が成り立つ。
\end{proof}

\section{形式代数空間の型}
\label{formal-spaces-section-types}

\noindent
本節では，形式代数空間が
「局所 Noether」，
「局所 adic*」，
「局所弱 adic」，
「局所可算添字付きかつ classical」，および
「局所可算添字付き」
であることを定義する。「局所 adic」，
「局所 classical」，および
「局所 McQuillan」
という型は扱わない。これらの場合については，以下の補題の類似を
どのように証明すればよいか分かっていないからである
（アフィン形式代数空間の間のエタール被覆に対して，これらの補題の
類似を証明すれば十分であろう）。

\begin{lemma}
\label{formal-spaces-lemma-iff-countably-indexed}
$S$ をスキームとする。$X\to Y$ を，代数空間により表現可能で，
全射かつ平坦なアフィン形式代数空間の射とする。このとき $X$ が
可算添字付きであることと，$Y$ が可算添字付きであることは同値である。
\end{lemma}

\begin{proof}
$X$ が可算添字付きであると仮定する。補題
\ref{formal-spaces-lemma-countable-affine-formal-algebraic-space} におけるように
$X=\colim X_n$ と書く。定義
\ref{formal-spaces-definition-affine-formal-algebraic-space} におけるように
$Y=\colim Y_\lambda$ と書く。各 $n$ に対して，
$X_n\to Y$ が $Y_{\lambda_n}$ を経由するような $\lambda_n$ を選べる
（補題 \ref{formal-spaces-lemma-factor-through-thickening} を参照せよ）。
他方，各 $\lambda$ に対してスキーム $Y_\lambda\times_YX$ はアフィンであり
（補題 \ref{formal-spaces-lemma-affine-representable-by-algebraic-spaces}），従って
$Y_\lambda\times_YX\to X$ はある $n$ に対する $X_n$ を経由する
（補題 \ref{formal-spaces-lemma-factor-through-thickening}）。図式は
$$
\xymatrix{
Y_\lambda \times_Y X \ar[r] \ar[d] & X_n \ar[r] \ar[d] & X \ar[d] \\
Y_\lambda \ar@{..>}[r] \ar@/_1pc/[rr] & Y_{\lambda_n} \ar[r] & Y
}
$$
である。点線の射が存在することを示せれば，
$Y=\colim Y_{\lambda_n}$ となり，$Y$ は可算添字付きであると結論できる。
そのため，$\mu\geq\lambda$ かつ $\mu\geq\lambda_n$ を満たす
$\mu$ を選ぶ。すると $Y_\lambda\to Y$ と
$Y_{\lambda_n}\to Y$ はともに $Y_\mu\to Y$ を経由する。
$Y_\mu=\Spec(B_\mu)$ と書き，閉部分スキーム $Y_\lambda$ が
$J\subset B_\mu$ に，閉部分スキーム $Y_{\lambda_n}$ が
$J'\subset B_\mu$ に対応するとする。示したいのは $J'\subset J$ である。
上の図式から $J'A_\mu\subset JA_\mu$ が分かる。ここで
$Y_\mu\times_YX=\Spec(A_\mu)$ である。$X\to Y$ は全射かつ平坦なので，
射 $Y_\lambda\times_YX\to Y_\lambda$ はアフィンスキームの
忠実平坦射であり，従って $B_\mu\to A_\mu$ は忠実平坦である。
ゆえに望む通り $J'\subset J$ である。

\medskip\noindent
$Y$ が可算添字付きであると仮定する。このとき補題
\ref{formal-spaces-lemma-property-goes-up-affine} により $X$ は可算添字付きである。
\end{proof}

\begin{lemma}
\label{formal-spaces-lemma-iff-cic}
$S$ をスキームとする。$X\to Y$ を，代数空間により表現可能で，
全射かつ平坦なアフィン形式代数空間の射とする。このとき
$X$ が可算添字付きかつ classical であることと，
$Y$ が可算添字付きかつ classical であることは同値である。
\end{lemma}

\begin{proof}
一方の向きの含意は補題
\ref{formal-spaces-lemma-property-goes-up-affine} で既に見た。逆向きについては，
補題 \ref{formal-spaces-lemma-iff-countably-indexed} により $X$ と $Y$ はともに
可算添字付きであると仮定してよい。従って補題
\ref{formal-spaces-lemma-countably-indexed} により，ある弱許容位相
$S$-代数 $A$ と $B$ に対して $X=\text{Spf}(A)$ および
$Y=\text{Spf}(B)$ と書ける。補題
\ref{formal-spaces-lemma-morphism-between-formal-spectra} により，射 $X\to Y$ は
連続な $S$-代数準同型 $\varphi:B\to A$ に対応する。補題
\ref{formal-spaces-lemma-representable-affine} から $\varphi$ は taut であることが分かる。
$J\subset B$ を開イデアルとし，$I\subset A$ を $JA$ の閉包とする。
補題 \ref{formal-spaces-lemma-fibre-product-affines-over-separated} および
\ref{formal-spaces-lemma-closure-image-ideal} により
$\Spec(B/J)\times_YX=\Spec(A/I)$ である。従って
$B/J\to A/I$ は忠実平坦である（$X\to Y$ が全射かつ平坦だからである）。
これは $\varphi:B\to A$ が節 \ref{formal-spaces-section-taut-descent} の状況に
あることを意味する（$A$ と $B$ の役割を入れ替える）。
補題 \ref{formal-spaces-lemma-taut-descent-weakly-admissible} により結論を得る。
\end{proof}

\begin{lemma}
\label{formal-spaces-lemma-iff-weakly-adic}
$S$ をスキームとする。$X\to Y$ を，代数空間により表現可能で，
全射かつ平坦なアフィン形式代数空間の射とする。このとき
$X$ が弱 adic であることと，$Y$ が弱 adic であることは同値である。
\end{lemma}

\begin{proof}
証明は補題 \ref{formal-spaces-lemma-iff-cic} の証明とまったく同じであるが，
最後に補題 \ref{formal-spaces-lemma-taut-descent-weakly-adic} を用いる。
\end{proof}

\begin{lemma}
\label{formal-spaces-lemma-iff-adic-star}
$S$ をスキームとする。$X\to Y$ を，代数空間により表現可能で，
全射かつ平坦なアフィン形式代数空間の射とする。このとき
$X$ が adic* であることと，$Y$ が adic* であることは同値である。
\end{lemma}

\begin{proof}
証明は補題 \ref{formal-spaces-lemma-iff-cic} の証明とまったく同じであるが，
最後に補題 \ref{formal-spaces-lemma-taut-descent-adic-star} を用いる。
\end{proof}

\begin{lemma}
\label{formal-spaces-lemma-iff-noetherian}
$S$ をスキームとする。$X\to Y$ を，代数空間により表現可能で，
全射かつ平坦かつ（局所）有限型なアフィン形式代数空間の射とする。
このとき $X$ が Noether であることと，$Y$ が Noether であることは同値である。
\end{lemma}

\begin{proof}
Noether アフィン形式代数空間は adic* であることに注意する
（補題 \ref{formal-spaces-lemma-implications-between-types} を参照せよ）。
従って補題 \ref{formal-spaces-lemma-iff-adic-star} により，$X$ と $Y$ はともに
adic* であると仮定してよい。補題
\ref{formal-spaces-lemma-characterize-noetherian-affine} の判定法を用いて
結論を得る。実際，補題
\ref{formal-spaces-lemma-countable-affine-formal-algebraic-space} におけるように
$Y=\colim Y_n$ と書き，各 $n$ に対して $X_n=Y_n\times_YX$ と置く。
このとき $X_n$ はアフィンスキームであり
（補題 \ref{formal-spaces-lemma-affine-representable-by-algebraic-spaces}），
$X=\colim X_n$ である。各射 $X_n\to Y_n$ は忠実平坦かつ有限型である。
従って，この状況で $X_n$ が Noether であることと $Y_n$ が
Noether であることは同値であるという事実から補題が従う。
下降には『代数』補題 \ref{algebra-lemma-descent-Noetherian} を，
上昇には『代数』補題 \ref{algebra-lemma-Noetherian-permanence} を参照せよ。
\end{proof}

\begin{lemma}
\label{formal-spaces-lemma-type-local}
$S$ をスキームとする。
$$
P \in
\left\{
\begin{matrix}
\text{可算添字付き},\\
\text{可算添字付きかつ classical},\\
\text{弱 adic},\ adic*,\ \text{Noether}
\end{matrix}
\right\}
$$
とする。$X$ を $S$ 上の形式代数空間とする。次の条件は同値である。
\begin{enumerate}
\item $Y$ がアフィン形式代数空間で，$f:Y\to X$ が代数空間により
表現可能かつエタールならば，$Y$ は性質 $P$ をもつ；
\item 定義 \ref{formal-spaces-definition-formal-algebraic-space} におけるような
ある $\{X_i\to X\}_{i\in I}$ に対して，各 $X_i$ は性質 $P$ をもつ。
\end{enumerate}
\end{lemma}

\begin{proof}
(1) から (2) が従うことは明らかである。(2) を仮定し，
$Y\to X$ を (1) におけるものとする。ファイバー積
$X_i\times_XY$ は形式代数空間なので（補題
\ref{formal-spaces-lemma-fibre-products-general}），定義
\ref{formal-spaces-definition-formal-algebraic-space} におけるような被覆
$\{X_{ij}\to X_i\times_XY\}$ を選べる。$Y$ は準コンパクトなので，
$$
X_{i_1 j_1} \amalg \ldots \amalg X_{i_n j_n} \longrightarrow Y
$$
が全射かつエタールとなる
$(i_1,j_1),\ldots,(i_n,j_n)$ が存在する。
このとき $X_{i_kj_k}\to X_{i_k}$ は代数空間により表現可能かつ
エタールなので，補題 \ref{formal-spaces-lemma-property-goes-up-affine} により
$X_{i_kj_k}$ は性質 $P$ をもつ。すると
$X_{i_1j_1}\amalg\ldots\amalg X_{i_nj_n}$ は性質 $P$ をもつ
アフィン形式代数空間である
（アフィン形式代数空間の有限非交和に関する小さな詳細は省略する）。
従って補題 \ref{formal-spaces-lemma-iff-countably-indexed}，
\ref{formal-spaces-lemma-iff-cic}，
\ref{formal-spaces-lemma-iff-weakly-adic}，
\ref{formal-spaces-lemma-iff-adic-star}，および
\ref{formal-spaces-lemma-iff-noetherian} のいずれかを適用して結論を得る。
\end{proof}

\noindent
前の補題により，次の定義を置く準備が整った。

\begin{definition}
\label{formal-spaces-definition-types-formal-algebraic-spaces}
$S$ をスキームとし，$X$ を $S$ 上の形式代数空間とする。
$X$ が
{\it 局所可算添字付き}，
{\it 局所可算添字付きかつ classical}，
{\it 局所弱 adic}，
{\it 局所 adic*}，または
{\it 局所 Noether}
であるとは，対応する性質について補題
\ref{formal-spaces-lemma-type-local} の同値な条件が成り立つことをいう。
\end{definition}

\noindent
局所 Noether 代数空間の閉部分集合に沿う形式完備化は，
局所 Noether 形式代数空間である。

\begin{lemma}
\label{formal-spaces-lemma-formal-completion-types}
$S$ をスキームとする。$X$ を $S$ 上の代数空間とし，
$T\subset|X|$ を閉部分集合とする。$X_{/T}$ を $X$ の $T$ に沿う
形式完備化とする。
\begin{enumerate}
\item $X\setminus T\to X$ が準コンパクトならば，
$X_{/T}$ は局所 adic* である。
\item $X$ が局所 Noether ならば，$X_{/T}$ は局所 Noether である。
\end{enumerate}
\end{lemma}

\begin{proof}
$U=\coprod U_i$ がアフィンスキームの非交和となる全射エタール射
$U\to X$ を選ぶ。『空間の性質』補題
\ref{spaces-properties-lemma-cover-by-union-affines} を参照せよ。
$T_i\subset U_i$ を $T$ の逆像とする。補題
\ref{formal-spaces-lemma-map-completions-representable} により
$X_{/T}\times_XU_i=(U_i)_{/T_i}$ である。
従って $\{(U_i)_{/T_i}\to X_{/T}\}$ は定義
\ref{formal-spaces-definition-formal-algebraic-space} におけるような被覆である。
さらに，$X\setminus T\to X$ が準コンパクトならば
$U_i\setminus T_i\to U_i$ も準コンパクトであり，$X$ が局所 Noether ならば
$U_i$ も局所 Noether である。従ってアフィンの場合，すなわち補題
\ref{formal-spaces-lemma-affine-formal-completion-types} から結論が従う。
\end{proof}

\begin{remark}[警告]
\label{formal-spaces-remark-warning-completion}
$X=\Spec(A)$ とし，$T\subset X$ を有限生成イデアル
$I\subset A$ の零点集合とする。$J=\sqrt{I}$ を $I$ の根基とする。
定義から，$X_{/T}=\text{Spf}(A^\wedge)$ であることが分かる。ここで
$A^\wedge=\lim A/I^n$ は $A$ の $I$-進完備化である。他方，
$I$-進完備化から $J$-進完備化への写像
$A^\wedge\to\lim A/J^n$ は環同型でないことがある。例として
$$
A = \bigcup\nolimits_{n \geq 1} \mathbf{C}[t^{1/n}]
$$
とし，$I=(t)$ とする。このとき $J=\mathfrak m$ は付値環
$A$ の極大イデアルで，$J^2=J$ である。従って $A$ の
$J$-進完備化は $\mathbf{C}$ であるが，$I$-進完備化は例
\ref{formal-spaces-example-david-hansen} で記述した付値環である
（特に $A\subset A^\wedge$ であることは容易に分かる）。
\end{remark}

\begin{lemma}
\label{formal-spaces-lemma-types-fibre-products}
$S$ をスキームとする。
% P10-E0325
$X\to Y$ および $Z\to Y$ を $S$ 上の形式代数空間の射とする。
このとき次が成り立つ。
\begin{enumerate}
\item $X$ と $Z$ が局所可算添字付きならば，$X\times_YZ$ は
局所可算添字付きである。
\item $X$ と $Z$ が局所可算添字付きかつ classical ならば，
$X\times_YZ$ は局所可算添字付きかつ classical である。
% P10-E0326
\item $X$ と $Z$ が弱 adic ならば，$X\times_YZ$ は弱 adic である。
\item $X$ と $Z$ が局所 adic* ならば，$X\times_YZ$ は局所 adic* である。
\item $X$ と $Z$ が局所 Noether で，$X_{red}\to Y_{red}$ が
局所有限型ならば，$X\times_YZ$ は局所 Noether である。
\end{enumerate}
\end{lemma}

\begin{proof}
定義 \ref{formal-spaces-definition-formal-algebraic-space} におけるような被覆
$\{Y_j\to Y\}$ を選ぶ。各 $j$ に対して，定義
\ref{formal-spaces-definition-formal-algebraic-space} におけるような被覆
$\{X_{ji}\to Y_j\times_YX\}$ を選ぶ。各 $j$ に対して，定義
\ref{formal-spaces-definition-formal-algebraic-space} におけるような被覆
$\{Z_{jk}\to Y_j\times_YZ\}$ を選ぶ。補題
\ref{formal-spaces-lemma-fibre-product-affines-over-separated} により
$X_{ji}\times_{Y_j}Z_{jk}$ はアフィン形式代数空間である。
従って
$$
\{X_{ji} \times_{Y_j} Z_{jk} \to X \times_Y Z\}
$$
は定義 \ref{formal-spaces-definition-formal-algebraic-space} におけるような被覆である。
% P10-E0330
ゆえに，$X$，$Y$，$Z$ がアフィン形式代数空間である場合に
(1)，(2)，(3)，(4) を証明すれば十分である。

\medskip\noindent
$X$ と $Z$ が可算添字付きであると仮定する。補題
\ref{formal-spaces-lemma-countable-affine-formal-algebraic-space} におけるように
$X=\colim X_n$ および $Z=\colim Z_m$ と書く。定義
\ref{formal-spaces-definition-affine-formal-algebraic-space} におけるように
$Y=\colim_{\lambda\in\Lambda}Y_\lambda$ と書く。各 $n$ と $m$ に対して，
$X_n\to Y$ と $Z_m\to Y$ が $Y_{\lambda_{n,m}}$ を経由するような
$\lambda_{n,m}\in\Lambda$ を見つけられる
（例えば補題 \ref{formal-spaces-lemma-factor-through-thickening} を参照せよ）。
$\lambda_0\in\Lambda$ を選ぶ。帰納的に各 $t\geq1$ に対して，
すべての $1\leq n,m\leq t$ について
$\lambda_t\geq\lambda_{n,m}$ かつ
$\lambda_t\geq\lambda_{t-1}$ となる元 $\lambda_t\in\Lambda$ を選ぶ。
$Y'=\colim Y_{\lambda_t}$ と置く。このとき $Y'\to Y$ は単射射であり，
$X\to Y$ と $Z\to Y$ は $Y'$ を経由する。従って $Y$ を $Y'$ で
置き換えてよい。すなわち，$Y$ が可算添字付きであると仮定してよい。

\medskip\noindent
$X$，$Y$，$Z$ が可算添字付きであると仮定する。補題
\ref{formal-spaces-lemma-countably-indexed} により，ある弱許容位相環 $A$，$B$，$C$ に対して
$X=\text{Spf}(A)$，$Y=\text{Spf}(B)$，$Z=\text{Spf}(C)$ と書ける。
% P10-E0327
射 $X\to Y$ と $Z\to Y$ は連続環写像 $B\to A$ と $B\to C$ により
与えられる（補題 \ref{formal-spaces-lemma-morphism-between-formal-spectra} を参照せよ）。
補題 \ref{formal-spaces-lemma-fibre-product-affines-over-separated} により
$X\times_YZ=\text{Spf}(A\widehat{\otimes}_BC)$ であり，
$A\widehat{\otimes}_BC$ は弱許容位相環であることが分かる。
特に補題 \ref{formal-spaces-lemma-completed-tensor-product} の (3) により
$X\times_YZ$ は可算添字付きである。これで (1) が証明された。

\medskip\noindent
(2) の証明。この場合 $X$ と $Z$ は可算添字付きなので，上の議論により
$X\times_YZ$ は $A\widehat{\otimes}_BC$ の形式スペクトルである。
ここで $A$ と $C$ は許容である。補題
\ref{formal-spaces-lemma-completed-tensor-product} の (2) により
$A\widehat{\otimes}_BC$ は許容である。

\medskip\noindent
(3) の証明。前と同様に，$X\times_YZ$ は
$A\widehat{\otimes}_BC$ の形式スペクトルである。ここで $A$ と $C$ は
弱 adic である。補題
\ref{formal-spaces-lemma-completed-tensor-product-weakly-adic} により
$A\widehat{\otimes}_BC$ は弱 adic である。

\medskip\noindent
(4) の証明。上と同様に論じれば，補題
\ref{formal-spaces-lemma-completed-tensor-product} の (4) から従う。

\medskip\noindent
% P10-E0328
(5) の証明。(5) を補題
\ref{formal-spaces-lemma-completed-tensor-product} の (5) から導くため，仮定が一致することを
示す必要がある。すなわち，証明の第1段落の記法において，
$X_{red}\to Y_{red}$ が局所有限型ならば，
$(X_{ji})_{red}\to(Y_j)_{red}$ は局所有限型である。
これは『空間の射』補題
\ref{spaces-morphisms-lemma-finite-type-local} と，次の可換図式
$$
\xymatrix{
(X_{ji})_{red} \ar[d] \ar[r] & (Y_j)_{red} \ar[d] \\
X_{red} \ar[r] & Y_{red}
}
$$
において縦の射がエタールであることから従う。実際，補題
\ref{formal-spaces-lemma-reduction-smooth} により
% P10-E0329
$(X_{ji})_{red}=X_{ij}\times_XX_{red}$ および
$(Y_j)_{red}=Y_j\times_YY_{red}$ である。従って上と同様に，
$X$，$Y$，$Z$ がアフィン形式代数空間，$X$，$Z$ が Noether，
かつ $X_{red}\to Y_{red}$ が有限型である場合に帰着する。
次に証明の第2段落では $Y$ を $Y'$ で置き換えたが，構成により
$Y_{red}=Y'_{red}$ なので，この置換で有限型という仮定は保たれる。
すると $X,Y,Z$ は $A,B,C$ に，$X\times_YZ$ は
$A\widehat{\otimes}_BC$ に対応し，$A$ と $C$ は Noether adic である。
最後に，被約化を取ることは位相的冪零元のイデアルで割ることに対応するので
（例 \ref{formal-spaces-example-reduction-affine-formal-spectrum}），
$X_{red}\to Y_{red}$ が有限型であるという事実は，まさに
$B/\mathfrak b\to A/\mathfrak a$ が有限型であることを意味する。
これで証明が完了する。
\end{proof}

\begin{lemma}
\label{formal-spaces-lemma-structure-locally-noetherian}
$S$ をスキームとする。$X$ を $S$ 上の局所 Noether 形式代数空間とする。
このとき，$S$ 上の局所 Noether 代数空間の有限次数の肥厚からなる系
$X_1\to X_2\to X_3\to\ldots$ に対して $X=\colim X_n$ と書ける。
ここで $X_1=X_{red}$ であり，すべての $m\geq n$ に対して
$X_n$ は $X_m$ における $X_1$ の第 $n$ 無限小近傍である。
\end{lemma}

\begin{proof}
証明の概略のみを述べ，いくつかの詳細を省略する。$X_1=X_{red}$ と置く。
部分関手 $X_n\subset X$ を次の規則で定義する。$T$ をスキームとする射
$f:T\to X$ が $X_n$ を経由するための必要十分条件は，閉埋め込み
$X_1\times_XT\to T$ のイデアル層の第 $n$ 冪が零であることである。
準連接加群の消滅は fppf 局所的に判定できるので，$X_n\subset X$ は部分層である。
$X_n\to X$ はスキームにより表現可能な閉埋め込みであり，
$X=\colim X_n$ が（fppf 層として）成り立つと主張する。
これを確かめるには $X$ 上エタール局所的に作業してよい。従って
$X=\text{Spf}(A)$ が Noether アフィン形式代数空間であると仮定してよい。
このとき $X_1=\Spec(A/\mathfrak a)$ である。ここで
$\mathfrak a\subset A$ は Noether adic 位相環 $A$ の位相的冪零元の
イデアルである。すると $X_n=\Spec(A/\mathfrak a^n)$ となり，
望む結果を得る。
\end{proof}

\section{射と連続環写像}
\label{formal-spaces-section-morphisms-rings}

\noindent
本節では，弱許容位相環と連続環準同型の圏を $\textit{WAdm}$ と記す。
次の充満部分圏を定義する。
$$
\begin{aligned}
\textit{WAdm} &\supset
\textit{WAdm}^{count} \supset
\textit{WAdm}^{cic} \\
&\supset \textit{WAdm}^{weakly\ adic} \\
&\supset \textit{WAdm}^{adic*} \supset
\textit{WAdm}^{Noeth}
\end{aligned}
$$
それぞれの対象は次の通りである。
\begin{enumerate}
\item $\textit{WAdm}^{count}$：可算な開イデアルの基本系をもつ弱許容位相環
$A$；
\item $\textit{WAdm}^{cic}$：可算な開イデアルの基本系をもつ許容位相環
$A$；
\item $\textit{WAdm}^{weakly\ adic}$：弱 adic 位相環
（節 \ref{formal-spaces-section-weakly-adic}）；
\item $\textit{WAdm}^{adic*}$：有限生成の定義イデアルをもつ adic 位相環；そして
\item $\textit{WAdm}^{Noeth}$：Noether である adic 位相環。
\end{enumerate}
これらの型の環の形式スペクトルは，
局所可算添字付き，局所可算添字付きかつ classical，
局所弱 adic，局所 adic*，および局所 Noether 形式代数空間の
基本構成要素である。

\medskip\noindent
可算添字付きアフィン形式代数空間の射と
$\textit{WAdm}^{count}$ の射との関係を簡単に復習する。
$S$ をスキームとし，$X,Y$ を可算添字付きアフィン形式代数空間とする。
$X=\text{Spf}(A)$ および
% P10-E0331
$Y=\text{Spf}(B)$ と書く。ここで $A$ と $B$ は
$\textit{WAdm}^{count}$ に属する位相 $S$-代数である
（補題 \ref{formal-spaces-lemma-countably-indexed} を参照せよ）。
補題 \ref{formal-spaces-lemma-morphism-between-formal-spectra} により，
射 $f:X\to Y$ と位相 $S$-代数の連続写像
$$
\varphi:B\longrightarrow A
$$
の間には一対一対応がある。この対応は
$f\mapsto f^\sharp$ および $\varphi\mapsto\text{Spf}(\varphi)$
で与えられる。

\medskip\noindent
$S$ をスキームとし，$f:X\to Y$ を局所可算添字付き形式代数空間の射とする。
可換図式
$$
\xymatrix{
U \ar[d] \ar[r] & V \ar[d] \\
X \ar[r] & Y
}
$$
を考える。ここで $U,V$ はアフィン形式代数空間であり，
$U\to X$ および $V\to Y$ は代数空間により表現可能かつエタールである。
定義 \ref{formal-spaces-definition-types-formal-algebraic-spaces}
（従って補題 \ref{formal-spaces-lemma-type-local}）により，
$U,V$ は可算添字付きアフィン形式代数空間である。
前段落の議論から $U\to V$ はある連続写像
$$
\varphi:B\longrightarrow A
$$
に対応する $\text{Spf}(\varphi)$ と同型であり，
ここでこれらは $\textit{WAdm}^{count}$ の位相 $S$-代数である。

\begin{lemma}
\label{formal-spaces-lemma-completion-in-sub}
$A\in\Ob(\textit{WAdm})$ とする。$A\to A'$ を環写像
（位相は考えない）とする。
$(A')^\wedge=\lim_{I\subset A\text{ w.i.d}}A'/IA'$
を，例 \ref{formal-spaces-example-representable-morphism-from-completion} で構成した
$\textit{WAdm}$ の対象とする。
\begin{enumerate}
\item $A$ が $\textit{WAdm}^{count}$ に属するならば，$(A')^\wedge$ もそうである；
\item $A$ が $\textit{WAdm}^{cic}$ に属するならば，$(A')^\wedge$ もそうである；
\item $A$ が $\textit{WAdm}^{weakly\ adic}$ に属するならば，$(A')^\wedge$ もそうである；
\item $A$ が $\textit{WAdm}^{adic*}$ に属するならば，$(A')^\wedge$ もそうである；
\item $A$ が $\textit{WAdm}^{Noeth}$ に属し，かつ $A'$ が Noether ならば，
$(A')^\wedge$ は $\textit{WAdm}^{Noeth}$ に属する。
\end{enumerate}
\end{lemma}

\begin{proof}
$A\to(A')^\wedge$ は taut であることを思い出そう。例
\ref{formal-spaces-example-representable-morphism-from-completion} の議論を参照せよ。
従って (1)，(2)，(3)，(4) はそれぞれ補題
\ref{formal-spaces-lemma-taut-ascent-countable}，
\ref{formal-spaces-lemma-taut-ascent-admissible}，
\ref{formal-spaces-lemma-taut-ascent-weakly-adic}，および
\ref{formal-spaces-lemma-taut-is-adic} から従う。
最後に $A$ が Noether かつ adic であると仮定する。(4) により
$(A')^\wedge$ は adic である。『代数』補題
\ref{algebra-lemma-completion-Noetherian-Noetherian} により
$(A')^\wedge$ は Noether である。従って (5) が成り立つ。
\end{proof}

\begin{situation}
\label{formal-spaces-situation-local-property}
$P$ を $\textit{WAdm}^{count}$ の射の性質とする。
次の可換図式を考える。
\begin{equation}
\label{formal-spaces-equation-localize}
\vcenter{
\xymatrix{
A \ar[r] & (A')^\wedge \\
B \ar[r] \ar[u]^\varphi & (B')^\wedge \ar[u]_{\varphi'}
}
}
\end{equation}
次の条件を満たすものとする。
\begin{enumerate}
\item $A$ と $B$ は $\textit{WAdm}^{count}$ の対象である；
\item $A\to A'$ と $B\to B'$ はエタール環写像である；
\item $(A')^\wedge=\lim A'/IA'$，および $(B')^\wedge=\lim B'/JB'$
であり，ここで $I\subset A$ および $J\subset B$ は
それぞれ $A$ および $B$ の弱許容な定義イデアルを走る；
\item $\varphi:B\to A$ および
$\varphi':(B')^\wedge\to(A')^\wedge$ は連続である。
\end{enumerate}
補題 \ref{formal-spaces-lemma-completion-in-sub} により，位相環
$(A')^\wedge$ と $(B')^\wedge$ は $\textit{WAdm}^{count}$ の対象である。
$P$ が局所的性質であるとは，次の公理が成り立つことをいう。
\begin{enumerate}
\item
\label{formal-spaces-item-axiom-1}
任意の図式 (\ref{formal-spaces-equation-localize}) に対して
$P(\varphi)\Rightarrow P(\varphi')$；
\item
\label{formal-spaces-item-axiom-2}
$A\to A'$ が忠実平坦である任意の図式
(\ref{formal-spaces-equation-localize}) に対して
$P(\varphi')\Rightarrow P(\varphi)$；
\item
\label{formal-spaces-item-axiom-3}
$i=1,\ldots,n$ に対する $P(B\to A_i)$ から
$P(B\to\prod_{i=1,\ldots,n}A_i)$ が従う。
\end{enumerate}
公理 (\ref{formal-spaces-item-axiom-3}) は，
$\textit{WAdm}^{count}$ が有限積をもつので意味をなす。
\end{situation}

\begin{lemma}
\label{formal-spaces-lemma-property-defines-property-morphisms}
$S$ をスキームとし，$f:X\to Y$ を $S$ 上の局所可算添字付き
形式代数空間の射とする。$P$ を $\textit{WAdm}^{count}$ の射の
局所的性質とする。次の条件は同値である。
\begin{enumerate}
\item 任意の可換図式
$$
\xymatrix{
U \ar[d] \ar[r] & V \ar[d] \\
X \ar[r] & Y
}
$$
であって，$U,V$ はアフィン形式代数空間，
$U\to X$ および $V\to Y$ は代数空間により表現可能かつエタールであるものに対し，
射 $U\to V$ は $P$ をもつ $\textit{WAdm}^{count}$ の射に対応する；
\item 定義 \ref{formal-spaces-definition-formal-algebraic-space} におけるような被覆
$\{Y_j\to Y\}$ が存在し，各 $j$ に対して定義
\ref{formal-spaces-definition-formal-algebraic-space} におけるような被覆
$\{X_{ji}\to Y_j\times_YX\}$ が存在し，各 $X_{ji}\to Y_j$ は
$P$ をもつ $\textit{WAdm}^{count}$ の射に対応する；そして
\item 定義 \ref{formal-spaces-definition-formal-algebraic-space} におけるような被覆
$\{X_i\to X\}$ が存在し，各 $i$ に対して分解
$X_i\to Y_i\to Y$ があり，$Y_i$ はアフィン形式代数空間，
$Y_i\to Y$ は代数空間により表現可能かつエタールで，
$X_i\to Y_i$ は $P$ をもつ $\textit{WAdm}^{count}$ の射に対応する。
\end{enumerate}
\end{lemma}

\begin{proof}
(1) から (2)，(2) から (3) が従うことは明らかである。(3) における
$\{X_i\to X\}$ および $X_i\to Y_i\to Y$ を仮定し，(1) における図式を取る。
$Y_i\times_YV$ は形式代数空間なので（補題
\ref{formal-spaces-lemma-fibre-products-general}），定義
\ref{formal-spaces-definition-formal-algebraic-space} におけるような被覆
$\{Y_{ij}\to Y_i\times_YV\}$ を選べる。各 $(i,j)$ に対して同様に，
定義 \ref{formal-spaces-definition-formal-algebraic-space} におけるような被覆
$\{X_{ijk}\to Y_{ij}\times_{Y_i}X_i\times_XU\}$ を選ぶ。
$U$ は準コンパクトなので，
$$
X_{i_1j_1k_1}\amalg\ldots\amalg X_{i_nj_nk_n}\longrightarrow U
$$
が全射となる $(i_1,j_1,k_1),\ldots,(i_n,j_n,k_n)$ を選べる。
各 $s=1,\ldots,n$ に対し，次の可換図式を考える。
$$
\xymatrix@C=1.5em{
& & & X_{i_sj_sk_s} \ar[ld] \ar[d] \ar[rd] \\
X \ar[d] & X_{i_s} \ar[l] \ar[d] &
X_{i_s}\times_XU \ar[l] \ar[d] & Y_{i_sj_s} \ar[ld] \ar[rd] &
X_{i_s}\times_XU \ar[d] \ar[r] &
U \ar[d] \ar[r] & X \ar[d] \\
Y & Y_{i_s} \ar[l] & Y_{i_s}\times_YV \ar[l] & &
Y_{i_s}\times_YV \ar[r] & V \ar[r] & Y
}
$$
可算添字付きアフィン形式代数空間の射について $P$ が成り立つとは，
対応する $\textit{WAdm}^{count}$ の射について成り立つこととする。
$X_{i_sj_sk_s}\to X_{i_s}$ および
$Y_{i_sj_s}\to Y_{i_s}$ はエタール環写像の完備化である
（補題 \ref{formal-spaces-lemma-etale} を参照せよ）。従って公理
(\ref{formal-spaces-item-axiom-1}) により
$P(X_{i_s}\to Y_{i_s})$ から
$P(X_{i_sj_sk_s}\to Y_{i_sj_s})$ が従う。
$Y_{i_sj_s}\to V$ もエタール環写像の完備化である
% P10-E0333
（同じ補題を参照せよ）。
図式
$$
\xymatrix{
X_{i_sj_sk_s} \ar@{=}[r] \ar[d] & X_{i_sj_sk_s} \ar[d] \\
Y_{i_sj_s} \ar[r] & V
}
$$
に公理 (\ref{formal-spaces-item-axiom-2}) を適用する
（恒等射は忠実平坦な環写像なのでこれが許される）。
\par\noindent
これにより $P(X_{i_sj_sk_s}\to V)$ が従う。
\par\noindent
公理 (\ref{formal-spaces-item-axiom-3}) により
$P(\coprod_{s=1,\ldots,n}X_{i_sj_sk_s}\to V)$ が従う。
\par\noindent
構成により $\coprod X_{i_sj_sk_s}\to U$ は全射であり，
対応する $\textit{WAdm}^{count}$ の射は忠実平坦なエタール環写像の
完備化である（補題 \ref{formal-spaces-lemma-etale-surjective} を参照せよ）。
さらに公理 (\ref{formal-spaces-item-axiom-2}) を一度適用すれば
（$B'=B$ とする），$P(U\to V)$ が成り立つ。以上で証明が完了する。
\end{proof}

\begin{remark}[adic-star に対する変種]
\label{formal-spaces-remark-variant-adic-star}
$P$ を $\textit{WAdm}^{adic*}$ の射の性質とする。
$P$ が局所的性質であるとは，公理
(\ref{formal-spaces-item-axiom-1})，(\ref{formal-spaces-item-axiom-2})，(\ref{formal-spaces-item-axiom-3}) が
Situation \ref{formal-spaces-situation-local-property} において
$\textit{WAdm}^{adic*}$ の射について成り立つことをいう。
同じ方法で，$S$ 上の局所 adic* 形式代数空間の間の射に対する
補題 \ref{formal-spaces-lemma-property-defines-property-morphisms} の変種を得る。
\end{remark}

\begin{remark}[Noether に対する変種]
\label{formal-spaces-remark-variant-Noetherian}
$P$ を $\textit{WAdm}^{Noeth}$ の射の性質とする。
$P$ が局所的性質であるとは，公理
(\ref{formal-spaces-item-axiom-1})，(\ref{formal-spaces-item-axiom-2})，(\ref{formal-spaces-item-axiom-3}) が
Situation \ref{formal-spaces-situation-local-property} において
$\textit{WAdm}^{Noeth}$ の射について成り立つことをいう。
同じ方法で，$S$ 上の局所 Noether 形式代数空間の間の射に対する
補題 \ref{formal-spaces-lemma-property-defines-property-morphisms} の変種を得る。
\end{remark}

\begin{situation}
\label{formal-spaces-situation-base-change-local-property}
$P$ を $\textit{WAdm}^{count}$ の射の局所的性質とし，
状況 \ref{formal-spaces-situation-local-property} を参照する。
$P$ が基底変換で安定であるとは，$\textit{WAdm}^{count}$ の
$B\to A$ および $B\to C$ に対して
$P(B\to A)\Rightarrow P(C\to A\widehat{\otimes}_BC)$
が成り立つことをいう。これは補題
\ref{formal-spaces-lemma-completed-tensor-product} により
$A\widehat{\otimes}_BC$ が $\textit{WAdm}^{count}$ の対象なので意味をなす。
\end{situation}

\begin{lemma}
\label{formal-spaces-lemma-base-change-property-morphisms}
$S$ をスキームとする。$P$ を $\textit{WAdm}^{count}$ の射の局所的性質で，
基底変換で安定なものとする。$f : X \to Y$ および $g : Z \to Y$ を
$S$ 上の局所可算添字付き形式代数空間の射とする。$f$ が補題
\ref{formal-spaces-lemma-property-defines-property-morphisms} の同値条件を満たすならば，
$\text{pr}_2 : X \times_Y Z \to Z$ もその条件を満たす。
\end{lemma}

\begin{proof}
定義 \ref{formal-spaces-definition-formal-algebraic-space} におけるような被覆
$\{Y_j \to Y\}$ を選ぶ。各 $j$ に対して，定義
\ref{formal-spaces-definition-formal-algebraic-space} におけるような被覆
$\{X_{ji} \to Y_j \times_Y X\}$ を選ぶ。各 $j$ に対して，定義
\ref{formal-spaces-definition-formal-algebraic-space} におけるような被覆
$\{Z_{jk} \to Y_j \times_Y Z\}$ を選ぶ。
$X_{ji} \times_{Y_j} Z_{jk}$ は可算添字付きアフィン形式代数空間である
（補題 \ref{formal-spaces-lemma-types-fibre-products} を参照せよ）。従って
$$
\{X_{ji} \times_{Y_j} Z_{jk} \to X \times_Y Z\}
$$
は定義 \ref{formal-spaces-definition-formal-algebraic-space} における被覆である。
さらに，射 $X_{ji} \times_{Y_j} Z_{jk} \to Z$ は $Z_{jk}$ を経由する。
仮定により，$X_{ji} \to Y_j$ は性質 $P$ をもつ
$\text{WAdm}^{count}$ の射 $B_j \to A_{ji}$ に対応する。
射 $Z_{jk} \to Y_j$ は $\text{WAdm}^{count}$ における射
$B_j \to C_{jk}$ に対応する。補題
\ref{formal-spaces-lemma-fibre-product-affines-over-separated} により
$X_{ji} \times_{Y_j} Z_{jk} =
\text{Spf}(A_{ji} \widehat{\otimes}_{B_j} C_{jk})$
であるから，$C_{jk} \to A_{ji} \widehat{\otimes}_{B_j} C_{jk}$ が
性質 $P$ をもつことを示せば十分である。これはまさに，$P$ が基底変換で
安定であるという条件が保証することである。
\end{proof}

\begin{remark}[adic-star に対する変種]
\label{formal-spaces-remark-base-change-variant-adic-star}
$P$ を $\textit{WAdm}^{adic*}$ の射の局所的性質とし，注記
\ref{formal-spaces-remark-variant-adic-star} を参照する。$P$ が
{\it 基底変換に関して安定} であるとは，$\textit{WAdm}^{adic*}$ における
$B \to A$ および $B \to C$ に対して
$P(B \to A) \Rightarrow P(C \to A \widehat{\otimes}_B C)$
が成り立つことをいう。これは補題
\ref{formal-spaces-lemma-completed-tensor-product} により
$A \widehat{\otimes}_B C$ が $\textit{WAdm}^{adic*}$ の対象なので意味をなす。
全く同じ方法で，$S$ 上の局所 adic* 形式代数空間の間の射に対する補題
\ref{formal-spaces-lemma-base-change-property-morphisms} の変種を得る。
\end{remark}

\begin{remark}[Noether に対する変種]
\label{formal-spaces-remark-base-change-variant-Noetherian}
$P$ を $\textit{WAdm}^{Noeth}$ の射の局所的性質とし，注記
\ref{formal-spaces-remark-variant-Noetherian} を参照する。$P$ が基底変換に関して安定であるとは，
$\textit{WAdm}^{Noeth}$ における $B \to A$ および $B \to C$ に対して，
$P(B \to A)$ から $A \widehat{\otimes}_B C$ が adic Noether であり，さらに
$P(C \to A \widehat{\otimes}_B C)$ が成り立つことが従うことをいう。
\footnote{判定基準については補題 \ref{formal-spaces-lemma-completed-tensor-product} を参照せよ。}
全く同じ方法で，$S$ 上の局所 Noether 形式代数空間の間の射に対する補題
\ref{formal-spaces-lemma-base-change-property-morphisms} の変種を得る。
\end{remark}

\begin{remark}[Noether に関する別の変種]
\label{formal-spaces-remark-base-change-variant-variant-Noetherian}
$P$ および $Q$ を $\textit{WAdm}^{Noeth}$ の射の局所的性質とし，注記
\ref{formal-spaces-remark-variant-Noetherian} を参照する。{\it $P$ が $Q$ による基底変換で
安定} であるとは，$\textit{WAdm}^{Noeth}$ における $B \to A$ および
$B \to C$ が $P(B \to A)$ および $Q(B \to C)$ を満たすとき，
$A \widehat{\otimes}_B C$ が adic Noether であり，
$P(C \to A \widehat{\otimes}_B C)$ が成り立つことをいう。
補題 \ref{formal-spaces-lemma-base-change-property-morphisms} の証明と全く同様に論じると，
次の主張を得る：$S$ 上の局所 Noether 形式代数空間の射
$f : X \to Y$ および $g : Y \to Z$ が，
\begin{enumerate}
\item $f$ と $P$ について補題 \ref{formal-spaces-lemma-property-defines-property-morphisms}
の同値条件を満たし，
\item $g$ と $Q$ について補題 \ref{formal-spaces-lemma-property-defines-property-morphisms}
の同値条件を満たすならば，
\end{enumerate}
$P$ について $\text{pr}_2 : X \times_Y Z \to Z$ も補題
\ref{formal-spaces-lemma-property-defines-property-morphisms} の同値条件を満たす。
\end{remark}

\begin{situation}
\label{formal-spaces-situation-composition-local-property}
$P$ を $\textit{WAdm}^{count}$ の射の局所的性質とし，状況
\ref{formal-spaces-situation-local-property} を参照する。$P$ が {\it 合成に関して安定} であるとは，
$\textit{WAdm}^{count}$ における $B \to A$ および $C \to B$ に対して
$P(B \to A) \wedge P(C \to B) \Rightarrow P(C \to A)$
が成り立つことをいう。
\end{situation}

\begin{lemma}
\label{formal-spaces-lemma-composition-property-morphisms}
$S$ をスキームとする。$P$ を，合成に関して安定な
$\textit{WAdm}^{count}$ の射の局所的性質とする。$f : X \to Y$ および
$g : Y \to Z$ を $S$ 上の局所可算添字付き形式代数空間の射とする。$f$ と $g$ が
補題 \ref{formal-spaces-lemma-property-defines-property-morphisms} の同値条件を満たすならば，
$g \circ f : X \to Z$ もそうである。
\end{lemma}

\begin{proof}
定義 \ref{formal-spaces-definition-formal-algebraic-space} におけるような被覆
$\{Z_k \to Z\}$ を選ぶ。各 $k$ に対して，定義
\ref{formal-spaces-definition-formal-algebraic-space} におけるような被覆
$\{Y_{kj} \to Z_k \times_Z Y\}$ を選ぶ。各 $k$ および $j$ に対して，定義
\ref{formal-spaces-definition-formal-algebraic-space} におけるような被覆
$\{X_{kji} \to Y_{kj} \times_Y X\}$ を選ぶ。
% P10-E0336
$f$ と $g$ が補題 \ref{formal-spaces-lemma-property-defines-property-morphisms} の同値条件を
満たすならば，射
% P10-E0337
$X_{kji} \to Y_{jk}$ および $Y_{jk} \to Z_k$ は，性質 $P$ をもつ
$\text{WAdm}^{count}$ の射
$B_{kj} \to A_{kji}$ および $C_k \to B_{kj}$ に対応する。
% P10-E0338
従ってそれらの合成も性質をもち，結論を得る。
\end{proof}

\begin{remark}[adic-star に対する変種]
\label{formal-spaces-remark-composition-variant-adic-star}
$P$ を $\textit{WAdm}^{adic*}$ の射の局所的性質とし，注記
\ref{formal-spaces-remark-variant-adic-star} を参照する。$P$ が {\it 合成に関して安定} であるとは，
$\textit{WAdm}^{adic*}$ における $B \to A$ および $C \to B$ に対して
$P(B \to A) \wedge P(C \to B) \Rightarrow P(C \to A)$
が成り立つことをいう。全く同じ方法で，$S$ 上の局所 adic* 形式代数空間の間の
射に対する補題 \ref{formal-spaces-lemma-composition-property-morphisms} の変種を得る。
\end{remark}

\begin{remark}[Noether に対する変種]
\label{formal-spaces-remark-composition-variant-Noetherian}
$P$ を $\textit{WAdm}^{Noeth}$ の射の局所的性質とし，注記
\ref{formal-spaces-remark-variant-Noetherian} を参照する。$P$ が {\it 合成に関して安定} であるとは，
$\textit{WAdm}^{Noeth}$ における $B \to A$ および $C \to B$ に対して
$P(B \to A) \wedge P(C \to B) \Rightarrow P(C \to A)$
が成り立つことをいう。全く同じ方法で，$S$ 上の局所 Noether 形式代数空間の間の
射に対する補題 \ref{formal-spaces-lemma-composition-property-morphisms} の変種を得る。
\end{remark}

\begin{situation}
\label{formal-spaces-situation-permanence-local-property}
$P$ を $\textit{WAdm}^{count}$ の射の局所的性質とし，状況
\ref{formal-spaces-situation-local-property} を参照する。$P$ が {\it 消去性質をもつ} とは，
$\textit{WAdm}^{count}$ における $B \to A$ および $C \to B$ に対して
$P(C \to B) \wedge P(C \to A) \Rightarrow P(B \to A)$
が成り立つことをいう。
\end{situation}

\begin{lemma}
\label{formal-spaces-lemma-permanence-property-morphisms}
$S$ をスキームとする。$P$ を消去性質をもつ
$\textit{WAdm}^{count}$ の射の局所的性質とする。$f : X \to Y$ および
$g : Y \to Z$ を $S$ 上の局所可算添字付き形式代数空間の射とする。$g \circ f$ と $g$
% P10-E0339
が補題 \ref{formal-spaces-lemma-property-defines-property-morphisms} の同値条件を満たすならば，
$f : X \to Y$ もそうである。
\end{lemma}

\begin{proof}
定義 \ref{formal-spaces-definition-formal-algebraic-space} におけるような被覆
$\{Z_k \to Z\}$ を選ぶ。各 $k$ に対して，定義
\ref{formal-spaces-definition-formal-algebraic-space} におけるような被覆
$\{Y_{kj} \to Z_k \times_Z Y\}$ を選ぶ。各 $k$ および $j$ に対して，定義
\ref{formal-spaces-definition-formal-algebraic-space} におけるような被覆
$\{X_{kji} \to Y_{kj} \times_Y X\}$ を選ぶ。
$\text{WAdm}^{count}$ の射 $B_{kj} \to A_{kji}$ および
$C_k \to B_{kj}$ に，それぞれ $X_{kji} \to Y_{jk}$ および
% P10-E0340
$Y_{jk} \to Z_k$ を対応させる。
% P10-E0341
$g \circ f$ と $g$ が補題 \ref{formal-spaces-lemma-property-defines-property-morphisms} の同値条件を
満たすならば，$C_k \to B_{kj}$ および $C_k \to A_{kji}$ は性質 $P$ を満たす。
従って $B_{kj} \to A_{kji}$ もそうであり，結論を得る。
\end{proof}

\begin{remark}[adic-star に対する変種]
\label{formal-spaces-remark-permanence-variant-adic-star}
$P$ を $\textit{WAdm}^{adic*}$ の射の局所的性質とし，注記
\ref{formal-spaces-remark-variant-adic-star} を参照する。$P$ が {\it 消去性質をもつ} とは，
$\textit{WAdm}^{adic*}$ における $B \to A$ および $C \to B$ に対して
$P(C \to A) \wedge P(C \to B) \Rightarrow P(B \to A)$
が成り立つことをいう。全く同じ方法で，$S$ 上の局所 adic* 形式代数空間の間の
射に対する補題
% P10-E0342
\ref{formal-spaces-lemma-composition-property-morphisms} の変種を得る。
\end{remark}

\begin{remark}[Noether に対する変種]
\label{formal-spaces-remark-permanence-variant-Noetherian}
$P$ を $\textit{WAdm}^{Noeth}$ の射の局所的性質とし，注記
\ref{formal-spaces-remark-variant-Noetherian} を参照する。$P$ が {\it 消去性質をもつ} とは，
$\textit{WAdm}^{Noeth}$ における $B \to A$ および $C \to B$ に対して
% P10-E0343
$P(C \to B) \wedge P(C \to A) \Rightarrow P(C \to B)$
が成り立つことをいう。全く同じ方法で，$S$ 上の局所 Noether 形式代数空間の間の
射に対する補題
% P10-E0344
\ref{formal-spaces-lemma-composition-property-morphisms} の変種を得る。
\end{remark}
\section{タウト環写像と代数空間による表現可能性}
\label{formal-spaces-section-taut}

\noindent
この節では，局所可算添字付き形式代数空間の間の射が，
エタール局所的に，開イデアルの可算な基本系をもつ弱許容位相環の間の
タウト連続環準同型に対応することを簡単に示す。実際，これは
補題 \ref{formal-spaces-lemma-representable-affine} からかなり明らかなので，
読者にはこの節を飛ばすことを勧める。

\begin{lemma}
\label{formal-spaces-lemma-representable-property-rings}
$B \to A$ を $\textit{WAdm}^{count}$ の矢印とする。
次は同値である。
\begin{enumerate}
\item[(a)] $B \to A$ はタウトである（定義 \ref{formal-spaces-definition-taut}）。
\item[(b)] $B \supset J_1 \supset J_2 \supset J_3 \supset \ldots$
が弱定義イデアルの基本系であるとき，
% P10-E0346
次の可換図式が存在する：
$$
\xymatrix{
A \ar[r] & \ldots \ar[r] & A_3 \ar[r] & A_2 \ar[r] & A_1 \\
B \ar[r] \ar[u] & \ldots \ar[r] & B/J_3 \ar[r] \ar[u] &
B/J_2 \ar[r] \ar[u] & B/J_1 \ar[u]
}
$$
ただし $A_{n + 1}/J_nA_{n + 1} = A_n$ かつ $A = \lim A_n$
を満たし，
% P10-E0347
位相環として考える。
\end{enumerate}
さらに，これらの同値な条件は局所的性質を定める。すなわち，公理
(\ref{formal-spaces-item-axiom-1})，(\ref{formal-spaces-item-axiom-2})，(\ref{formal-spaces-item-axiom-3}) を満たす。
\end{lemma}

\begin{proof}
(a) と (b) の同値性は明らかである。以下で公理の代数的証明を与えるが，
実はそれらは既に証明されている。すなわち補題
\ref{formal-spaces-lemma-representable-affine} を用いると，条件 (a) と (b) は，
対応するアフィン形式代数空間の射が代数空間により表現可能であることを
意味する。
この条件は補題
\ref{formal-spaces-lemma-representable-by-algebraic-spaces-local} によれば
「始域と終域についてエタール局所的」であるから，直ちに公理
(\ref{formal-spaces-item-axiom-1})，(\ref{formal-spaces-item-axiom-2})，(\ref{formal-spaces-item-axiom-3}) を得る。

\medskip\noindent
(\ref{formal-spaces-item-axiom-1})，(\ref{formal-spaces-item-axiom-2})，(\ref{formal-spaces-item-axiom-3}) の
直接的な代数的証明を与える。図式
\ref{formal-spaces-equation-localize} が状況
\ref{formal-spaces-situation-local-property} において与えられているとする。
例 \ref{formal-spaces-example-representable-morphism-from-completion} により，
写像 $A \to (A')^\wedge$ および $B \to (B')^\wedge$ は
(a) と (b) を満たす。

\medskip\noindent
$\varphi$ について (a) と (b) が成り立つと仮定する。
$J \subset B$ を弱定義イデアルとする。このとき $JA$ の閉包，
それぞれ $J(B')^\wedge$ の閉包は，弱定義イデアル
$I \subset A$，それぞれ $J' \subset (B')^\wedge$ である。
さらに $I(A')^\wedge$ の閉包は弱定義イデアル
$I' \subset (A')^\wedge$ である。位相的議論により，$I'$ は
$J(A')^\wedge$ の閉包および $J'(A')^\wedge$ の閉包でもある。
最後に，$J$ が $B$ の弱定義イデアルの基本系を走るとき，
$A$ および $(A')^\wedge$ におけるイデアル $I$ および $I'$ も
基本系をなす。したがって (a) は $\varphi'$ に対して成り立つ。
これは (\ref{formal-spaces-item-axiom-1}) を証明する。

\medskip\noindent
$A \to A'$ が忠実平坦であり，$\varphi'$ について (a) と (b) が
成り立つと仮定する。$J \subset B$ を弱定義イデアルとする。
写像 $B \to (B')^\wedge \to (A')^\wedge$ に対する (a) と (b) を用いると，
$J(A')^\wedge$ の閉包 $I'$ は弱定義イデアルであることが分かる。
特に $I'$ は開であり，従って $A$ における $I'$ の逆像も開である。
そこで（以下で説明するように）
\begin{align*}
A \cap I'
& =
A \cap \bigcap (J(A')^\wedge + \Ker((A')^\wedge \to A'/I_0A')) \\
& =
A \cap \bigcap \Ker((A')^\wedge \to A'/JA' + I_0 A') \\
& = \bigcap (JA + I_0)
\end{align*}
となる。これは補題 \ref{formal-spaces-lemma-closed} により $JA$ の閉包である。
交差は弱定義イデアル $I_0 \subset A$ 全体にわたる。
% P10-E0348
第1の等号は，$(A')^\wedge$ における $0$ の近傍の基本系が，
写像 $(A')^\wedge \to A'/I_0A'$ の核であることによる。
第2の等号は自明である。第3の等号は $A \to A'$ が忠実平坦であること，
すなわち代数，補題
\ref{algebra-lemma-faithfully-flat-universally-injective} による。
したがって $JA$ の閉包は開である。補題
\ref{formal-spaces-lemma-topologically-nilpotent} により $JA$ の閉包は
$A$ の弱定義イデアルである。最後に，弱定義イデアル
$I \subset A$ を一つ取ると，$J$ であって
$J(A')^\wedge$ が $I(A')^\wedge$ の閉包に含まれるものを，
$B \to (B')^\wedge$ に対する性質 (a) と $\varphi'$ により
見いだせる。
したがって (a) は $\varphi$ に対して成り立つ。これは
(\ref{formal-spaces-item-axiom-2}) を証明する。

\medskip\noindent
(\ref{formal-spaces-item-axiom-3}) の証明は省略する。
\end{proof}

\begin{lemma}
\label{formal-spaces-lemma-representable-local-property}
$S$ をスキームとする。$f : X \to Y$ を $S$ 上の局所可算添字付き
形式代数空間の射とする。次は同値である。
\begin{enumerate}
\item 任意の可換図式
$$
\xymatrix{
U \ar[d] \ar[r] & V \ar[d] \\
X \ar[r] & Y
}
$$
で，$U$ と $V$ がアフィン形式代数空間であり，$U \to X$ と $V \to Y$ が
代数空間により表現可能かつエタールであるものに対し，射 $U \to V$ は
$\textit{WAdm}^{count}$ のタウト写像 $B \to A$ に対応する。
\item 補題 \ref{formal-spaces-definition-formal-algebraic-space} におけるような被覆
$\{Y_j \to Y\}$ が存在し，各 $j$ に対して補題
\ref{formal-spaces-definition-formal-algebraic-space} におけるような被覆
$\{X_{ji} \to Y_j \times_Y X\}$ が存在して，各 $X_{ji} \to Y_j$ が
$\textit{WAdm}^{count}$ のタウト環写像に対応する。
\item 補題 \ref{formal-spaces-definition-formal-algebraic-space} におけるような被覆
$\{X_i \to X\}$ が存在し，各 $i$ に対して因数分解
$X_i \to Y_i \to Y$ が存在する。ここで $Y_i$ はアフィン形式代数空間，
$Y_i \to Y$ は代数空間により表現可能かつエタールであり，
$X_i \to Y_i$ は $\textit{WAdm}^{count}$ のタウト環写像に対応する。
\item $f$ は代数空間により表現可能である。
\end{enumerate}
\end{lemma}

\begin{proof}
$\textit{WAdm}^{count}$ の写像が「タウト」であるという性質は，
補題 \ref{formal-spaces-lemma-representable-property-rings} により局所的性質である。
したがって補題 \ref{formal-spaces-lemma-property-defines-property-morphisms} は，
(1)，(2)，(3) が同値であることを正確に述べている。
一方，補題 \ref{formal-spaces-lemma-representable-affine} により，
$\textit{WAdm}^{count}$ の写像が「タウト」であることは，
可算添字付きアフィン形式代数空間の対応する射が
「代数空間により表現可能」であることと正確に対応する。
したがって補題
\ref{formal-spaces-lemma-representable-by-algebraic-spaces-local} の含意
% P10-E0350
(1) $\Rightarrow$ (2) は，(4) がこの補題の (1) を含意することを示す。
同様に，補題 \ref{formal-spaces-lemma-representable-by-algebraic-spaces-local} の含意 (4) $\Rightarrow$ (1) は，
(2) がこの補題の (4) を含意することを示す。
\end{proof}

\section{adic 射}
\label{formal-spaces-section-adic}

\noindent
この節では，局所 adic* 形式代数空間 $X$ および $Y$ の射
「adic 射」$f : X \to Y$ という時折用いられる概念を，一方では
$f$ が代数空間により表現可能であることと，他方では我々の
タウト連続環準同型の概念と結び付ける。まず，有限生成の定義イデアルを
もつ adic 環では，タウト性と adic 性が同値であることを思い出そう。

\begin{lemma}
\label{formal-spaces-lemma-adic-homomorphism}
$A$ と $B$ を pre-adic 位相環とする。$\varphi : A \to B$ を
連続環準同型とする。
\begin{enumerate}
\item $\varphi$ が adic ならば，$\varphi$ はタウトである。
\item $B$ が完備で，$A$ が有限生成の定義イデアルをもち，$\varphi$ が
タウトならば，$\varphi$ は adic である。
\end{enumerate}
特に，条件「$\varphi$ は adic」と「$\varphi$ はタウト」は
$\textit{WAdm}^{adic*}$ の圏上で同値である。
\end{lemma}

\begin{proof}
(1) は補題 \ref{formal-spaces-lemma-adic-taut} である。
(2) は補題 \ref{formal-spaces-lemma-taut-is-adic} である。
最後の主張は (1) と (2) の帰結である。
\end{proof}

\noindent
$S$ をスキームとする。$f : X \to Y$ を $S$ 上の局所 adic* 形式代数空間の
射とする。補題 \ref{formal-spaces-lemma-representable-local-property} により，
次は同値である。
\begin{enumerate}
\item $f$ は代数空間により表現可能である（すなわち，補題
\ref{formal-spaces-lemma-representable-by-algebraic-spaces-local} の同値条件が成り立つ）。
\item 任意の可換図式
$$
\xymatrix{
U \ar[d] \ar[r] & V \ar[d] \\
X \ar[r] & Y
}
$$
で，$U$ と $V$ がアフィン形式代数空間であり，$U \to X$ と $V \to Y$ が
代数空間により表現可能かつエタールであるものに対し，射 $U \to V$ は
\footnote{補題 \ref{formal-spaces-lemma-adic-homomorphism} により，タウトと同値である。}
$\textit{WAdm}^{adic*}$ の adic 写像に対応する。
\end{enumerate}
この状況で，$f$ は {\it adic 射} であるという（形式的な定義は以下で与える）。
この概念および用語は，局所 adic* である形式代数空間の間の射に対してのみ
定義・使用する。そうでなければ，上の (1) と (2) の同値性をもたないからである。

\begin{definition}
\label{formal-spaces-definition-adic-morphism}
$S$ をスキームとする。$f : X \to Y$ を $S$ 上の形式代数空間の射とする。
$X$ と $Y$ は局所 adic* であると仮定する。$f$ が
{\it adic 射} であるとは，$f$ が代数空間により表現可能であることをいう。
上の議論を参照。
\end{definition}

\section{有限型の射}
\label{formal-spaces-section-finite-type}

\noindent
% P10-E0351
Stacks プロジェクトでの設定方法により，以下が本当に正しい構成であり，
より強い概念には別の名前を与えるべきである。

\begin{definition}
\label{formal-spaces-definition-finite-type}
$S$ をスキームとする。$f : Y \to X$ を $S$ 上の形式代数空間の射とする。
\begin{enumerate}
\item $f$ が {\it 局所有限型} であるとは，$f$ が代数空間により表現可能で，
Bootstrap，定義
\ref{bootstrap-definition-property-transformation} の意味で局所有限型であることをいう。
\item $f$ が {\it 有限型} であるとは，$f$ が局所有限型であり，かつ準コンパクト
（定義 \ref{formal-spaces-definition-quasi-compact-morphism}）であることをいう。
\end{enumerate}
\end{definition}

\noindent
ある形式代数空間の有限型射と，Section \ref{formal-spaces-section-tft} で扱う
位相的有限型である連続環写像 $A \to B$ との関係を論じる。

\begin{lemma}
\label{formal-spaces-lemma-characterize-finite-type}
$S$ をスキームとする。$f : X \to Y$ を $S$ 上の形式代数空間の射とする。
次は同値である。
\begin{enumerate}
\item $f$ は有限型である。
\item $f$ は代数空間により表現可能であり，Bootstrap，定義
\ref{bootstrap-definition-property-transformation} の意味で有限型である。
\end{enumerate}
\end{lemma}

\begin{proof}
これは Bootstrap，補題
\ref{bootstrap-lemma-transformations-property-implication}，
代数空間の射についての「準コンパクト $+$ 局所有限型
$\Rightarrow$ 有限型」
という含意，および補題 \ref{formal-spaces-lemma-quasi-compact-representable} から従う。
\end{proof}

\begin{lemma}
\label{formal-spaces-lemma-composition-finite-type}
有限型射の合成は有限型である。同じことが局所有限型射についても成り立つ。
\end{lemma}

\begin{proof}
Bootstrap，補題
\ref{bootstrap-lemma-composition-transformation-property} を参照し，
空間の射，補題 \ref{spaces-morphisms-lemma-composition-finite-type} を用いる。
\end{proof}

\begin{lemma}
\label{formal-spaces-lemma-base-change-finite-type}
有限型射の基底変換は有限型である。同じことが局所有限型射についても成り立つ。
\end{lemma}

\begin{proof}
Bootstrap，補題
\ref{bootstrap-lemma-base-change-transformation-property} を参照し，
空間の射，補題 \ref{spaces-morphisms-lemma-base-change-finite-type} を用いる。
\end{proof}

\begin{lemma}
\label{formal-spaces-lemma-permanence-finite-type}
$S$ をスキームとする。$f : X \to Y$ および $g : Y \to Z$ を
$S$ 上の形式代数空間の射とする。$g \circ f : X \to Z$ が局所有限型ならば，
$f : X \to Y$ も局所有限型である。
\end{lemma}

\begin{proof}
補題 \ref{formal-spaces-lemma-permanence-representable} により，$f$ は代数空間により
表現可能である。$T$ をスキームとし，$T \to Z$ を射とする。このとき空間の射，
補題 \ref{spaces-morphisms-lemma-permanence-finite-type} を
代数空間の射
$T \times_Z X \to T \times_Z Y \to T$ に適用すれば結論を得る。
\end{proof}

\noindent
局所有限型であることは始域および終域上局所的である。

\begin{lemma}
\label{formal-spaces-lemma-finite-type-local}
$S$ をスキームとする。$f : X \to Y$ を $S$ 上の形式代数空間の射とする。
次は同値である：
\begin{enumerate}
\item 射 $f$ は局所有限型である。
\item 次の可換図式が存在する：
$$
\xymatrix{
U \ar[d] \ar[r] & V \ar[d] \\
X \ar[r] & Y
}
$$
ここで $U$ と $V$ は形式代数空間で，鉛直射は代数空間により表現可能かつ
エタールであり，$U \to X$ は全射，$U \to V$ は局所有限型である。
\item 任意の次の可換図式に対して：
$$
\xymatrix{
U \ar[d] \ar[r] & V \ar[d] \\
X \ar[r] & Y
}
$$
ここで $U$ と $V$ は形式代数空間で，鉛直射は代数空間により表現可能かつ
エタールであるとき，射 $U \to V$ は局所有限型である。
\item 補題 \ref{formal-spaces-definition-formal-algebraic-space} におけるような被覆
$\{Y_j \to Y\}$ が存在し，各 $j$ に対して補題
\ref{formal-spaces-definition-formal-algebraic-space} におけるような被覆
$\{X_{ji} \to Y_j \times_Y X\}$ が存在して，各 $j$ と $i$ について
$X_{ji} \to Y_j$ が局所有限型である。
\item 補題 \ref{formal-spaces-definition-formal-algebraic-space} におけるような被覆
$\{X_i \to X\}$ が存在し，各 $i$ に対して因数分解
$X_i \to Y_i \to Y$ が存在する。ここで $Y_i$ はアフィン形式代数空間，
$Y_i \to Y$ は代数空間により表現可能かつエタールであり，
$X_i \to Y_i$ は局所有限型であり，かつ
\item
% P10-E0352
ここにさらに追加。
\end{enumerate}
\end{lemma}

\begin{proof}
% P10-E0353
5つのケースのそれぞれで，射 $f : X \to Y$ は代数空間により表現可能である。
補題 \ref{formal-spaces-lemma-representable-by-algebraic-spaces-local} を参照。
以下ではこれをさらに言及せずに用いる。

\medskip\noindent
(1) が (2) を含意することは明らかである。$U = X$，$V = Y$ と取ればよい。
逆に，(2) の図式が与えられているとする。$T$ をスキームとし，
$T \to Y$ を射とする。このとき
$$
\xymatrix{
U \times_Y T \ar[d] \ar[r] & V \times_Y T \ar[d] \\
X \times_Y T \ar[r] & T
}
$$
を考えることができる。鉛直射はエタールであり，% P10-E0354
上の水平射はそのような射の
基底変換として局所有限型である。したがって空間の射，補題
\ref{spaces-morphisms-lemma-finite-type-local} により，
$X \times_Y T \to T$ は局所有限型である。これは (1) が成り立つということである。

\medskip\noindent
(1) が真であり，(3) の図式を考えるとする。このとき
$U \to Y$ は局所有限型である（合成 $U \to X \to Y$ として，
Bootstrap，補題
\ref{bootstrap-lemma-composition-transformation-property} による）。
$T$ をスキームとし，$T \to V$ を射とする。このとき射影
$T \times_V U \to T$ は
$$
T \times_V U = (T \times_Y U) \times_{(V \times_Y V)} V
\to T \times_Y U \to T
$$
と因数分解する。第2の射は（$U \to X \to Y$ の合成の基底変換として）
局所有限型であり，第1の射は対角射 $V \to V \times_Y V$ の基底変換である。
後者は補題 \ref{formal-spaces-lemma-diagonal-morphism-formal-algebraic-spaces} により
局所有限型である。

\medskip\noindent
(3) が (2) を含意することは明らかである。したがって今や (1)--(3) は同値である。

\medskip\noindent
(4) の条件が意味をもつことに注意する。$Y_j \times_Y X$ は補題
\ref{formal-spaces-lemma-fibre-products} により形式代数空間だからである。(4) が (5) を
含意することは明らかである。

\medskip\noindent
(5) のような $X_i \to Y_i \to Y$ が与えられたとする。このとき
$V = \coprod Y_i$ および $U = \coprod X_i$ と置くと，(5) が (2) を含意する。

\medskip\noindent
最後に (1)--(3) が真であると仮定する。このとき補題
\ref{formal-spaces-definition-formal-algebraic-space} における任意の被覆
$\{Y_j \to Y\}$ と，各 $j$ に対する補題
\ref{formal-spaces-definition-formal-algebraic-space} における任意の被覆
$\{X_{ji} \to Y_j \times_Y X\}$ を選べる。
% P10-E0355
すると $X_{ij} \to Y_j$ は (3) により局所有限型であり，
(4) が成り立つことが分かる。これで証明が終わる。
\end{proof}

\begin{example}
\label{formal-spaces-example-finite-type-from-finite-type-ring-map}
$S$ をスキームとする。$A$ を $S$ 上の弱許容位相環とする。
$A \to A'$ を有限型環写像とする。このとき
$$
(A')^\wedge = \lim_{I \subset A\ w.i.d.} A'/IA'
$$
は弱許容環であり，対応する射
$\text{Spf}((A')^\wedge) \to \text{Spf}(A)$ は表現可能である。
例 \ref{formal-spaces-example-representable-morphism-from-completion} を参照せよ。
$T \to \text{Spf}(A)$ が $T$ を準コンパクトなスキームとする射ならば，
この射は，ある弱定義イデアル $I \subset A$ に対する
$\Spec(A/I)$ を通る（補題 \ref{formal-spaces-lemma-factor-through-thickening}）。
すると
$T \times_{\text{Spf}(A)} \text{Spf}((A')^\wedge)$
は $T \times_{\Spec(A/I)} \Spec(A'/IA')$ に等しく，
$\text{Spf}((A')^\wedge) \to \text{Spf}(A)$ が有限型であることが分かる。
\end{example}

\begin{lemma}
\label{formal-spaces-lemma-locally-finite-type-locally-noetherian}
$S$ をスキームとする。$f : X \to Y$ を
$S$ 上の形式代数空間の射とする。$Y$ が局所 Noether であり，
% P10-E0356
$f$ が局所有限型ならば，$X$ は局所 Noether である。
\end{lemma}

\begin{proof}
補題 \ref{formal-spaces-lemma-finite-type-local} におけるような
$\{Y_j \to Y\}$ および
$\{X_{ij} \to Y_j \times_Y X\}$ を選ぶ。
すると補題 \ref{formal-spaces-lemma-property-goes-up-affine-morphism} により，
各 $X_{ij}$ は Noether である。これで補題が証明された。
\end{proof}

\begin{lemma}
\label{formal-spaces-lemma-fibre-product-Noetherian}
$S$ をスキームとする。$f : X \to Y$ および $Z \to Y$ を
$S$ 上の形式代数空間の射とする。$Z$ が局所 Noether であり，
$f$ が局所有限型ならば，$Z \times_Y X$ は局所 Noether である。
\end{lemma}

\begin{proof}
射 $Z \times_Y X \to Z$ は補題
\ref{formal-spaces-lemma-base-change-finite-type} により局所有限型である。
したがって補題 \ref{formal-spaces-lemma-locally-finite-type-locally-noetherian} から従う。
\end{proof}

\section{全射な射}
\label{formal-spaces-section-surjective}

\noindent
補題 \ref{formal-spaces-lemma-reduction-surjective} により，次の定義は，代数空間の射または
代数空間により表現可能な形式代数空間の射について既にある定義と衝突しない。

\begin{definition}
\label{formal-spaces-definition-surjective}
$S$ をスキームとする。$S$ 上の形式代数空間の射 $f : X \to Y$ が
基礎となる被約代数空間上に全射 $X_{red} \to Y_{red}$ を誘導するとき，
この射は {\it 全射} であるという。
\end{definition}

\begin{lemma}
\label{formal-spaces-lemma-composition-surjective}
2つの全射な射の合成は全射な射である。
\end{lemma}

\begin{proof}
省略する。
\end{proof}

\begin{lemma}
\label{formal-spaces-lemma-base-change-surjective}
全射な射の基底変換は全射な射である。
\end{lemma}

\begin{proof}
省略する。
\end{proof}

\begin{lemma}
\label{formal-spaces-lemma-characterize-surjective}
$S$ をスキームとする。$f : X \to Y$ を $S$ 上の形式代数空間の射とする。
次は同値である：
\begin{enumerate}
\item $f$ は全射である。
\item 任意のスキーム $T$ と射 $T \to Y$ に対し，射影
$X \times_Y T \to T$ は形式代数空間の全射な射である。
\item 任意のアフィンスキーム $T$ と射 $T \to Y$ に対し，射影
$X \times_Y T \to T$ は形式代数空間の全射な射である。
\item 定義 \ref{formal-spaces-definition-formal-algebraic-space} におけるような被覆
$\{Y_j \to Y\}$ が存在し，各 $X \times_Y Y_j \to Y_j$ は
形式代数空間の全射な射である。
\item 形式代数空間の全射な射 $Z \to Y$ が存在し，
$X \times_Y Z \to Z$ は全射であり，かつ
\item
% P10-E0357
ここにさらに追加。
\end{enumerate}
\end{lemma}

\begin{proof}
省略する。
\end{proof}

\section{モノ射}
\label{formal-spaces-section-monomorphisms}

\noindent
定義は次の通りである。

\begin{definition}
\label{formal-spaces-definition-monomorphism}
$S$ をスキームとする。
$S$ 上の形式代数空間の射が層の単射であるとき，これを
{\it モノ射} と呼ぶ。
\end{definition}

\noindent
次はその一例である。$X$ を代数空間とし，$T \subset |X|$ を閉部分集合とする。
このとき，$T$ に沿う $X$ の形式完備化から $X$ への射
$X_{/T} \to X$ はモノ射である。特に，形式代数空間のモノ射は一般には
表現可能でない。

\begin{lemma}
\label{formal-spaces-lemma-composition-monomorphism}
2つのモノ射の合成はモノ射である。
\end{lemma}

\begin{proof}
省略する。
\end{proof}

\begin{lemma}
\label{formal-spaces-lemma-base-change-monomorphism}
モノ射の基底変換はモノ射である。
\end{lemma}

\begin{proof}
省略する。
\end{proof}

\begin{lemma}
\label{formal-spaces-lemma-characterize-monomorphisms}
$S$ をスキームとする。$f : X \to Y$ を $S$ 上の形式代数空間の射とする。
次は同値である：
\begin{enumerate}
\item $f$ はモノ射である。
\item 任意のスキーム $T$ と射 $T \to Y$ に対し，射影
$X \times_Y T \to T$ は形式代数空間のモノ射である。
\item 任意のアフィンスキーム $T$ と射 $T \to Y$ に対し，射影
$X \times_Y T \to T$ は形式代数空間のモノ射である。
\item 定義 \ref{formal-spaces-definition-formal-algebraic-space} におけるような被覆
$\{Y_j \to Y\}$ が存在し，各 $X \times_Y Y_j \to Y_j$ は
形式代数空間のモノ射であり，かつ
\item 射の族 $\{Y_j \to Y\}$ が存在して，
$\coprod Y_j \to Y$ は $(\Sch/S)_{fppf}$ 上の層の全射であり，
すべての $j$ について各 $X \times_Y Y_j \to Y_j$ はモノ射である。
\item 代数空間により表現可能，全射，平坦，かつ局所有限表示な
形式代数空間の射 $Z \to Y$ が存在し，
% P10-E0359
$X \times_Y Z \to X$ はモノ射であり，かつ
\item
% P10-E0358
ここにさらに追加。
\end{enumerate}
\end{lemma}

\begin{proof}
省略する。
\end{proof}

\section{閉埋め込み}
\label{formal-spaces-section-closed-immersions}

\noindent
定義は次の通りである。

\begin{definition}
\label{formal-spaces-definition-closed-immersion}
$S$ をスキームとする。$f : Y \to X$ を $S$ 上の形式代数空間の射とする。
$f$ が代数空間により表現可能であり，Bootstrap，定義
\ref{bootstrap-definition-property-transformation} の意味で閉埋め込みであるとき，
$f$ は {\it 閉埋め込み} であるという。
\end{definition}

\noindent
% P10-E0360
この節を読む際には，冒頭の必須補題を飛ばしてよい。

\begin{lemma}
\label{formal-spaces-lemma-composition-closed-immersion}
2つの閉埋め込みの合成は閉埋め込みである。
\end{lemma}

\begin{proof}
省略する。
\end{proof}

\begin{lemma}
\label{formal-spaces-lemma-base-change-closed-immersion}
閉埋め込みの基底変換は閉埋め込みである。
\end{lemma}

\begin{proof}
省略する。
\end{proof}

\begin{lemma}
\label{formal-spaces-lemma-characterize-closed-immersion}
$S$ をスキームとする。$f : X \to Y$ を $S$ 上の形式代数空間の射とする。
次は同値である：
\begin{enumerate}
\item $f$ は閉埋め込みである。
\item 任意のスキーム $T$ と射 $T \to Y$ に対し，射影
$X \times_Y T \to T$ は閉埋め込みである。
\item 任意のアフィンスキーム $T$ と射 $T \to Y$ に対し，射影
$X \times_Y T \to T$ は閉埋め込みである。
\item 定義 \ref{formal-spaces-definition-formal-algebraic-space} におけるような被覆
$\{Y_j \to Y\}$ が存在し，各 $X \times_Y Y_j \to Y_j$ は
閉埋め込みであり，かつ
\item 代数空間により表現可能，全射，平坦，かつ局所有限表示な
形式代数空間の射 $Z \to Y$ が存在し，
% P10-E0361
$X \times_Y Z \to X$ は閉埋め込みであり，かつ
\item
% P10-E0362
ここにさらに追加。
\end{enumerate}
\end{lemma}

\begin{proof}
省略する。
\end{proof}

\begin{lemma}
\label{formal-spaces-lemma-closed-immersion-into-McQuillan}
$S$ をスキームとする。$X$ を $S$ 上の McQuillan アフィン形式代数空間とする。
$f : Y \to X$ を $S$ 上の形式代数空間の閉埋め込みとする。
このとき $Y$ は McQuillan アフィン形式代数空間であり，$f$ は，
タウトで，閉な核と稠密な像をもつ弱許容位相 $S$-代数の連続準同型
$A \to B$ に対応する。
\end{lemma}

\begin{proof}
$A$ を弱許容位相環として $X = \text{Spf}(A)$ と書く。
$I_\lambda$ を $A$ の弱許容な定義イデアルの基本系とする。このとき
$Y \times_X \Spec(A/I_\lambda)$ は $\Spec(A/I_\lambda)$ の閉部分スキームであり，
従ってアフィンである（定義 \ref{formal-spaces-definition-closed-immersion}）。
$Y \times_X \Spec(A/I_\lambda) = \Spec(B_\lambda)$ と書く。
環写像 $A/I_\lambda \to B_\lambda$ は全射である。従って合成
$A \to B \to B_\lambda$ が全射なので，射影
$$
B = \lim B_\lambda \longrightarrow B_\lambda
$$
は全射である。補題 \ref{formal-spaces-lemma-mcquillan-affine-formal-algebraic-space} により，
$Y$ が McQuillan であることが従う。環写像 $A \to B$ は補題
\ref{formal-spaces-lemma-representable-affine} によりタウトである。
$B$ は完備で $A \to B$ は連続なので，核は閉である。最後に，すべての
$\lambda$ について $A \to B_\lambda$ が全射なので，$B$ における
$A$ の像は稠密である。
\end{proof}

\noindent
上の結果があるにもかかわらず，終域が McQuillan アフィン形式代数空間であるときに
閉埋め込みがどのように振る舞うかは一般には分からない。
注記 \ref{formal-spaces-remark-questions} を参照せよ。

\begin{example}
\label{formal-spaces-example-closed-immersion-from-quotient}
$S$ をスキームとする。$A$ を $S$ 上の弱許容位相環とする。
$K \subset A$ を閉イデアルとする。ここで
$$
B = (A/K)^\wedge = \lim_{I \subset A\ w.i.d.} A/(I + K)
$$
と置くと，射 $\text{Spf}(B) \to \text{Spf}(A)$ は表現可能である。
例 \ref{formal-spaces-example-representable-morphism-from-completion} を参照せよ。
$T \to \text{Spf}(A)$ が $T$ を準コンパクトなスキームとする射ならば，
この射は，ある弱定義イデアル $I \subset A$ に対する
$\Spec(A/I)$ を通る（補題 \ref{formal-spaces-lemma-factor-through-thickening}）。
すると $T \times_{\text{Spf}(A)} \text{Spf}(B)$ は
$T \times_{\Spec(A/I)} \Spec(A/(K + I))$ に等しく，
$\text{Spf}(B) \to \text{Spf}(A)$ が閉埋め込みであることが分かる。
$K$ は閉なので，$A \to B$ の核は $K$ である。ただし，一般には環写像
$A \to B = (A/K)^\wedge$ が全射である必要はないことに注意せよ。
\end{example}

\begin{lemma}
\label{formal-spaces-lemma-monomorphism-iso-over-red}
$S$ をスキームとする。$f : X \to Y$ を形式代数空間の射とする。
次を仮定する：
\begin{enumerate}
\item $f$ は代数空間により表現可能である。
\item $f$ はモノ射である。
\item 包含 $Y_{red} \to Y$ は $f$ を通って因数分解する。
\item $f$ は局所有限型であるか，または $Y$ は局所 Noether である。
\end{enumerate}
このとき $f$ は閉埋め込みである。
\end{lemma}

\begin{proof}
仮定 (2) と (3) は
$X_{red} = X \times_Y Y_{red} = Y_{red}$
を含意する。以下ではこのことを断りなく用いる。

\medskip\noindent
$Y' \to Y$ が $S$ 上の形式代数空間のエタール射ならば，基底変換
$f' : X \times_Y Y' \to Y'$ は条件 (1)--(4) を満たす。
従って補題 \ref{formal-spaces-lemma-characterize-closed-immersion} により，
$Y$ がアフィン形式代数空間であると仮定してよい。

\medskip\noindent
定義 \ref{formal-spaces-definition-affine-formal-algebraic-space} におけるように
$Y = \colim_{\lambda \in \Lambda} Y_\lambda$ と書く。
すると $X_\lambda = X \times_Y Y_\lambda$ は，モノ射
$f_\lambda : X_\lambda \to Y_\lambda$ を備えた代数空間であり，この射は同型
$X_{\lambda, red} \to Y_{\lambda, red}$ を誘導する。
従って $X_\lambda$ は『空間の極限』命題
\ref{spaces-limits-proposition-affine} によりアフィンスキームである
（$X_{\lambda, red} \to X_\lambda$ は全射かつ整だからである）。
証明を終えるには $X_\lambda \to Y_\lambda$ が閉埋め込みであることを示せば十分であり，
次の段落でこれを行う。

\medskip\noindent
$X \to Y$ を，$X_{red} = X \times_Y Y_{red} = Y_{red}$ を満たす
アフィンスキームのモノ射とする。一般には，この条件は $X \to Y$ が
閉埋め込みであることを含意しない。『例』節
\ref{examples-section-epimorphism-not-surjective} を参照せよ。
しかし仮定 (4) により，% P10-E0363
前段落において $X_\lambda \to Y_\lambda$ が有限型であるか，または
$Y_\lambda$ が Noether であることが分かっている。これは，$X \to Y$ が環写像
$R \to A$ に対応し，$I \subset R$ を冪零根基（すなわち $R$ の最大局所冪零イデアル）
とすると $R/I \to A/IA$ が同型であり，さらに $R \to A$ が有限型であるか，
または $R$ が Noether であることを意味する。前者の場合，
$R \to A$ は『代数』補題
\ref{algebra-lemma-surjective-mod-locally-nilpotent} により全射である。
後者の場合，$I$ は有限生成であり，従って冪零であるから，
$R \to A$ は Nakayama の補題，すなわち『代数』補題
\ref{algebra-lemma-NAK} の (11) により全射である。
\end{proof}

\section{制限冪級数}
\label{formal-spaces-section-restricted-power-series}

\noindent
$A$ を線形位相に関して完備な位相環とする（『代数続論』定義
\ref{more-algebra-definition-topological-ring}）。
$I_\lambda$ を開イデアルの基本系とする。$r \geq 0$ を整数とする。
この状況では，極限位相を備えた環をしばしば
$$
A\{x_1, \ldots, x_r\} =
\lim_\lambda A/I_\lambda[x_1, \ldots, x_r] =
\lim_\lambda (A[x_1, \ldots, x_r]/I_\lambda A[x_1, \ldots, x_r])
$$
と表す。言い換えれば，これはイデアル $I_\lambda$ に関する多項式環の完備化である。
$A\{x_1, \ldots, x_r\}$ の元は，$x_1, \ldots, x_r$ に関する冪級数
$$
f = \sum\nolimits_{E = (e_1, \ldots, e_r)} a_E x_1^{e_1} \ldots x_r^{e_r}
$$
で，係数 $a_E \in A$ が $A$ の位相で 0 に収束するものと考えられる。
言い換えれば，任意の $\lambda$ に対し，有限個を除くすべての $a_E$ は
$I_\lambda$ に属する。このため，$A\{x_1, \ldots, x_r\}$ の元は
{\it 制限冪級数} と呼ばれることがある。この環を
$A\langle x_1, \ldots, x_r\rangle$ と表すこともあるが，ここではこの記法を用いない。

\begin{remark}[制限冪級数の普遍性]
\label{formal-spaces-remark-universal-property}
\begin{reference}
\cite[Chapter 0, 7.5.3]{EGA}
\end{reference}
$A \to C$ を完備線形位相環の連続写像とする。このとき，任意の
$A$-代数写像 $A[x_1, \ldots x_r] \to C$ は，制限冪級数上の連続写像
$A\{x_1, \ldots, x_r\} \to C$ へ一意的に延長される。
\end{remark}

\begin{remark}
\label{formal-spaces-remark-I-adic-completion-and-restricted-power-series}
$A$ を環とし，$I \subset A$ をイデアルとする。$A$ が $I$-進完備ならば，
$A[x_1, \ldots, x_r]$ の $I$-進完備化 $A[x_1, \ldots, x_r]^\wedge$ は，
環としては $A$ 上の制限冪級数環である。しかし，
$A[x_1, \ldots, x_r]^\wedge$ が $I$-進完備であることは明らかでない。
$A\{x_1, \ldots, x_r\}$ 上の位相は極限位相（常に完備）と考える一方，
$A[x_1, \ldots, x_r]^\wedge$ 上の位相はしばしば $I$-進位相
（必ずしも完備ではない）と考える。$I$ が有限生成ならば，位相環として
$A\{x_1, \ldots, x_r\} = A[x_1, \ldots, x_r]^\wedge$
である。『代数』補題 \ref{algebra-lemma-hathat-finitely-generated} を参照せよ。
\end{remark}

\section{位相的有限型代数}
\label{formal-spaces-section-tft}

\noindent
ここでは次の定義を用いる。この定義は一般に合意されているものではない。
多くの著者は追加の条件を課すが，それは多くの場合，最も一般的な場合ではなく
特定の型の環だけに関心があるためである。

\begin{definition}
\label{formal-spaces-definition-topologically-finite-type}
$A \to B$ を位相環の連続写像とする（『代数続論』定義
\ref{more-algebra-definition-topological-ring}）。
像が $B$ に稠密である $A$-代数写像
$A[x_1, \ldots, x_n] \to B$ が存在するとき，$B$ は $A$ 上
{\it 位相的有限型} であるという。
\end{definition}

\noindent
$A$ が完備線形位相環ならば，制限冪級数環
$A\{x_1, \ldots, x_r\}$ は $A$ 上位相的有限型である。
$k$ が体ならば，冪級数環 $k[[x_1, \ldots, x_r]]$ は
$k$ 上位相的有限型である。

\medskip\noindent
弱許容位相環の連続タウト写像について，位相的有限型であることは，
対応するアフィン形式代数空間の間の射が有限型であることとちょうど対応する。

\begin{lemma}
\label{formal-spaces-lemma-topologically-finite-type-finite-type}
$S$ をスキームとする。$\varphi : A \to B$ を $S$ 上の弱許容位相環の
連続写像とする。次は同値である：
\begin{enumerate}
\item $\text{Spf}(\varphi) : Y = \text{Spf}(B) \to \text{Spf}(A) = X$
は有限型である。
\item $\varphi$ はタウトであり，$B$ は $A$ 上位相的有限型である。
\end{enumerate}
\end{lemma}

\begin{proof}
補題 \ref{formal-spaces-lemma-representable-affine} を用いて，$\varphi$ のタウト性を
$\text{Spf}(\varphi)$ の表現可能性に結び付けられる。以下ではこのことを
断りなく用いる。従って
$X = \colim \Spec(A/I)$ および $Y = \colim \Spec(B/J(I))$ である。
ここで $I \subset A$ は $A$ の弱定義イデアルを走り，$J(I)$ は $B$ における
$IB$ の閉包である。

\medskip\noindent
(2) を仮定する。像が稠密である環写像
$A[x_1, \ldots, x_r] \to B$ を選ぶ。すると
$A[x_1, \ldots, x_r] \to B \to B/J(I)$ も稠密な像をもち，従って全射である。
ゆえに $B/J(I)$ は $A/I$ 上有限型である。$T \to X$ を，$T$ が
準コンパクトなスキームである射とする。このとき $T \to X$ は，ある $I$ に対する
$\Spec(A/I)$ を通って因数分解する（補題
\ref{formal-spaces-lemma-factor-through-thickening}）。補題
\ref{formal-spaces-lemma-representable-affine} の証明を参照すれば，
$T \times_X Y = T \times_{\Spec(A/I)} \Spec(B/J(I))$ である。
% P10-E0364
従って $T \times_Y X \to T$ は，有限型射
$\Spec(B/J(I)) \to \Spec(A/I)$ の基底変換として有限型である。
ゆえに (1) が成り立つ。

\medskip\noindent
(1) を仮定する。上のような任意の $I \subset A$ を選ぶ。
$\Spec(A/I) \times_X Y = \Spec(B/J(I))$ なので，
$A/I \to B/J(I)$ は有限型である。$B/J(I)$ の $A/I$ 上の生成元に写る
$b_1, \ldots, b_r \in B$ を選ぶ。
\par\noindent
$x_i$ を $b_i$ に写す環写像
$A[x_1, \ldots, x_r] \to B$ の像が稠密であると主張する。
これを示すため，第2の弱定義イデアル $I' \subset I$ を取る。このとき
$$
B/(J(I') + IB) = B/J(I)
$$
である。実際，$J(I)$ は $IB$ の閉包であり，$J(I')$ は開だからである。
従って『代数』補題
\ref{algebra-lemma-surjective-mod-locally-nilpotent} を適用して，
$(A/I')[x_1, \ldots, x_r] \to B/J(I')$ が全射であることが分かる。
ゆえに (2) が成り立ち，証明が完了する。
\end{proof}

\noindent
$A$ を線形位相に関して完備な位相環とする。$(I_\lambda)$ を開イデアルの
基本系とする。$\mathcal{C}$ を，次を満たす逆系 $(B_\lambda)$ の圏とする：
\begin{enumerate}
\item $B_\lambda$ は有限型 $A/I_\lambda$-代数である。
\item $B_\mu \to B_\lambda$ は，
$B_\mu/I_\lambda B_\mu \to B_\lambda$ という同型を誘導する
$A/I_\mu$-代数準同型である。
\end{enumerate}
$\mathcal{C}$ の射は，準同型の両立する系により与えられる。

\begin{lemma}
\label{formal-spaces-lemma-category-affine-over}
$S$ をスキームとする。$X$ を $S$ 上のアフィン形式代数空間とする。
$X$ が McQuillan であると仮定し，$A$ を $X$ に付随する弱許容位相環とする。
このとき，次の圏の間に反同値がある：
\begin{enumerate}
\item 上で導入した圏 $\mathcal{C}$。
\item アフィン形式代数空間の有限型射 $Y \to X$ の圏。
\end{enumerate}
\end{lemma}

\begin{proof}
$(I_\lambda)$ を $A$ の弱許容な定義イデアルの基本系とする。
(2) における $Y$ を考える。このとき
$Y \times_X \Spec(A/I_\lambda)$ はアフィンである
（定義 \ref{formal-spaces-definition-finite-type} および補題
\ref{formal-spaces-lemma-affine-representable-by-algebraic-spaces}）。
$Y \times_X \Spec(A/I_\lambda) = \Spec(B_\lambda)$ と書く。
$\Spec(B_\lambda) \to \Spec(A/I_\lambda)$ が有限型なので
（定義 \ref{formal-spaces-definition-finite-type} による），環写像
$A/I_\lambda \to B_\lambda$ は有限型である。従って
$(B_\lambda)$ は $\mathcal{C}$ の対象である。

\medskip\noindent
逆に，$\mathcal{C}$ の対象 $(B_\lambda)$ が与えられたとき，
$Y = \colim \Spec(B_\lambda)$ と置ける。これはアフィン形式代数空間である。
次を主張する：
$$
Y \times_X \Spec(A/I_\lambda) =
\left(\colim_\mu \Spec(B_\mu)\right) \times_X \Spec(A/I_\lambda) =
\Spec(B_\lambda)
$$
これを示すには，準コンパクトなスキーム $U$ 上で評価したときに同じ値を得ることを
示せば十分である。射
\par\noindent
$U \to \left(\colim_\mu \Spec(B_\mu)\right) \times_X \Spec(A/I_\lambda)$
\par\noindent
は，ある $\mu \geq \lambda$ に対する射
\par\noindent
$U \to \Spec(B_\mu) \times_{\Spec(A/I_\mu)} \Spec(A/I_\lambda)$
\par\noindent
から来る（補題 \ref{formal-spaces-lemma-factor-through-thickening} を2回用いる）。
$\mathcal{C}$ の対象に関する第2の仮定により
\par\noindent
$\Spec(B_\mu) \times_{\Spec(A/I_\mu)} \Spec(A/I_\lambda) = \Spec(B_\lambda)$
\par\noindent
なので，主張が従う。これを用いて，射 $Y \to X$ が有限型であることを示せる。
実際，$U$ を準コンパクトなスキームとする任意の射 $U \to X$ に対して，
ある $\lambda$ に関する因数分解 $U \to \Spec(A/I_\lambda) \to X$ を得る
（上で引用した補題を参照）。従って
% P10-E0365
$$
Y \times_X U =
Y \times_X \Spec(A/I_\lambda)) \times_{\Spec(A/I_\lambda)} U =
\Spec(B_\lambda) \times_{\Spec(A/I_\lambda)} U
$$
は望み通り $U$ 上有限型なスキームである。従って構成
\par\noindent
$(B_\lambda) \mapsto \colim \Spec(B_\lambda)$ は圏 (1) から圏 (2) への関手を与える。

\medskip\noindent
証明を終えるため，上の構成が圏 (1) と圏 (2) の間の互いに準逆な関手を
定めることを示す。一方の向きでは，圏 $\mathcal{C}$ の任意の対象
$(B_\lambda)$ に対して
$$
\left(\colim_\mu \Spec(B_\mu)\right) \times_X \Spec(A/I_\lambda) =
\Spec(B_\lambda)
$$
を示さなければならない。これは上で証明した。他方の向きでは，圏 (2) の
$Y$ に対して
$$
Y = \colim (Y \times_X \Spec(A/I_\lambda))
$$
を示さなければならない。これも準コンパクトな試験対象上で評価し，
$X = \colim \Spec(A/I_\lambda)$ であることを用いれば従う。
\end{proof}

\begin{remark}
\label{formal-spaces-remark-questions}
$A$ を弱許容位相環とし，$(I_\lambda)$ を弱定義イデアルの基本系とする。
$X = \text{Spf}(A)$ と置く。言い換えれば，$X$ は McQuillan アフィン形式代数空間である。
$f : Y \to X$ をアフィン形式代数空間の射とする。一般には $Y$ が McQuillan であるとは限らない。
より具体的には，次の問いを考えられる：
\begin{enumerate}
\item $f : Y \to X$ が閉埋め込みであると仮定する。このとき
$Y$ は McQuillan であり，$f$ は，タウトで，核 $K \subset A$ が閉イデアルであり，
像 $\varphi(A)$ が $B$ に稠密である弱許容位相環の連続写像
$\varphi : A \to B$ に対応する。補題
\ref{formal-spaces-lemma-closed-immersion-into-McQuillan} を参照せよ。
どのような $A$ に関する条件が，例
\ref{formal-spaces-example-closed-immersion-from-quotient} におけるように
$B = (A/K)^\wedge$ であることを保証するか？
\item どのような $A$ に関する条件が，閉埋め込み $f : Y \to X$ が
閉イデアルによる $A$ の商 $A/K$ に対応すること，言い換えれば，対応する連続写像
$\varphi$ が全射かつ開であることを保証するか？
\item $f : Y \to X$ が有限型であると仮定する。このとき補題
\ref{formal-spaces-lemma-category-affine-over} により，$(B_\lambda)$ を $\mathcal{C}$ の対象として
$Y = \colim \Spec(B_\lambda)$ を得る。この場合，すべての $\lambda$ に対し
$B_\lambda$ が $A/I_\lambda$ 上 $r$ 個の元で生成されるような固定した整数 $r$ が
存在する（議論は本質的に補題
\ref{formal-spaces-lemma-topologically-finite-type-finite-type} の (1) $\Rightarrow$ (2) の証明で
既に与えられている）。しかし，射影 $\lim B_\lambda \to B_\lambda$ が全射であること，
すなわち $Y$ が McQuillan であることは明らかでない。
$Y$ が McQuillan でない例はあるか？
\item $f : Y \to X$ が有限型で，$Y$ が McQuillan であると仮定する。
すると $f$ は弱許容位相環の連続写像 $\varphi : A \to B$ に対応する。
実際，$\varphi$ はタウトで，$B$ は $A$ 上位相的有限型である。
補題 \ref{formal-spaces-lemma-topologically-finite-type-finite-type} を参照せよ。
言い換えれば，$f$ は
$$
Y \longrightarrow \mathbf{A}^r_X \longrightarrow X
$$
と因数分解する。ここで第1の射は McQuillan アフィン形式代数空間の閉埋め込みである。
しかしこの場合，$Y \to \mathbf{A}^r_X$ について問い (1) と (2) が生じる。
\end{enumerate}
以下では，$X$ が可算添字付き，すなわち $A$ が開イデアルの可算な基本系をもつ場合に
これらの問いに答える。より一般にこれらの問いへの答え，または反例を知っているならば，
% P10-E0367
\href{mailto:stacks.project@gmail.com}{stacks.project@gmail.com} へメールを送ってほしい。
\end{remark}

\begin{lemma}
\label{formal-spaces-lemma-closed-immersion-into-countably-indexed}
$S$ をスキームとする。$X$ を $S$ 上の可算添字付きアフィン形式代数空間とする。
$f : Y \to X$ を $S$ 上の形式代数空間の閉埋め込みとする。
このとき $Y$ は可算添字付きアフィン形式代数空間であり，$f$ は
$A \to A/K$ に対応する。ここで $A$ は $\textit{WAdm}^{count}$ の対象
（節 \ref{formal-spaces-section-morphisms-rings}）であり，$K \subset A$ は閉イデアルである。
\end{lemma}

\begin{proof}
補題 \ref{formal-spaces-lemma-countably-indexed} により，$A$ を
$\textit{WAdm}^{count}$ の対象として $X = \text{Spf}(A)$ と書ける。
閉埋め込みは表現可能かつアフィンなので，補題
\ref{formal-spaces-lemma-property-goes-up-affine-morphism} により，$Y$ はアフィン形式代数空間かつ
% P10-E0366
可算添字付きである。従って補題 \ref{formal-spaces-lemma-countably-indexed} を再び適用すると，
$B$ を $\textit{WAdm}^{count}$ の対象として $Y = \text{Spf}(B)$ と書ける。
補題 \ref{formal-spaces-lemma-closed-immersion-into-McQuillan} により，$f$ は
$\textit{WAdm}^{count}$ の射 $A \to B$ で，タウトかつ稠密な像をもつものから
与えられる。証明を終えるため，補題 \ref{formal-spaces-lemma-dense-image-surjective} を適用する。
\end{proof}

\begin{lemma}
\label{formal-spaces-lemma-quotient-restricted-power-series}
$B \to A$ を $\textit{WAdm}^{count}$ の矢印とする。節
\ref{formal-spaces-section-morphisms-rings} を参照せよ。次は同値である：
\begin{enumerate}
\item[(a)] $B \to A$ はタウトであり，任意の弱定義イデアル $J \subset B$ に対し，
$I \subset A$ を $JA$ の閉包とすると $B/J \to A/I$ は有限型である。
\item[(b)] $B \to A$ はタウトであり，$B$ の弱定義イデアルの余終系
$(J_\lambda)$ に対し，$I_\lambda \subset A$ を $J_\lambda A$ の閉包とすると
$B/J_\lambda \to A/I_\lambda$ は有限型である。
\item[(c)] $B \to A$ はタウトであり，$A$ は $B$ 上位相的有限型である。
\item[(d)] $A$ は位相 $B$-代数として，$B\{x_1, \ldots, x_n\}$ の
閉イデアルによる商と同型である。
\end{enumerate}
さらに，これらの同値な条件は局所的性質を定める。すなわち，公理
(\ref{formal-spaces-item-axiom-1})，(\ref{formal-spaces-item-axiom-2})，(\ref{formal-spaces-item-axiom-3}) を満たす。
\end{lemma}

\begin{proof}
(a) $\Rightarrow$ (b)，(c) $\Rightarrow$ (a)，(d) $\Rightarrow$ (c) は
定義から直ちに従う。(b) が成り立つと仮定し，$J \subset B$ および
$I \subset A$ を (a) におけるものとする。次の可換図式を選ぶ：
$$
\xymatrix{
A \ar[r] & \ldots \ar[r] & A_3 \ar[r] & A_2 \ar[r] & A_1 \\
B \ar[r] \ar[u] & \ldots \ar[r] & B/J_3 \ar[r] \ar[u] &
B/J_2 \ar[r] \ar[u] & B/J_1 \ar[u]
}
$$
これは補題 \ref{formal-spaces-lemma-representable-property-rings} におけるように
$A_{n + 1}/J_nA_{n + 1} = A_n$ および $A = \lim A_n$ を満たす。
任意の $m$ に対し，$J_\lambda \subset J_m$ となる $\lambda$ が存在する。
$B/J_\lambda \to A/I_\lambda$ は有限型なので，
$B/J_m \to A/I_m$ も有限型である。$A_1$ の $B/J_1$ 上の生成元
$\alpha_1, \ldots, \alpha_n \in A_1$ を取る。$A$ は全射な遷移写像をもつ系の
可算極限なので，$A_1$ の $\alpha_1, \ldots, \alpha_n$ に写る
$a_1, \ldots, a_n \in A$ を見つけられる。注記
\ref{formal-spaces-remark-universal-property} により，$x_i$ を $a_i$ に写す連続写像
$B\{x_1, \ldots, x_n\} \to A$ を得る。この写像は『代数』補題
\ref{algebra-lemma-surjective-mod-locally-nilpotent} により全射
$(B/J_m)[x_1, \ldots, x_n] \to A_m$ を誘導する。
$m \geq 1$ に対し，短完全列
$$
0 \to K_m \to (B/J_m)[x_1, \ldots, x_n] \to A_m \to 0
$$
を得る。$A_{m + 1}/J_mA_{m + 1} = A_m$ なので，誘導された遷移写像
$K_{m + 1} \to K_m$ は全射である。従って，これらの短完全列の逆極限は完全である。
『代数』補題 \ref{algebra-lemma-ML-exact-sequence} を参照せよ。
$B\{x_1, \ldots, x_n\} = \lim (B/J_m)[x_1, \ldots, x_n]$ かつ
$A = \lim A_m$ なので，$B\{x_1, \ldots, x_n\} \to A$ は全射かつ開である。
$A$ は完備なので，核は閉イデアルである。このようにして
(a)，(b)，(c)，(d) が同値であることが分かる。

\medskip\noindent
図式 (\ref{formal-spaces-equation-localize}) が状況
\ref{formal-spaces-situation-local-property} におけるように与えられたとする。例
\ref{formal-spaces-example-finite-type-from-finite-type-ring-map} により，写像
$A \to (A')^\wedge$ および $B \to (B')^\wedge$ は
(a)，(b)，(c)，(d) を満たす。さらに，補題
\ref{formal-spaces-lemma-representable-property-rings} により，公理
(\ref{formal-spaces-item-axiom-1}) と (\ref{formal-spaces-item-axiom-2}) を証明するためには，
$B \to A$ と $(B')^\wedge \to (A')^\wedge$ がともにタウトであると仮定してよい。
弱定義イデアル $J \subset B$ を選ぶ。$J' \subset (B')^\wedge$，
$I \subset A$，$I' \subset (A')^\wedge$ をそれぞれ
$J(B')^\wedge$，$JA$，$J(A')^\wedge$ の閉包とする。
上で述べたことにより，次の可換図式を考えれば十分である：
$$
\xymatrix{
A/I \ar[r] & (A')^\wedge/I' \\
B/J \ar[r] \ar[u]^{\overline{\varphi}} &
(B')^\wedge/J' \ar[u]_{\overline{\varphi}'}
}
$$
そして (1) $\overline{\varphi}$ が有限型
$\Rightarrow \overline{\varphi}'$ が有限型であること，および (2) $A \to A'$ が忠実平坦ならば，
$\overline{\varphi}'$ が有限型 $\Rightarrow \overline{\varphi}$ が有限型であることを
示せばよい。$(B')^\wedge$ と $(A')^\wedge$ 上の位相の構成により，
$(B')^\wedge/J' = B'/JB'$ かつ $(A')^\wedge/I' = A'/IA'$ であることに注意する。
特に，図式の水平写像はエタールである。ここで (1) は『代数』補題
\ref{algebra-lemma-compose-finite-type} から従い，(2) は，環写像
$A/I \to (A')^\wedge/I' = A'/IA'$ が忠実平坦かつエタールなので，
『降下』補題 \ref{descent-lemma-finite-type-local-source-fppf-algebra} から従う。

\medskip\noindent
公理 (\ref{formal-spaces-item-axiom-3}) の証明は省略する。
\end{proof}

\begin{lemma}
\label{formal-spaces-lemma-quotient-restricted-power-series-admissible}
補題 \ref{formal-spaces-lemma-quotient-restricted-power-series} において，$B$ が許容
（例えば adic）ならば，同値な条件 (a)--(d) はさらに次とも同値である：
\begin{enumerate}
\item[(e)] $B \to A$ はタウトであり，ある定義イデアル $J \subset B$ に対し，
$I \subset A$ を $JA$ の閉包とすると $B/J \to A/I$ は有限型である。
\end{enumerate}
\end{lemma}

\begin{proof}
(e) が (a) を含意することを示せば十分である。$J' \subset B$ を弱定義イデアルとし，
$I' \subset A$ を $J'A$ の閉包とする。$B/J' \to A/I'$ が有限型であることを
示さなければならない。より小さい弱定義イデアル $J'' = J' \cap J$ に対して
対応する主張が成り立てば，$J'$ に対しても成り立つ。従って
$J' \subset J$ と仮定してよい。$J$ は単なる弱定義イデアルではなく
定義イデアルなので，ある $n \geq 1$ に対し $J^n \subset J'$ である。
従って，完全な行をもつ図式
$$
\xymatrix{
0 \ar[r] & I/I' \ar[r] & A/I' \ar[r] & A/I \ar[r] & 0 \\
0 \ar[r] & J/J' \ar[r] \ar[u] & B/J' \ar[r] \ar[u] & B/J \ar[r] \ar[u] & 0
}
$$
を考えられる。$I' \subset A$ は開であり，$I$ は $J A$ の閉包なので，
$I/I' = (J/J') \cdot A/I'$ である。$J/J'$ は冪零イデアルであり，
$B/J \to A/I$ は有限型なので，『代数』補題
\ref{algebra-lemma-finite-type-mod-nilpotent} により，望み通り
$A/I'$ は $B/J'$ 上有限型である。
\end{proof}

\begin{lemma}
\label{formal-spaces-lemma-representable-affine-finite-type}
$S$ をスキームとする。$f : X \to Y$ をアフィン形式代数空間の射とする。
% P10-E0640
$Y$ が可算添字付きであると仮定する。次は同値である：
\begin{enumerate}
\item $f$ は局所有限型である。
\item $f$ は有限型である。
\item $f$ は $\textit{WAdm}^{count}$ の射 $B \to A$ に対応する
（節 \ref{formal-spaces-section-morphisms-rings}）。この射は補題
\ref{formal-spaces-lemma-quotient-restricted-power-series} の同値な条件を満たす。
\end{enumerate}
\end{lemma}

\begin{proof}
$X$ と $Y$ はアフィンなので，条件 (1) と (2) が同値であることは明らかである。
(1) と (2) の場合，射 $f$ は定義により代数空間により表現可能であり，従って補題
\ref{formal-spaces-lemma-affine-representable-by-algebraic-spaces} によりアフィンである。
ゆえに (1) または (2) が成り立つならば，補題
\ref{formal-spaces-lemma-property-goes-up-affine-morphism} により $X$ は可算添字付きである。
補題 \ref{formal-spaces-lemma-countably-indexed} を参照し，$A$ と $B$ を
$\textit{WAdm}^{count}$ に属する位相 $S$-代数として
$X = \text{Spf}(A)$ および $Y = \text{Spf}(B)$ と書く。
補題 \ref{formal-spaces-lemma-morphism-between-formal-spectra} により，$f$ は連続写像
$B \to A$ に対応する。従って結果は補題
\ref{formal-spaces-lemma-topologically-finite-type-finite-type} から従う。
\end{proof}

\begin{lemma}
\label{formal-spaces-lemma-finite-type-local-property}
$S$ をスキームとする。$f : X \to Y$ を $S$ 上の局所可算添字付き
形式代数空間の射とする。次は同値である：
\begin{enumerate}
\item 任意の可換図式
$$
\xymatrix{
U \ar[d] \ar[r] & V \ar[d] \\
X \ar[r] & Y
}
$$
で，$U$ と $V$ がアフィン形式代数空間，$U \to X$ と $V \to Y$ が
代数空間により表現可能かつエタールであるものに対し，射 $U \to V$ は，
タウトかつ位相的有限型である $\textit{WAdm}^{count}$ の射に対応する。
\item 定義 \ref{formal-spaces-definition-formal-algebraic-space} におけるような被覆
$\{Y_j \to Y\}$ が存在し，各 $j$ に対して定義
\ref{formal-spaces-definition-formal-algebraic-space} におけるような被覆
$\{X_{ji} \to Y_j \times_Y X\}$ が存在し，各 $X_{ji} \to Y_j$ は，
タウトかつ位相的有限型である $\textit{WAdm}^{count}$ の射に対応する。
\item 定義 \ref{formal-spaces-definition-formal-algebraic-space} におけるような被覆
$\{X_i \to X\}$ が存在し，各 $i$ に対して因数分解
$X_i \to Y_i \to Y$ が存在する。ここで $Y_i$ はアフィン形式代数空間，
$Y_i \to Y$ は代数空間により表現可能かつエタールであり，
% P10-E0368
$X_i \to Y_i$ は，タウトかつ位相的有限型である
$\textit{WAdm}^{count}$ の射に対応する。
\item $f$ は局所有限型である。
\end{enumerate}
\end{lemma}

\begin{proof}
補題 \ref{formal-spaces-lemma-quotient-restricted-power-series} により，性質
$P(\varphi)=$「$\varphi$ はタウトかつ位相的有限型である」は
$\text{WAdm}^{count}$ 上局所的である。従って補題
\ref{formal-spaces-lemma-property-defines-property-morphisms} により，条件 (1)，(2)，(3) は同値である。
一方，補題 \ref{formal-spaces-lemma-representable-affine-finite-type} により，
$\textit{WAdm}^{count}$ の射に関する条件 $P$ は，可算添字付きアフィン形式代数空間の
射が局所有限型であることとちょうど対応する。従って補題
\ref{formal-spaces-lemma-finite-type-local} の含意 (1) $\Rightarrow$ (3) により，
(4) はこの補題の (1) を含意する。同様に，補題
\ref{formal-spaces-lemma-finite-type-local} の含意 (4) $\Rightarrow$ (1) により，
(2) はこの補題の (4) を含意する。
\end{proof}

\section{射の分離公理}
\label{formal-spaces-section-separation-axioms}

\noindent
この節は，形式代数空間の射に対する『空間の射』節
\ref{spaces-morphisms-section-separation-axioms} の類似である。

\begin{definition}
\label{formal-spaces-definition-separated-morphism}
$S$ をスキームとする。$f : X \to Y$ を $S$ 上の形式代数空間の射とする。
$\Delta_{X/Y} : X \to X \times_Y X$ を対角射とする。
\begin{enumerate}
\item $\Delta_{X/Y}$ が閉埋め込みであるとき，$f$ は {\it 分離} であるという。
\item $\Delta_{X/Y}$ が準コンパクトであるとき，$f$ は {\it 準分離} であるという。
\end{enumerate}
\end{definition}

\noindent
$\Delta_{X/Y}$ は補題
\ref{formal-spaces-lemma-diagonal-morphism-formal-algebraic-spaces} により（スキームにより）
表現可能なので，アフィンスキーム $T$ からの射 $T \to X \times_Y X$ を考え，
$$
E = T \times_{X \times_Y X} X \longrightarrow T
$$
が準コンパクトまたは閉埋め込みであるかを調べれば判定できる。補題
\ref{formal-spaces-lemma-quasi-compact-representable} または定義
\ref{formal-spaces-definition-closed-immersion} を参照せよ。スキーム $E$ は，$Y$ への射として
一致する2つの射 $a, b : T \to X$ の等化子であり，$E \to T$ は
モノ射かつ局所有限型であることに注意する。

\begin{lemma}
\label{formal-spaces-lemma-base-change-separated}
定義 \ref{formal-spaces-definition-separated-morphism} に列挙されたすべての分離公理は
基底変換の下で安定である。
\end{lemma}

\begin{proof}
$f : X \to Y$ および $Y' \to Y$ を形式代数空間の射とする。
$f' : X' \to Y'$ を $Y' \to Y$ による $f$ の基底変換とする。
すると $\Delta_{X'/Y'}$ は，射
$X' \times_{Y'} X' \to X \times_Y X$ による $\Delta_{X/Y}$ の基底変換である。
定義 \ref{formal-spaces-definition-separated-morphism} で用いられる対角射の各性質は
基底変換の下で安定である。従って補題が成り立つ。
\end{proof}

\begin{lemma}
\label{formal-spaces-lemma-fibre-product-after-map}
$S$ をスキームとする。$f : X \to Z$，$g : Y \to Z$，$Z \to T$ を
$S$ 上の形式代数空間の射とする。誘導される射
$i : X \times_Z Y \to X \times_T Y$ を考える。このとき：
\begin{enumerate}
\item $i$ は（スキームにより）表現可能，局所有限型，局所準有限，分離，かつモノ射である。
\item $Z \to T$ が分離ならば，$i$ は閉埋め込みである。
\item $Z \to T$ が準分離ならば，$i$ は準コンパクトである。
\end{enumerate}
\end{lemma}

\begin{proof}
一般の圏論により，次の図式はファイバー積図式である：
$$
\xymatrix{
X \times_Z Y \ar[r]_i \ar[d] & X \times_T Y \ar[d] \\
% P10-E0369
Z \ar[r]^-{\Delta_{Z/T}} \ar[r] & Z \times_T Z
}
$$
従って $i$ は対角射 $\Delta_{Z/T}$ の基底変換である。
ゆえに補題は補題 \ref{formal-spaces-lemma-diagonal-morphism-formal-algebraic-spaces} から従う。
\end{proof}

\begin{lemma}
\label{formal-spaces-lemma-composition-separated}
定義 \ref{formal-spaces-definition-separated-morphism} に列挙されたすべての分離公理は
射の合成の下で安定である。
\end{lemma}

\begin{proof}
$f : X \to Y$ および $g : Y \to Z$ を，問題の公理が適用される形式代数空間の射とする。
対角射 $\Delta_{X/Z}$ は合成
$$
X \longrightarrow X \times_Y X \longrightarrow X \times_Z X.
$$
である。分離公理は，対角射がある性質 $\mathcal{P}$ をもつことを要求して定義される。
上の補題 \ref{formal-spaces-lemma-fibre-product-after-map} により，第2の射もこの性質をもつ。
性質 $\mathcal{P}$ をもつ（表現可能な）射の合成も性質 $\mathcal{P}$ をもつ射なので，
補題が従う。
\end{proof}

\begin{lemma}
\label{formal-spaces-lemma-separated-local}
$S$ をスキームとする。$f : X \to Y$ を $S$ 上の形式代数空間の射とする。
$\mathcal{P}$ を定義 \ref{formal-spaces-definition-separated-morphism} の分離公理のいずれかとする。
次は同値である：
\begin{enumerate}
\item $f$ は $\mathcal{P}$ をもつ。
\item 任意のスキーム $Z$ と射 $Z \to Y$ に対し，$f$ の基底変換
$Z \times_Y X \to Z$ は $\mathcal{P}$ をもつ。
\item 任意のアフィンスキーム $Z$ と任意の射 $Z \to Y$ に対し，$f$ の基底変換
$Z \times_Y X \to Z$ は $\mathcal{P}$ をもつ。
\item 任意のアフィンスキーム $Z$ と任意の射 $Z \to Y$ に対し，
形式代数空間 $Z \times_Y X$ は $\mathcal{P}$ をもつ
（定義 \ref{formal-spaces-definition-separated} を参照）。
\item 定義 \ref{formal-spaces-definition-formal-algebraic-space} におけるような被覆
$\{Y_j \to Y\}$ が存在し，すべての $j$ について基底変換
$Y_j \times_Y X \to Y_j$ は $\mathcal{P}$ をもつ。
\end{enumerate}
\end{lemma}

\begin{proof}
補題 \ref{formal-spaces-lemma-base-change-separated} を，以下では断りなく繰り返し用いる。
特に，(1) が (2) を，(2) が (3) を含意することは明らかである。

\medskip\noindent
(3) を仮定し，$Z$ がアフィンスキームである射 $Z \to Y$ を取る。
$U$，$V$ をアフィンスキームとし，射
$a : U \to Z \times_Y X$ および $b : V \to Z \times_Y X$ を取る。このとき
$$
U \times_{Z \times_Y X} V =
(Z \times_Y X) \times_{\Delta, (Z \times_Y X) \times_Z (Z \times_Y X)}
(U \times_Z V)
$$
である。$\mathcal{P} =$「準分離」ならばこれは準コンパクトであり，
$\mathcal{P} =$「分離」ならば $U \times_Z V$ への閉埋め込みを備えた
アフィンスキームであることが分かる。従って (4) が成り立つ。

\medskip\noindent
(4) を仮定し，$Z$ がアフィンスキームである射 $Z \to Y$ を取る。
$U$，$V$ をアフィンスキームとし，射
$a : U \to Z \times_Y X$ および $b : V \to Z \times_Y X$ を取る。
上の議論を逆にたどると，$\mathcal{P} =$「準分離」ならば
$U \times_{Z \times_Y X} V \to U \times_Z V$ は準コンパクトであり，
$\mathcal{P} =$「分離」ならば閉埋め込みであることが分かる。
$U$ が $Z \times_Y X$ のエタール被覆を走るように上の $U$ と $V$ を選べるので，
対応する射
$$
U \times_Z V \to (Z \times_Y X) \times_Z (Z \times_Y X)
$$
はアフィンによるエタール被覆をなす。従って
$\Delta : (Z \times_Y X) \to (Z \times_Y X) \times_Z (Z \times_Y X)$
は，それぞれ準コンパクトまたは閉埋め込みであると結論できる。
ゆえに (3) が成り立つ。

\medskip\noindent
(3) が (5) を含意することを証明する。(3) を仮定し，
$\{Y_j \to Y\}$ を定義 \ref{formal-spaces-definition-formal-algebraic-space} におけるものとする。
射
$$
\Delta_j :
Y_j \times_Y X
\longrightarrow
(Y_j \times_Y X) \times_{Y_j} (Y_j \times_Y X) =
Y_j \times_Y X \times_Y X
$$
% P10-E0370
が対応する性質（すなわち準コンパクトまたは閉埋め込み）をもつことを示さなければならない。
定義 \ref{formal-spaces-definition-affine-formal-algebraic-space} におけるように
$Y_j = \colim Y_{j, \lambda}$ と書く。上の式で $Y_j$ を
$Y_{j, \lambda}$ に置き換えると，(3) が成り立つという仮定によりこの性質をもつ。
表示された射は射 $\Delta_{j, \lambda}$ の余極限であり，$\Delta_j$ が対応する性質を
もつかどうかは，$Y_j \times_Y X \times_Y X$ へ写るアフィンスキームによる
基底変換後に判定できるので，補題 \ref{formal-spaces-lemma-factor-through-thickening} から結論を得る。

\medskip\noindent
(5) が (1) を含意することを証明する。$\{Y_j \to Y\}$ を (5) におけるものとする。
このときファイバー積図式
$$
\xymatrix{
\coprod Y_j \times_Y X \ar[r] \ar[d] &
X \ar[d] \\
\coprod Y_j \times_Y X \times_Y X \ar[r] &
X \times_Y X
}
$$
を得る。仮定により左の鉛直射は準コンパクトまたは閉埋め込みである。
『空間』補題
\ref{spaces-lemma-descent-representable-transformations-property} により，
右の鉛直射も準コンパクトまたは閉埋め込みである。
\end{proof}

\section{固有射}
\label{formal-spaces-section-proper}

\noindent
ここでは次の定義を用いる。

\begin{definition}
\label{formal-spaces-definition-proper}
$S$ をスキームとする。$f : Y \to X$ を $S$ 上の形式代数空間の射とする。
$f$ が代数空間により表現可能であり，Bootstrap，定義
\ref{bootstrap-definition-property-transformation} の意味で固有であるとき，
$f$ は {\it 固有} であるという。
\end{definition}

\noindent
定義から，固有射は有限型であることが従う。

\begin{lemma}
\label{formal-spaces-lemma-proper-local}
$S$ をスキームとする。$f : X \to Y$ を $S$ 上の形式代数空間の射とする。
次は同値である：
\begin{enumerate}
\item $f$ は固有である。
\item 任意のスキーム $Z$ と射 $Z \to Y$ に対し，$f$ の基底変換
$Z \times_Y X \to Z$ は固有である。
\item 任意のアフィンスキーム $Z$ と任意の射 $Z \to Y$ に対し，$f$ の基底変換
$Z \times_Y X \to Z$ は固有である。
\item 任意のアフィンスキーム $Z$ と任意の射 $Z \to Y$ に対し，
形式代数空間 $Z \times_Y X$ は $Z$ 上固有な代数空間である。
\item 定義 \ref{formal-spaces-definition-formal-algebraic-space} におけるような被覆
$\{Y_j \to Y\}$ が存在し，すべての $j$ について基底変換
$Y_j \times_Y X \to Y_j$ は固有である。
\end{enumerate}
\end{lemma}

\begin{proof}
省略する。
\end{proof}

\begin{lemma}
\label{formal-spaces-lemma-base-change-proper}
形式代数空間の固有射は基底変換により保たれる。
\end{lemma}

\begin{proof}
これは補題 \ref{formal-spaces-lemma-proper-local} と基底変換の推移性の直ちに従う帰結である。
\end{proof}

\section{形式代数空間と fpqc 被覆}
\label{formal-spaces-section-fpqc}

\noindent
この節は『空間の性質』節 \ref{spaces-properties-section-fpqc} の類似である。
まずそちらの節を読まれたい。

\begin{lemma}
\label{formal-spaces-lemma-sheaf-fpqc}
\begin{slogan}
形式代数空間は fpqc 層である
\end{slogan}
$S$ をスキームとする。$X$ を $S$ 上の形式代数空間とする。
このとき $X$ は fpqc 位相に対する層条件を満たす。
\end{lemma}

\begin{proof}
証明は『空間の性質』命題
\ref{spaces-properties-proposition-sheaf-fpqc} の証明と {\bf 同一} である。
$X$ は Zariski 位相に対する層なので，次を示せば十分である。
アフィンの全射平坦射 $f : T' \to T$ が与えられたとき，
$X(T)$ は2つの写像 $X(T') \to X(T' \times_T T')$ の等化子である。
『位相』補題 \ref{topologies-lemma-sheaf-property-fpqc} を参照せよ。

\medskip\noindent
$a, b : T \to X$ を，$a \circ f = b \circ f$ を満たす2つの射とする。
$a = b$ を示さなければならない。ファイバー積
$$
E = X \times_{\Delta_{X/S}, X \times_S X, (a, b)} T.
$$
を考える。補題 \ref{formal-spaces-lemma-diagonal-formal-algebraic-space} により，射
$\Delta_{X/S}$ は表現可能なモノ射である。従って $E \to T$ はスキームのモノ射である。
仮定 $a \circ f = b \circ f$ は，$T' \to T$ が（一意的に）$E$ を通って
因数分解することを含意する。可換図式
$$
\xymatrix{
T' \times_T E \ar[r] \ar[d] & E \ar[d] \\
T' \ar[r] \ar@/^5ex/[u] \ar[ru] & T
}
$$
を考える。射影 $T' \times_T E \to T'$ は切断をもつモノ射なので，同型である。
従って『降下』補題 \ref{descent-lemma-descending-property-isomorphism} により，
$E \to T$ は同型である。これは望み通り $a = b$ であることを意味する。

\medskip\noindent
次に，$c : T' \to X$ を，2つの合成
$T' \times_T T' \to T' \to X$ が同じである射とする。
$T' \to T$ との合成が $c$ である射 $a : T \to X$ を見つけなければならない。
アフィン形式スキーム $U$ とエタール射 $U \to X$ を，$|U| \to |X_{red}|$ の像が
$|c| : |T'| \to |X_{red}|$ の像を含むように選ぶ。これは定義
\ref{formal-spaces-definition-formal-algebraic-space}，『空間の性質』補題
\ref{spaces-properties-lemma-topology-points}，アフィン形式代数空間の有限和が
アフィン形式代数空間であるという事実，および $|T'|$ が準コンパクトであるという
事実により可能である（小さな議論は省略する）。射 $U \to X$ はスキームにより
表現可能であり（補題 \ref{formal-spaces-lemma-presentation-representable}），分離である
（補題 \ref{formal-spaces-lemma-separated-from-separated}）。従って
$$
V = U \times_{X, c} T' \longrightarrow T'
$$
はスキームのエタールかつ分離な射である。また，$U \to X$ の選び方により全射でもある。
% P10-E0641
この点を論じたくなければ，$U$ をアフィン形式代数空間の非交和で置き換えて
$U \to X$ を全射にしてもよく，他の議論はすべて同様に働く。
$c \circ \text{pr}_0 = c \circ \text{pr}_1$ であることは，
$V/T'/T$ 上に降下データを得ることを意味する
（『降下』定義 \ref{descent-definition-descent-datum}）。実際，
\begin{align*}
V \times_{T'} (T' \times_T T')
& =
U \times_{X, c \circ \text{pr}_0} (T' \times_T T') \\
& =
(T' \times_T T') \times_{c \circ \text{pr}_1, X} U \\
& =
(T' \times_T T') \times_{T'} V
\end{align*}
である。射 $V \to T'$ は『射続論』補題
\ref{more-morphisms-lemma-etale-separated-ind-quasi-affine} により ind-準アフィンである
（エタール射は局所準有限である。『射』補題
\ref{morphisms-lemma-etale-locally-quasi-finite} を参照せよ）。
『亜群続論』補題 \ref{more-groupoids-lemma-ind-quasi-affine} により，降下データは有効である。
$V$ 上の与えられた降下データおよび $T' \times_T W$ 上の標準的降下データと両立する同型
$\alpha : T' \times_T W \to V$ をもつ射 $W \to T$ があるとする。
すると $W \to T$ は全射かつエタールである（『降下』補題
\ref{descent-lemma-descending-property-surjective} および
\ref{descent-lemma-descending-property-etale}）。合成
$$
b' : T' \times_T W \longrightarrow V = U \times_{X, c} T' \longrightarrow U
$$
を考える。2つの合成
$b' \circ (\text{pr}_0, 1), 
b' \circ (\text{pr}_1, 1) :
(T' \times_T T') \times_T W \to T' \times_T W \to U$
は，$\alpha$ の選び方と $c$ の対応する性質により一致する
（計算は省略する）。従って『降下』補題
\ref{descent-lemma-fpqc-universal-effective-epimorphisms} により，
$b'$ は射 $b : W \to U$ に降下する。図式
$$
\xymatrix{
T' \times_T W \ar[r] \ar[d] & W \ar[r]_b & U \ar[d] \\
T' \ar[rr]^c &  & X
}
$$
は可換である。これは，$T$ 上エタール局所的に $a$ の存在を証明したこと，すなわち
$a' : W \to X$ をもつことを意味する。しかし第1段落で一意性を証明したので，
このエタール局所解は貼り合わせ条件を満たす。すなわち
$\text{pr}_0^*a' = \text{pr}_1^*a'$ が $X(W \times_T W)$ の元として成り立つ。
$X$ はエタール層なので，$W$ 上で $a'$ に制限される一意な $a \in X(T)$ を得る。
\end{proof}

\section{アフィン形式スキームからの射}
\label{formal-spaces-section-map-out-of}

\noindent
後で有用となるいくつかの結果を証明する。
同様の性格をもつ非常に一般的な結果は，論文
\cite{Bhatt-Algebraize} に見いだせる。

\begin{lemma}
\label{formal-spaces-lemma-map-into-affine}
$S$ をスキームとする。$A$ を弱許容位相
$S$-代数とする。$X$ を $S$ 上のアフィンスキームとする。このとき自然な写像
$$
\Mor_S(\Spec(A), X)
\longrightarrow
\Mor_S(\text{Spf}(A), X)
$$
は全単射である。
\end{lemma}

\begin{proof}
$X$ がアフィンで，$X = \Spec(B)$ と書けるとする。このとき補題
\ref{formal-spaces-lemma-morphism-between-formal-spectra} より，
射 $\text{Spf}(A) \to \Spec(B)$ は，$B$ に離散位相を入れたときの
連続な $S$-代数写像 $B \to A$ に対応する。
これらは単なる $S$-代数写像であり，従って射
$\Spec(A) \to \Spec(B)$ に対応する。
\end{proof}

\begin{lemma}
\label{formal-spaces-lemma-map-into-scheme}
$S$ をスキームとする。$A$ を弱許容位相
$S$-代数とし，ある弱定義イデアル $I \subset A$ に対して
$A/I$ が局所環であると仮定する。$X$ を $S$ 上のスキームとする。
このとき自然な写像
$$
\Mor_S(\Spec(A), X)
\longrightarrow
\Mor_S(\text{Spf}(A), X)
$$
は全単射である。
\end{lemma}

\begin{proof}
$\varphi : \text{Spf}(A) \to X$ を射とする。$\Spec(A/I)$ は局所的なので，
$\varphi$ は $\Spec(A/I)$ をアフィン開集合 $U \subset X$ に写す。
しかしこれは，任意の定義イデアル $J$ に対して $\Spec(A/J)$ が
$U$ に写ることを含意する。従って補題 \ref{formal-spaces-lemma-map-into-affine} を適用すると，
$\varphi$ は射 $\Spec(A) \to X$ から生じることが分かる。
これで写像の全射性が証明された。単射性の証明は省略する。
\end{proof}

\begin{lemma}
\label{formal-spaces-lemma-map-into-algebraic-space}
$S$ をスキームとする。$R$ を完備局所ネーター $S$-代数とする。
$X$ を $S$ 上の代数空間とする。このとき自然な写像
$$
\Mor_S(\Spec(R), X)
\longrightarrow
\Mor_S(\text{Spf}(R), X)
$$
は全単射である。
\end{lemma}

\begin{proof}
$\mathfrak m$ を $R$ の極大イデアルとする。上のような $R$ に対して
$$
\Mor_S(\Spec(R), X) \longrightarrow \lim \Mor_S(\Spec(R/\mathfrak m^n), X)
$$
が全単射であることを示さなければならない。

\medskip\noindent
単射性：$x, x' : \Spec(R) \to X$ を，右辺の同じ元に写る2つの射とする。
ファイバー積
$$
T = \Spec(R) \times_{(x, x'), X \times_S X, \Delta} X
$$
を考える。このとき $T$ はスキームであり，$T \to \Spec(R)$ は局所有限型，
モノ射，分離，かつ局所準有限である。『空間の射』補題
\ref{spaces-morphisms-lemma-properties-diagonal} を参照せよ。
特に $T$ は局所ネーターである。『射』補題
\ref{morphisms-lemma-finite-type-noetherian} を参照せよ。
$t \in T$ を $\Spec(R)$ の閉点に写る一意な点とする。この点が存在するのは，
$x$ と $x'$ が $R/\mathfrak m$ 上で一致するからである。このとき
$R \to \mathcal{O}_{T, t}$ はネーター環の局所環写像であり，すべての $n$ に対して
$R/\mathfrak m^n \to \mathcal{O}_{T, t}/\mathfrak m^n\mathcal{O}_{T, t}$
は同型である（すべての $n$ に対して $x$ と $x'$ が
$\Spec(R/\mathfrak m^n)$ 上で一致するからである）。
$\mathcal{O}_{T, t}$ はその完備化へ単射的に写るので
（『代数』補題 \ref{algebra-lemma-intersect-powers-ideal-module-zero} を参照せよ），
$R = \mathcal{O}_{T, t}$ と結論できる。従って $x$ と $x'$ は $R$ 上で一致する。

\medskip\noindent
全射性：$(x_n)$ を右辺の元とする。スキーム $U$ と全射エタール射
$U \to X$ を選ぶ。$k = R/\mathfrak m$ と置き，剰余体上に誘導される射を
$x_0 : \Spec(k) \to X$ と書く。スキームの射
$U \times_{X, x_0} \Spec(k) \to \Spec(k)$ は全射エタールである。
従って $U \times_{X, x_0} \Spec(k)$ は $k$ の有限分離体拡大のスペクトルの
空でない非交和である。『射』補題
\ref{morphisms-lemma-etale-over-field} を参照せよ。
従って有限分離体拡大 $k'/k$ と $k'$-点
$u_0 : \Spec(k') \to U$ であって，図式
$$
\xymatrix{
\Spec(k') \ar[d] \ar[r]_-{u_0} & U \ar[d] \\
\Spec(k) \ar[r]^-{x_0} & X
}
$$
を可換にするものが見つかる。$R \subset R'$ を，剰余体上で $k'/k$ を誘導する
ネーター完備局所環の有限エタール拡大とする
（『代数』補題 \ref{algebra-lemma-henselian-cat-finite-etale} および
\ref{algebra-lemma-complete-henselian} を参照せよ）。
$x_n$ の $\Spec(R'/\mathfrak m^nR')$ への制限を $x'_n$ と書く。
『空間の射続論』補題
\ref{spaces-more-morphisms-lemma-etale-formally-etale}
により，$(x'_n)$ に写る元
$(u'_n) \in \lim \Mor_S(\Spec(R'/\mathfrak m^nR'), U)$
が見つかる。補題 \ref{formal-spaces-lemma-map-into-scheme} により，族 $(u'_n)$ は一意な射
$u' : \Spec(R') \to U$ から生じる。その合成を
$x' : \Spec(R') \to X$ と書く。
% P10-E0642
$R' \otimes_R R'$ は，この議論を適用できるネーター完備局所環の
スペクトルの有限積であることに注意する。従って図式
$$
\xymatrix{
\Spec(R' \otimes_R R') \ar[r] \ar[d] & \Spec(R') \ar[d]^{x'} \\
\Spec(R') \ar[r]^{x'} & X
}
$$
は可換である。これは，上で示した単射性と，$x'_n$ が
$R/\mathfrak m^n$ 上で定義された $x_n$ の制限であることによる。
$\{\Spec(R') \to \Spec(R)\}$ は fppf 被覆なので，$x'$ は射
$x : \Spec(R) \to X$ に降下すると結論できる。$x_n$ が
$x$ の $\Spec(R/\mathfrak m^n)$ への制限であることの証明は省略する。
\end{proof}

\begin{lemma}
\label{formal-spaces-lemma-adic-into-completion}
$S$ をスキームとする。$X$ を $S$ 上の代数空間とする。
$T \subset |X|$ を，$X \setminus T \to X$ が準コンパクトとなるような
閉部分集合とする。$R$ を完備局所ネーター $S$-代数とする。
このとき adic 射 $p : \text{Spf}(R) \to X_{/T}$ は，
$g^{-1}(T) = \{\mathfrak m_R\}$ を満たす一意な射
$g : \Spec(R) \to X$ に対応する。
\end{lemma}

\begin{proof}
この主張が意味をなすのは，補題 \ref{formal-spaces-lemma-formal-completion-types} により
$X_{/T}$ が adic* だからである（従って射について adic という用語を
用いることが許される。定義 \ref{formal-spaces-definition-adic-morphism} を参照せよ）。
$p$ が与えられたとする。補題 \ref{formal-spaces-lemma-map-into-algebraic-space} により，
合成 $\text{Spf}(R) \to X_{/T} \to X$ に対応する一意な射
$g : \Spec(R) \to X$ を得る。
$T$ 上の被約被誘導閉部分空間構造を $Z \subset X$ とする。
% P10-E0371
包含射 $Z \to X$ は射 $Z \to X_{/T}$ に対応する。
$p$ は adic なので代数空間により表現可能であり，
$$
\text{Spf}(R) \times_{X_{/T}} Z = \text{Spf}(R) \times_X Z
$$
は $\text{Spf}(R)$ への閉埋め込みを備えた代数空間であることが分かる。
（等号は $X_{/T} \to X$ がモノ射であることによる。）
従ってこのファイバー積は，あるイデアル $J \subset R$ に対する
$\Spec(R/J)$ に等しい。このイデアルは $\mathfrak m_R^{n_0}$ を含み，
ここで $n_0 \geq 1$ である。これは，
$\Spec(R) \times_X Z$ が $\Spec(R)$ の閉部分スキームであることを含意する。
$\Spec(R) \times_X Z = \Spec(R/I)$ と書くと，
$n \geq n_0$ に対する $\Spec(R/\mathfrak m_R^n)$ との交わりは
$\Spec(R/J)$ に等しい。代数的には，これはすべての $n \geq n_0$ に対して
$I + \mathfrak m_R^n = J + \mathfrak m_R^n = J$
であることを意味する。Krull の交叉定理により，これは $I = J$ を含意し，
結論を得る。
\end{proof}
\section{形式代数空間の小エタール・サイト}
\label{formal-spaces-section-etale-sites}

\noindent
次の定義の動機は古典的な形式スキームから来る。形式スキーム
$(\mathfrak X, \mathcal{O}_\mathfrak X)$
の台位相空間は，被約 $\mathfrak X_{red}$ の台位相空間である。

\medskip\noindent
次は重要な注意である。$X$ を，被約 $X_{red}$ をもつ代数空間とする
（『空間の性質』定義
\ref{spaces-properties-definition-reduced-induced-space}）。
このとき『空間の射続論』定理
\ref{spaces-more-morphisms-theorem-topological-invariance} および補題
\ref{spaces-more-morphisms-lemma-topological-invariance} により
$$
\begin{aligned}
X_{spaces, \etale} &= X_{red, spaces, \etale},\quad \\
X_\etale &= X_{red, \etale},\quad \\
X_{affine, \etale} &= X_{red, affine, \etale}
\end{aligned}
$$
である。従って，形式代数空間がたまたま代数空間である場合にも，
次の定義は既存の概念と衝突しない。

\begin{definition}
\label{formal-spaces-definition-etale-sites}
$S$ をスキームとする。$X$ を，被約 $X_{red}$ をもつ形式代数空間とする
（補題 \ref{formal-spaces-lemma-reduction-formal-algebraic-space}）。
\begin{enumerate}
\item $X$ の {\it 小エタール・サイト} $X_\etale$ とは，
『空間の性質』定義
\ref{spaces-properties-definition-etale-site} のサイト
$X_{red, \etale}$ である。
\item サイト $X_{spaces, \etale}$ とは，『空間の性質』定義
\ref{spaces-properties-definition-spaces-etale-site} のサイト
$X_{red, spaces, \etale}$ である。
\item サイト $X_{affine, \etale}$ とは，『空間の性質』補題
\ref{spaces-properties-lemma-alternative} のサイト
$X_{red, affine, \etale}$ である。
\end{enumerate}
\end{definition}

\noindent
補題 \ref{formal-spaces-lemma-identify-spaces-etale} において，
% P10-E0372
$X_{spaces, \etale}$ は，代数空間により表現可能かつエタールな
形式代数空間の射を用いて記述できることを見る。
『空間の性質』補題
\ref{spaces-properties-lemma-compare-etale-sites} および
\ref{spaces-properties-lemma-alternative}
により，同一視
\begin{equation}
\label{formal-spaces-equation-etale-topos}
\Sh(X_\etale) = \Sh(X_{spaces, \etale}) = \Sh(X_{affine, \etale})
\end{equation}
を得る。これを $X$ の {\it （小）エタール・トポス} と呼ぶ。

\begin{lemma}
\label{formal-spaces-lemma-functoriality-etale-site}
$S$ をスキームとする。
$f : X \to Y$ を $S$ 上の形式代数空間の射とする。
\begin{enumerate}
\item 連続関手
$Y_{spaces, \etale} \to X_{spaces, \etale}$
があり，これはサイトの射
$$
f_{spaces, \etale} : X_{spaces, \etale} \to Y_{spaces, \etale}.
$$
を誘導する。
\item 規則 $f \mapsto f_{spaces, \etale}$ は合成と両立する。言い換えると
$(f \circ g)_{spaces, \etale}
= f_{spaces, \etale} \circ g_{spaces, \etale}$ である
（『サイト』定義 \ref{sites-definition-composition-morphisms-sites} を参照せよ）。
\item $f_{spaces, \etale}$ に付随するトポスの射は，
(\ref{formal-spaces-equation-etale-topos}) を介してトポスの射
$f_{small} : \Sh(X_\etale) \to \Sh(Y_\etale)$
を誘導し，その構成は合成と両立する。
\end{enumerate}
\end{lemma}

\begin{proof}
ここでの唯一の点は，補題 \ref{formal-spaces-lemma-reduction-formal-algebraic-space} により
$f$ が被約の射 $X_{red} \to Y_{red}$ を誘導することである。
従ってこの補題は，代数空間の射に対する対応する補題
（『空間の性質』補題
\ref{spaces-properties-lemma-functoriality-etale-site}）から直ちに従う。
\end{proof}

\noindent
形式代数空間の射 $X \to Y$ がエタールならば，トポスの射
$\Sh(X_\etale) \to \Sh(Y_\etale)$ は局所化である。
これを次に述べる。

\begin{lemma}
\label{formal-spaces-lemma-etale-morphism-topoi}
$S$ をスキームとし，$f : X \to Y$ を $S$ 上の形式代数空間の射とする。
$f$ が代数空間により表現可能かつエタールであると仮定する。
この場合，余連続関手 $j : X_\etale \to Y_\etale$ が存在する。
トポスの射 $f_{small}$ は $j$ に付随するトポスの射である。
『サイト』補題 \ref{sites-lemma-cocontinuous-morphism-topoi} を参照せよ。
さらに $j$ は連続でもあるので，『サイト』補題
\ref{sites-lemma-when-shriek} が適用できる。
\end{lemma}

\begin{proof}
誘導される射 $X_{red} \to Y_{red}$ がエタールであることを示せば，
代数空間の場合（『空間の性質』補題
\ref{spaces-properties-lemma-etale-morphism-topoi}）から直ちに従う。
$X \times_Y Y_{red}$ は被約代数空間 $Y_{red}$ 上エタールな代数空間であり，
従ってそれ自身被約であることに注意する。
% P10-E0643
これは『空間の性質』節
\ref{spaces-properties-section-types-properties} における被約代数空間の定義による。
従って望むとおり $X_{red} = X \times_Y Y_{red}$ である。
\end{proof}

\begin{lemma}
\label{formal-spaces-lemma-affine-identify-affine-etale}
$S$ をスキームとする。$X$ を $S$ 上のアフィン形式代数空間とする。
このとき $X_{affine, \etale}$ は，対象が次を満たす形式代数空間の射
$\varphi : U \to X$ である圏と同値である。
\begin{enumerate}
\item $U$ はアフィン形式代数空間である。
\item $\varphi$ は代数空間により表現可能かつエタールである。
\end{enumerate}
\end{lemma}

\begin{proof}
補題で導入された圏を $\mathcal{C}$ と書く。$\mathcal{C}$ の
$\varphi : U \to X$ に対して，射 $\varphi$ は（スキームにより）表現可能かつ
アフィンである。補題
\ref{formal-spaces-lemma-affine-representable-by-algebraic-spaces} を参照せよ。
$X_{affine, \etale} = X_{red, affine, \etale}$ であることを思い出そう。
従って $U \times_X X_{red}$ はアフィンスキームなので，関手
$$
\mathcal{C} \longrightarrow X_{affine, \etale},\quad
(U \to X) \longmapsto U \times_X X_{red}
$$
を定義できる。

\medskip\noindent
証明を終えるため，準逆を構成する。すなわち，定義
\ref{formal-spaces-definition-affine-formal-algebraic-space} のように
$X = \colim X_\lambda$ と書く。各 $\lambda$ に対して
$X_{red} \subset X_\lambda$ は厚化である。従って，例えば『代数続論』補題
\ref{more-algebra-lemma-locally-nilpotent-henselian} により，各 $\lambda$ に対して
同値
$$
X_{red, affine, \etale} = X_{\lambda, affine, \etale}
$$
を得る。従って $U_{red} \to X_{red}$ がエタール射で
$U_{red}$ がアフィンならば，$X$ を定義する系の遷移射と両立する
アフィンスキームのエタール射の系 $U_\lambda \to X_\lambda$ を得る。
従って
$$
U = \colim U_\lambda
$$
を $X$ 上のアフィン形式代数空間として取れる。この構成から
$U \times_X X_\lambda = U_\lambda$ を得る。これは $U \to X$ が
表現可能かつエタールであることを示す。
2つの構成が互いに逆であることの検証は省略する。
\end{proof}

\begin{lemma}
\label{formal-spaces-lemma-affine-etale-mcquillan}
$S$ をスキームとする。$X$ を $S$ 上のアフィン形式代数空間とする。
$X$ が McQuillan，すなわちある弱許容位相 $S$-代数 $A$ に対して
$\text{Spf}(A)$ に等しいと仮定する。
このとき $(X_{affine, \etale})^{opp}$ は，次のような圏と同値である。
\begin{enumerate}
\item 対象は，$A$-代数 $B^\wedge = \lim B/JB$ である。ここで
$A \to B$ はエタール環写像であり，$J$ は $A$ の弱定義イデアルを走る。
\item 射は連続な $A$-代数準同型である。
\end{enumerate}
\end{lemma}

\begin{proof}
補題 \ref{formal-spaces-lemma-affine-identify-affine-etale} および
\ref{formal-spaces-lemma-etale} を組み合わせればよい。
\end{proof}

\begin{lemma}
\label{formal-spaces-lemma-identify-spaces-etale}
$S$ をスキームとする。$X$ を $S$ 上の形式代数空間とする。
このとき $X_{spaces, \etale}$ は，対象が形式代数空間の射
$\varphi : U \to X$ で，$\varphi$ が代数空間により表現可能かつ
エタールであるものからなる圏と同値である。
\end{lemma}

\begin{proof}
補題で導入された圏を $\mathcal{C}$ と書く。
$X_{spaces, \etale} = X_{red, spaces, \etale}$ であることを思い出そう。
従って $U \times_X X_{red}$ は $X_{red}$ 上エタールな代数空間なので，関手
$$
\mathcal{C} \longrightarrow X_{spaces, \etale},\quad
(U \to X) \longmapsto U \times_X X_{red}
$$
を定義できる。

\medskip\noindent
証明を終えるため，準逆を構成する。
$X_{red, spaces, \etale}$ の対象
$\psi : V \to X_{red}$ を選ぶ。次で与えられる関手
$U_{V, \psi} : (\Sch/S)_{fppf} \to \textit{Sets}$
を考える。
$$
U_{V, \psi}(T) = \{(a, b) \mid
a : T \to X,
\ b : T \times_{a, X} X_{red} \to V,
\ \psi \circ b = a|_{T \times_{a, X} X_{red}}\}
$$
変換 $U_{V, \psi} \to X$，$(a, b) \mapsto a$ は圏
$\mathcal{C}$ の対象を定めると主張する。
まず $U_{V, \psi}$ が形式代数空間であることを証明しよう。
$U_{V, \psi}$ は fppf 位相に対する層であることに注意する
（いくつかの詳細は省略する）。次に，定義
\ref{formal-spaces-definition-formal-algebraic-space} のように，
$X_i \to X$ がアフィン形式代数空間によるエタール被覆であると仮定する。
$V_i = V \times_{X_{red}} X_{i, red}$ と置き，射影を
$\psi_i : V_i \to X_{i, red}$ と書く。このとき
$$
U_{V, \psi} \times_X X_i = U_{V_i, \psi_i}
$$
である。これは $X_{i, red} = X_i \times_X X_{red}$ なので
（$X_i \to X$ は代数空間により表現可能かつエタールである），
形式的議論から従う。従って $U_{V_i, \psi_i}$ がアフィン形式代数空間であることを
示せば十分である。実際，そのとき定義
\ref{formal-spaces-definition-formal-algebraic-space} のような被覆
$U_{V_i, \psi_i} \to U_{V, \psi}$ を得る。
一方，補題 \ref{formal-spaces-lemma-etale-morphism-topoi} の証明で，
$\psi_i : V_i \to X_i$ はアフィン形式代数空間の表現可能かつエタールな射
$U_i \to X_i$ の基底変換であることを見た。
このとき，望むとおり $U_i = U_{V_i, \psi_i}$ であることは難しくない。

\medskip\noindent
$U_{V, \psi} \to X$ が代数空間により表現可能かつエタールであることの
検証は省略する。従って逆向きの関手
$(V, \psi) \mapsto (U_{V, \psi} \to X)$ を得る。
2つの構成が互いに逆であることの検証は省略する。
\end{proof}

\begin{lemma}
\label{formal-spaces-lemma-identify-affine-etale}
$S$ をスキームとする。$X$ を $S$ 上の形式代数空間とする。
このとき $X_{affine, \etale}$ は，対象が次を満たす形式代数空間の射
$\varphi : U \to X$ である圏と同値である。
\begin{enumerate}
\item $U$ はアフィン形式代数空間である。
\item $\varphi$ は代数空間により表現可能かつエタールである。
\end{enumerate}
\end{lemma}

\begin{proof}
これは補題 \ref{formal-spaces-lemma-identify-spaces-etale} および
\ref{formal-spaces-lemma-characterize-affine} を組み合わせることで従う。
\end{proof}

\section{構造層}
\label{formal-spaces-section-structure-sheaf}


\noindent
節 \ref{formal-spaces-section-etale-sites} と同様に，本節の議論の動機は古典的な形式スキームから来る。
これらは対 $(\mathfrak X, \mathcal{O}_\mathfrak X)$ として定義され，ここで
$\mathcal{O}_\mathfrak X$ は位相環の層である。

\medskip\noindent
$X$ を形式代数空間とする。$X$ の {\it 構造層} とは，
エタール・サイト $X_\etale$（節 \ref{formal-spaces-section-etale-sites} で定義した）上の
位相環の層 $\mathcal{O}_X$ であって，次の条件が成り立つとき，位相環として
$$
\mathcal{O}_X(U_{red}) = \lim \Gamma(U_\lambda, \mathcal{O}_{U_\lambda})
$$
を満たすものをいう。
\begin{enumerate}
\item $\varphi : U \to X$ は形式代数空間の射である。
\item $U$ はアフィン形式代数空間である。
\item $\varphi$ は代数空間により表現可能かつエタールである。
\item $U_{red} \to X_{red}$ は，対応する $X_\etale$ のアフィン対象である。
補題 \ref{formal-spaces-lemma-identify-affine-etale} を参照せよ。
\item $U = \colim U_\lambda$ は，定義
\ref{formal-spaces-definition-affine-formal-algebraic-space} のような $U$ の余極限表示である。
\end{enumerate}
% P10-E0644
構造層は存在するが，予期しない仕方で振る舞うことがある。

\begin{lemma}
\label{formal-spaces-lemma-structure-sheaf}
任意の形式代数空間は構造層をもつ。
\end{lemma}

\begin{proof}
$S$ をスキームとする。$X$ を $S$ 上の形式代数空間とする。
(\ref{formal-spaces-equation-etale-topos}) により，$\mathcal{O}_X$ を
$X_{affine, \etale}$ 上の位相環の層として構成すれば十分である。
対象が，$U$ はアフィン形式代数空間で $\varphi$ は代数空間により表現可能かつ
エタールであるような形式代数空間の射 $\varphi : U \to X$ である圏を
$\mathcal{C}$ と書く。補題 \ref{formal-spaces-lemma-identify-affine-etale} により，
関手 $U \mapsto U_{red}$ は圏の同値
$\mathcal{C} \to X_{affine, \etale}$ である。従って補題の上で与えた規則により，
既に $\mathcal{O}_X$ を $X_{affine, \etale}$ 上の位相環の前層としてもつ。
従って層条件を確認すれば十分である。

\medskip\noindent
$X_{affine, \etale}$ の定義により，被覆は $\mathcal{C}$ の射の有限族
$\{g_i : U_i \to U\}_{i = 1, \ldots, n}$ で，
$\{U_{i, red} \to U_{red}\}$ がエタール被覆であるものに対応する。
% P10-E0373
射 $g_i$ は代数空間により表現可能であり
（補題 \ref{formal-spaces-lemma-permanence-representable}），従ってアフィンである
（補題 \ref{formal-spaces-lemma-affine-representable-by-algebraic-spaces}）。
さらに $g_i$ はエタールである（『空間の性質』補題
\ref{spaces-properties-lemma-etale-permanence} から形式的に従う。
実際，$U_i$ と $U$ は『ブートストラップ』節
\ref{bootstrap-section-representable-by-spaces-properties} の意味で
$X$ 上エタールである）。
最後に，定義 \ref{formal-spaces-definition-affine-formal-algebraic-space} のように
$U = \colim U_\lambda$ と書く。

\medskip\noindent
以上の準備により，層の性質を次のように証明できる。
各 $\lambda$ に対して
$U_{i, \lambda} = U_i \times_U U_\lambda$ および
$U_{ij, \lambda} = (U_i \times_U U_j) \times_U U_\lambda$ と置く。
上の議論により，これらはアフィンスキームであり，
$\{U_{i, \lambda} \to U_\lambda\}$ はエタール被覆であり，
$U_{ij, \lambda} = U_{i, \lambda} \times_{U_\lambda} U_{j, \lambda}$ である。
また $U_i = \colim U_{i, \lambda}$ および
$U_i \times_U U_j = \colim U_{ij, \lambda}$ である。
各 $\lambda$ に対して完全列
$$
0 \to
\Gamma(U_\lambda, \mathcal{O}_{U_\lambda}) \to
\prod\nolimits_i \Gamma(U_{i, \lambda}, \mathcal{O}_{U_{i, \lambda}}) \to
\prod\nolimits_{i, j} \Gamma(U_{ij, \lambda}, \mathcal{O}_{U_{ij, \lambda}})
$$
を得る。これは $U_\lambda$ 上の構造層とエタール位相に対する層条件による
（『エタール・コホモロジー』命題
\ref{etale-cohomology-proposition-quasi-coherent-sheaf-fpqc} を参照せよ）。
極限は極限と可換なので，これらの完全列の逆極限は完全列
$$
0 \to
\lim \Gamma(U_\lambda, \mathcal{O}_{U_\lambda}) \to
\prod\nolimits_i \lim \Gamma(U_{i, \lambda}, \mathcal{O}_{U_{i, \lambda}}) \to
\prod\nolimits_{i, j} \lim
\Gamma(U_{ij, \lambda}, \mathcal{O}_{U_{ij, \lambda}})
$$
である。これは
$$
0 \to
\mathcal{O}_X(U_{red}) \to
\prod\nolimits_i \mathcal{O}_X(U_{i, red}) \to
\prod\nolimits_{i, j} \mathcal{O}_X((U_i \times_U U_j)_{red})
$$
がアーベル群の列として完全であることを意味する。なお
$\mathcal{O}_X(U_{red})$ 上の極限位相が
$\prod_{i = 1, \ldots, n} \mathcal{O}_X(U_{i, red})$
上の極限位相から誘導されることを示さなければならない。
写像 $\mathcal{O}_X(U_{red}) \to \mathcal{O}_{U_\lambda}(U_\lambda)$
の核は開部分加群の基本系をなす。これらは写像
$\prod_i \mathcal{O}_X(U_{i, red}) \to
\prod_i \mathcal{O}_{U_{i, \lambda}}(U_{i, \lambda})$
の核の逆像である。
% P10-E0374
そしてこれらの核は
$\mathcal{O}_X(U_{red}) \to \mathcal{O}_{U_\lambda}(U_\lambda)$
の開部分加群の基本系をなすので，証明は終わる。
\end{proof}

\begin{remark}
\label{formal-spaces-remark-not-enough-sections}
構造層は必ずしも「十分な切断」をもたない。『例』節
\ref{examples-section-affine-formal-algebraic-space} で，
McQuillan でないアフィン形式代数空間が存在し，正則関数が点を分離しない例さえ
存在することを見た。
\end{remark}

\noindent
次の補題では，可算添字付きアフィン形式スキームの構造層の高次コホモロジーが
消滅することを証明する。可算添字付きでないものについては，一般にはこれは
成り立たないと思われる。

\begin{lemma}
\label{formal-spaces-lemma-higher-vanishing-structure-sheaf}
$X$ が可算添字付きアフィン形式代数空間ならば，
$n > 0$ に対して $H^n(X_\etale, \mathcal{O}_X) = 0$ である。
\end{lemma}

\begin{proof}
$X_{affine, \etale}$ は同じトポスを与えるので，これを用いてよい。
消滅を示すため，『サイト上のコホモロジー』補題
\ref{sites-cohomology-lemma-cech-vanish-collection} を適用する。
$X_{affine, \etale}$ は有限非交和をもつので，これは
\par\noindent
$X_{affine, \etale} = X_{red, affine, \etale}$ における1本の射からなる被覆
$\{U_{red} \to V_{red}\}$ の {\v C}ech 複体へ帰着される
（『エタール・コホモロジー』補題
\ref{etale-cohomology-lemma-cech-complex} を参照せよ）。
従って
$$
0 \to \mathcal{O}_X(V_{red}) \to \mathcal{O}_X(U_{red}) \to
\mathcal{O}_X(U_{red} \times_{V_{red}} U_{red}) \to \ldots
$$
が完全であることを示さなければならない。以下では
$V_{red} = X_{red}$ の場合にこれを行う。一般の場合も全く同じ方法で証明される。

\medskip\noindent
$X = \text{Spf}(A)$ であり，$A$ は弱定義イデアルの可算基本系をもつ
弱許容位相環であることを思い出そう。補題
\ref{formal-spaces-lemma-affine-identify-affine-etale} および
\ref{formal-spaces-lemma-affine-etale-mcquillan} において，$X_{affine, \etale}$ の対象
$U_{red}$ は，代数空間により表現可能かつエタールなアフィン形式代数空間の射
$U \to X$ に対応し，
% P10-E0645
$B$ がエタール $A$-代数であるとき $U = \text{Spf}(B^\wedge)$ となることを見た。
構造層に対する規則により
$$
\mathcal{O}_X(U_{red}) = B^\wedge
$$
である。ここで $B^\wedge = \lim B/JB$ であり，極限は
弱定義イデアル $J \subset A$ を走ることを思い出そう。
定義をたどると
$$
\mathcal{O}_X(U_{red} \times_{X_{red}} U_{red}) = (B \otimes_A B)^\wedge
$$
を得る。以下も同様である。$U \to X$ は被覆なので，写像 $A \to B$ は忠実平坦である。
補題 \ref{formal-spaces-lemma-etale-surjective} を参照せよ。従って複体
$$
0 \to A \to B \to B \otimes_A B \to B \otimes_A B \otimes_A B \to \ldots
$$
は {\bf 普遍的に} 完全である。『降下』補題
\ref{descent-lemma-ff-exact} を参照せよ。我々の目標は
$$
H^n(0 \to A^\wedge \to B^\wedge \to (B \otimes_A B)^\wedge
\to (B \otimes_A B \otimes_A B)^\wedge \to \ldots)
$$
が $n > 0$ に対してゼロであることを示すことである。
状況を明らかにするため，完備化前の完全複体を短完全列に分割しよう。
% P10-E0375
$$
0 \to A \to B \to M_1 \to 0,\quad
0 \to M_i \to B^{\otimes_A i + 1} \to M_{i + 1} \to 0
$$
上で述べたことにより，これらは普遍的に完全な短完全列である。
従って $A$ の任意のイデアル $J$ に対して
$JM_i = M_i \cap J(B^{\otimes_A i + 1})$ である。
特に，$B^{\otimes_A i + 1}$ の部分加群としての $M_i$ 上の位相は，
$B^{\otimes_A i}$ の商加群としての $M_i$ 上の位相と同じである。
従って $A$ には弱定義イデアルの可算基本系が存在するので，補題
\ref{formal-spaces-lemma-ses} により列
$$
0 \to A^\wedge \to B^\wedge \to M_1^\wedge \to 0,\quad
0 \to M_i^\wedge \to (B^{\otimes_A i + 1})^\wedge \to M_{i + 1}^\wedge \to 0
$$
は完全のままである。これで補題が証明された。
\end{proof}

\begin{remark}
\label{formal-spaces-remark-bad-quasi-coherent}
構造層がよい性質をもっていても，準連接加群のよい理論があることを意味しない。
例えば『例』節 \ref{examples-section-nonabelian-QCoh} で，
% P10-E0646
ほとんど任意のネーター・アフィン形式代数空間に対し，準連接加群の最も自然な概念が
アーベルでない加群の圏を導くことを見た。
\end{remark}

\section{形式代数空間の余極限}
\label{formal-spaces-section-colimits-of-formal}

\noindent
本節では，節 \ref{formal-spaces-section-global-colimits} の結果を
形式代数空間の射の系の場合へ一般化する。
以下の補題における条件
「$f_{\lambda \mu} : X_\lambda \to X_\mu$ は閉埋め込みであり，
同型 $X_{\lambda, red} \to X_{\mu, red}$ を誘導する」は，
「$f_{\lambda \mu}$ は表現可能な厚化である」と言い換えられることに注意する。

\begin{lemma}
\label{formal-spaces-lemma-colimit-affine-formal}
$S$ をスキームとする。有向集合 $\Lambda$ と，
$\Lambda$ 上のアフィン形式代数空間の系
$(X_\lambda, f_{\lambda \mu})$ が与えられ，各
$f_{\lambda \mu} : X_\lambda \to X_\mu$ が閉埋め込みで，同型
$X_{\lambda, red} \to X_{\mu, red}$ を誘導すると仮定する。
このとき $X = \colim_{\lambda \in \Lambda} X_\lambda$ は
$S$ 上のアフィン形式代数空間である。
\end{lemma}

\begin{proof}
これを，有向集合 $\Omega_\lambda$ 上のアフィンスキームの余極限
$X_\lambda = \colim_{\omega \in \Omega_\lambda} X_{\lambda, \omega}$
として書ける。ここで遷移射
$X_{\lambda, \omega} \to X_{\lambda, \omega'}$ は厚化である。
各 $\lambda, \mu \in \Lambda$ と $\omega \in \Omega_\lambda$ に対し，
$\mu \geq \lambda$ ならば，射 $X_{\lambda, \omega} \to X_\mu$ が
$X_{\mu, \omega'}$ を経由するような $\omega' \in \Omega_\mu$ が存在する。
補題 \ref{formal-spaces-lemma-factor-through-thickening} を参照せよ。
すると射 $X_{\lambda, \omega} \to X_{\mu, \omega'}$ は
被約上に同型を誘導する閉埋め込みであり，従って厚化である。
$\Omega = \coprod_{\lambda \in \Lambda} \Omega_\lambda$ と置き，
$(\lambda, \omega) \leq (\mu, \omega')$ であることを，
$\lambda \leq \mu$ かつ $X_{\lambda, \omega} \to X_\mu$ が
$X_{\mu, \omega'}$ を経由することと定める。以上から
$\Omega$ は有向集合であり，
$X = \colim_{\lambda \in \Lambda} X_\lambda =
\colim_{(\lambda, \omega) \in \Omega} X_{\lambda, \omega}$ であることが従う。
これで証明が終わる。
\end{proof}

\begin{lemma}
\label{formal-spaces-lemma-colimit-formal-spaces-is-formal-space}
$S$ をスキームとする。有向集合 $\Lambda$ と，
$\Lambda$ 上の形式代数空間の系
$(X_\lambda, f_{\lambda \mu})$ が与えられ，各
$f_{\lambda \mu} : X_\lambda \to X_\mu$ が閉埋め込みで，同型
$X_{\lambda, red} \to X_{\mu, red}$ を誘導すると仮定する。
このとき $X = \colim_{\lambda \in \Lambda} X_\lambda$ は
$S$ 上の形式代数空間である。
\end{lemma}

\begin{proof}
余極限を fppf 層の圏で取るので，$X$ は層である。
$\lambda \in \Lambda$ を1つ選んで固定する。定義
\ref{formal-spaces-definition-formal-algebraic-space} のような被覆
$\{X_{i, \lambda} \to X_\lambda\}$ を選ぶ。
特に，$\{X_{i, \lambda, red} \to X_{\lambda, red}\}$ は
アフィンスキームによるエタール被覆である。
各 $\mu \geq \lambda$ に対して，エタールな垂直射をもつ
カルテジアン図式
$$
\xymatrix{
X_{i, \lambda} \ar[r] \ar[d] & X_{i, \mu} \ar[d] \\
X_\lambda \ar[r] & X_\mu
}
$$
が存在する。実際，エタール射
$X_{i, \lambda, red} \to X_{\lambda, red} = X_{\mu, red}$ は，
$X_{i, \mu}$ がアフィン形式代数空間であるような形式代数空間のエタール射
$X_{i, \mu} \to X_\mu$ に対応する。補題
\ref{formal-spaces-lemma-affine-identify-affine-etale} を参照せよ。
同じ補題により，$X_{i, \mu}$ の $X_\lambda$ への基底変換は
$X_{i, \lambda}$ と一致する。また
$\mu' \geq \mu \geq \lambda$ に対して
$X_{i, \mu} = X_\mu \times_{X_{\mu'}} X_{i, \mu'}$ であることも従う。
$X_i = \colim X_{i, \mu}$ と置く。このとき（関手として）
$X_{i, \mu} = X_i \times_X X_\mu$ である。
アフィン（または準コンパクト）スキーム $T$ からの任意の射
$T \to X = \colim X_\mu$ は，ある $\mu$ に対して $X_\mu$ に写るので，
% P10-E0376
$\colim X_{i, \mu} \to \colim X_\mu$ はエタールであると結論できる。
従って $\colim X_{i, \mu}$ がアフィン形式代数空間であることを示せば，
補題が成り立つ。射 $X_{i, \mu} \to X_{i, \mu'}$ は，閉埋め込み
$X_\mu \to X_{\mu'}$ の基底変換なので閉埋め込みであることに注意する。
最後に，$X_{\mu, red} \to X_{\mu', red}$ は同型なので，
射 $X_{i, \mu, red} \to X_{i, \mu', red}$ は同型である。
従って補題 \ref{formal-spaces-lemma-colimit-affine-formal} で論じた場合へ帰着される。
\end{proof}

\section{再完備化}
\label{formal-spaces-section-recompletion}

\noindent
本節では，形式代数空間をその被約の閉部分集合に沿って完備化することを定義する。
これは節 \ref{formal-spaces-section-completion} の自然な一般化である。

\begin{lemma}
\label{formal-spaces-lemma-completion-affine-formal-is-affine-formal}
$S$ をスキームとする。$X$ を $S$ 上のアフィン形式代数空間とする。
$T \subset |X_{red}|$ を閉部分集合とする。このとき関手
$$
X_{/T} : (\Sch/S)_{fppf} \longrightarrow \textit{Sets},\quad
U \longmapsto \{f : U \to X : f(|U|) \subset T\}
$$
はアフィン形式代数空間である。
\end{lemma}

\begin{proof}
定義 \ref{formal-spaces-definition-affine-formal-algebraic-space} のように
$X = \colim X_\lambda$ と書く。このとき
$X_{\lambda, red} = X_{red}$ であり，$T$ を
$|X_\lambda| = |X_{\lambda, red}|$ の閉部分集合とみなす。
補題 \ref{formal-spaces-lemma-completion-affine-is-affine-formal-algebraic-space} により，
各 $\lambda$ に対する完備化 $(X_\lambda)_{/T}$ はアフィン形式代数空間である。
遷移射 $(X_\lambda)_{/T} \to (X_\mu)_{/T}$ は，遷移射
$X_\lambda \to X_\mu$ の基底変換なので閉埋め込みである。
補題 \ref{formal-spaces-lemma-map-completions-representable} を参照せよ。
また補題 \ref{formal-spaces-lemma-reduction-completion} により，射
$((X_\lambda)_{/T})_{red} \to ((X_\mu)_/T)_{red}$ は同型である。
$X_{/T} = \colim (X_\lambda)_{/T}$ なので，補題
\ref{formal-spaces-lemma-colimit-affine-formal} により結論を得る。
\end{proof}

\begin{lemma}
\label{formal-spaces-lemma-completion-fas-is-fas}
$S$ をスキームとする。$X$ を $S$ 上の形式代数空間とする。
$T \subset |X_{red}|$ を閉部分集合とする。このとき関手
$$
X_{/T} : (\Sch/S)_{fppf} \longrightarrow \textit{Sets},\quad
U \longmapsto \{f : U \to X \mid f(|U|) \subset T\}
$$
は形式代数空間である。
\end{lemma}

\begin{proof}
関手 $X_{/T}$ は fppf 層である。実際，$\{U_i \to U\}$ が fppf 被覆ならば，
$\coprod |U_i| \to |U|$ は全射である。

\medskip\noindent
定義 \ref{formal-spaces-definition-formal-algebraic-space} のような被覆
$\{g_i : X_i \to X\}_{i \in I}$ を選ぶ。射
$X_i \times_X X_{/T} \to X_{/T}$ はエタールであり
（『空間』補題
\ref{spaces-lemma-base-change-representable-transformations-property} を参照せよ），
写像 $\coprod X_i \times_X X_{/T} \to X_{/T}$ は層の全射である。
従って $X_{/T} \times_X X_i$ がアフィン形式代数空間であることを示せば十分である。
$X_i \times_X X_{/T}$ の $U$-値点は，像が閉部分集合
$|g_{i, red}|^{-1}(T) \subset |X_{i, red}|$ に含まれる射
$U \to X_i$ である。従って補題
\ref{formal-spaces-lemma-completion-affine-formal-is-affine-formal} から従う。
\end{proof}

\begin{definition}
\label{formal-spaces-definition-completion-formal-algebraic-space}
$S$ をスキームとする。$X$ を $S$ 上の形式代数空間とする。
$T \subset |X_{red}|$ を閉部分集合とする。
% P10-E0647
補題 \ref{formal-spaces-lemma-completion-is-formal-algebraic-space} の形式代数空間
$X_{/T}$ を，{\it $T$ に沿う $X$ の完備化} と呼ぶ。
\end{definition}

\noindent
$f : X \to X'$ をスキーム $S$ 上の形式代数空間の射とする。
$T \subset |X_{red}|$ および $T' \subset |X'_{red}|$ を，
$|f_{red}|(T) \subset T'$ を満たす閉部分集合とする。
すると $f$ は完備化の間の形式代数空間の射
$$
X_{/T} \longrightarrow X'_{/T'}
$$
を定めることは明らかである。

\begin{lemma}
\label{formal-spaces-lemma-map-recompletions}
$S$ をスキームとする。$f : X' \to X$ を $S$ 上の形式代数空間の射とする。
$T \subset |X_{red}|$ を閉部分集合とし，
$T' = |f_{red}|^{-1}(T) \subset |X'_{red}|$ と置く。このとき
$$
\xymatrix{
X'_{/T'} \ar[r] \ar[d] & X' \ar[d]^f \\
X_{/T} \ar[r] & X
}
$$
は $S$ 上の形式代数空間のカルテジアン図式である。
\end{lemma}

\begin{proof}
実際，構成により水平射はモノ射である。従って，スキーム $U$ からの射
% P10-E0377
$g : U \to X'$ が $X'_{/T}$ の点を定めることと，
$f \circ g$ が $X_{/T}$ の点を定めることが同値であると示せば十分である。
言い換えると，$g(U)$ が $T' \subset |X'_{red}|$ に含まれることと，
$(f \circ g)(U)$ が $T \subset |X_{red}|$ に含まれることが同値であると
示さなければならない。これは $T'$ を $T$ の逆像として選んだことから直ちに従う。
\end{proof}

\begin{lemma}
\label{formal-spaces-lemma-reduction-recompletion}
$S$ をスキームとする。$X$ を $S$ 上の形式代数空間とする。
$T \subset |X_{red}|$ を閉部分集合とする。
$T$ に沿う $X$ の完備化 $X_{/T}$ の被約 $(X_{/T})_{red}$ は，
$T$ に対応する $X_{red}$ の被約被誘導閉部分空間 $Z$ である。
\end{lemma}

\begin{proof}
補題 \ref{formal-spaces-lemma-reduction-formal-algebraic-space}，
『空間の性質』定義
\ref{spaces-properties-definition-reduced-induced-space}
（これは $Z$ を構成するため『空間の性質』補題
\ref{spaces-properties-lemma-reduced-closed-subspace} を用いる），
および $X_{/T}$ の定義から，$Z$ と $(X_{/T})_{red}$ は
% P10-E0378
同じ写像性質によって特徴付けられる被約代数空間であることが従う。
すなわち，始域が被約代数空間である射 $g : Y \to X$ は，
$|Y|$ が $T \subset |X|$ に写るとき，かつそのときに限り，これらを経由する。
\end{proof}

\begin{lemma}
\label{formal-spaces-lemma-recompletion-affine-types}
$S$ をスキームとする。$X$ を $S$ 上のアフィン形式代数空間とする。
$T \subset X_{red}$ を閉部分集合とし，$X_{/T}$ を
$T$ に沿う $X$ の形式完備化とする。このとき
\begin{enumerate}
\item $X_{/T}$ はアフィン形式代数空間である。
\item $X$ が McQuillan ならば，$X_{/T}$ は McQuillan である。
\item $|X_{red}| \setminus T$ が準コンパクトで $X$ が可算添字付きならば，
$X_{/T}$ は可算添字付きである。
\item $|X_{red}| \setminus T$ が準コンパクトで $X$ が adic* ならば，
$X_{/T}$ は adic* である。
\item $X$ が Noether ならば，$X_{/T}$ は Noether である。
\end{enumerate}
\end{lemma}

\begin{proof}
(1) は補題 \ref{formal-spaces-lemma-completion-affine-formal-is-affine-formal} である。
$X$ が McQuillan ならば，ある弱許容位相環 $A$ に対して
$X = \text{Spf}(A)$ である。このとき $X_{/T} \to X \to \Spec(A)$ は
補題 \ref{formal-spaces-lemma-mcquillan-affine-formal-algebraic-space} の性質 (2) を満たすので，
$X_{/T}$ は McQuillan である。定義
\ref{formal-spaces-definition-types-affine-formal-algebraic-space} を参照せよ。

\medskip\noindent
$X$ と $T$ が (3) のとおりであると仮定する。このとき
$X = \text{Spf}(A)$ であり，$A$ は弱定義イデアルの基本系
$A \supset I_1 \supset I_2 \supset I_3 \supset \ldots$
をもつ。補題 \ref{formal-spaces-lemma-countably-indexed} を参照せよ。
『代数』補題 \ref{algebra-lemma-qc-open} により，有限生成イデアル
$\overline{J} = (\overline{f}_1, \ldots, \overline{f}_r) \subset A/I_1$
で，$\Spec(A/I_1) = |X_{red}|$ の内部で $T$ を $\overline{J}$ により
切り出すものを見つけられる。
$\overline{f}_i$ の持ち上げ $f_i \in A$ を選ぶ。
$Z = \Spec(B)$ をアフィンスキームとし，$g : Z \to X$ を
$g(Z) \subset T$ を満たす射とする（集合論的に）。
すると $g^\sharp : A \to B$ はある $n$ に対して $A/I_n$ を経由し，
各 $i$ に対して $g^\sharp(f_i)$ は $B$ の冪零元である。
従ってある $n, m \geq 1$ に対して
$J_{m, n} = (f_1, \ldots, f_r)^m + I_n$ は $B$ でゼロに写る。
従って $X_{/T}$ は $\lim_{n, m} A/J_{m, n}$ の形式スペクトルであり，
可算添字付きである。(3) が証明された。

\medskip\noindent
(4) の証明。この場合も議論は (3) と同じである。
ただし，ある有限生成イデアル $I \subset A$ に対して
$I_n = I^n$ と選べる。このとき $J = (f_1, \ldots, f_r) + I$ と置けば，
$X_{/T}$ が $\lim A/J^n$ の形式スペクトルであることは明らかである。
いくつかの詳細は省略する。

\medskip\noindent
(5) の証明。この場合 $X_{red}$ はネーター環のスペクトルなので，
$|X_{red}| \setminus T$ が準コンパクトであるという仮定は満たされる。
従って (4) の証明と同様に，
% P10-E0648
$X_{/T}$ は $\lim A/J^n$ のスペクトルであり，後者はネーター adic 位相環である。
『代数』補題 \ref{algebra-lemma-completion-Noetherian-Noetherian} を参照せよ。
\end{proof}

\begin{lemma}
\label{formal-spaces-lemma-recompletion-types}
$S$ をスキームとする。$X$ を $S$ 上の形式代数空間とする。
$T \subset X_{red}$ を閉部分集合とし，$X_{/T}$ を
$T$ に沿う $X$ の形式完備化とする。このとき
\begin{enumerate}
\item $X_{red} \setminus T \to X_{red}$ が準コンパクトで，$X$ が
局所可算添字付きならば，$X_{/T}$ は局所可算添字付きである。
\item $X_{red} \setminus T \to X_{red}$ が準コンパクトで，$X$ が
局所 adic* ならば，$X_{/T}$ は局所 adic* である。
\item $X$ が局所 Noether ならば，$X_{/T}$ は局所 Noether である。
\end{enumerate}
\end{lemma}

\begin{proof}
定義 \ref{formal-spaces-definition-formal-algebraic-space} のような被覆
$\{X_i \to X\}$ を選ぶ。$T$ の逆像を $T_i \subset X_{i, red}$ とする。
$X_i \times_X X_{/T} = (X_i)_{/T_i}$ である
（補題 \ref{formal-spaces-lemma-map-recompletions}）。
従って $\{(X_i)_{/T_i} \to X_{/T}\}$ は定義
\ref{formal-spaces-definition-formal-algebraic-space} のような被覆である。
さらに，$X_{red} \setminus T \to X_{red}$ が準コンパクトならば，
$X_{i, red} \setminus T_i \to X_{i, red}$ も準コンパクトである。
% P10-E0379
また $X$ が局所可算添字付き，局所 adic*，または局所 Noether ならば，
$X_i$ はそれぞれ可算添字付き，adic*，または Noether である。
従ってアフィンの場合，すなわち補題
\ref{formal-spaces-lemma-recompletion-affine-types} から補題が従う。
\end{proof}







\section{閉部分空間に沿う完備化}
\label{formal-spaces-section-completion-subspace}

\noindent
% P10-E0380
本節は，閉部分空間に関する完備化についての節
\ref{formal-spaces-section-completion} の類似である。

\begin{definition}
\label{formal-spaces-definition-completion-subspace}
$S$ をスキームとする。$X$ を $S$ 上の代数空間とする。
$Z \subset X$ を閉部分空間とし，$n$ 次無限小近傍を
$Z_n \subset X$ と書く。形式代数空間
$$
X^\wedge_Z = \colim Z_n
$$
（補題 \ref{formal-spaces-lemma-colimit-formal-spaces-is-formal-space} を参照せよ）を
{\it $Z$ に沿う $X$ の完備化}と呼ぶ。
\end{definition}

\noindent
例えば $X = \Spec(A)$ で，$Z$ がイデアル $I \subset A$ により
切り出されるならば，
$$
X^\wedge_Z = \colim \Spec(A/I^n) = \text{Spf}(B)
$$
である。ここで $B = \lim A/I^n$ は極限位相を備える。
$B$ は弱 adic 位相環ではあるが，一般には adic ではないことに注意せよ。

\medskip\noindent
一般の場合に戻る。$T = |Z|$ と置くと，$Z$ に沿う完備化と
$T$ に沿う完備化（節 \ref{formal-spaces-section-completion}）を比較する標準射
$X^\wedge_Z \to X_{/T}$ が存在するが，これは同型とは限らない
（補題 \ref{formal-spaces-lemma-when-isomorphism} を参照せよ）。

\medskip\noindent
$f : X \to X'$ をスキーム $S$ 上の代数空間の射とする。
$Z \subset X$ および $Z' \subset X'$ を閉部分空間とし，
$f|_Z$ が $Z$ を $Z'$ に写して射 $Z \to Z'$ を誘導すると仮定する。
すると $f$ が完備化の間の形式代数空間の射
$$
X^\wedge_Z \longrightarrow (X')^\wedge_{Z'}
$$
を定めることは明らかである。

\begin{lemma}
\label{formal-spaces-lemma-map-completions-subspaces-representable}
$S$ をスキームとする。$f : X' \to X$ を $S$ 上の代数空間の射とする。
$Z \subset X$ を閉部分空間とし，
$Z' = f^{-1}(Z) = X' \times_X Z$ と置く。このとき
$$
\xymatrix{
(X')^\wedge_{Z'} \ar[r] \ar[d] & X' \ar[d]^f \\
X^\wedge_Z \ar[r] & X
}
$$
は層のカルテジアン図式である。特に射
$(X')^\wedge_{Z'} \to X^\wedge_Z$ は代数空間により表現可能である。
\end{lemma}

\begin{proof}
実際，$Y \to X$ をスキームから $X$ への射で，射 $Y \to X$ が
$Z$ を経由するものとする。
すると $Y \times_X X' \to X$ は代数空間の射であり，
$Y \times_X X' \to X'$ は $Z'$ を経由する。
すべての $n \geq 1$ に対して $Z'_n = X' \times_X Z_n$ なので，
無限小近傍についても同じことが成り立つ。
従って関手のカルテジアン正方形は，公式
$X^\wedge_Z = \colim Z_n$ および
$(X')^\wedge_{Z'} = \colim Z'_n$ から従う。
\end{proof}

\begin{lemma}
\label{formal-spaces-lemma-reduction-completion-subspace}
$S$ をスキームとする。$X$ を $S$ 上の代数空間とする。
$Z \subset X$ を閉部分空間とする。$Z$ に沿う $X$ の完備化
$X^\wedge_Z$ の被約 $(X^\wedge_Z)_{red}$ は $Z_{red}$ である。
\end{lemma}

\begin{proof}
省略する。
\end{proof}

\begin{lemma}
\label{formal-spaces-lemma-affine-formal-completion-subspace-types}
$S$ をスキームとする。$X = \Spec(A)$ を $S$ 上のアフィンスキームとする。
$Z \subset X$ を，イデアル $I \subset A$ に対応する閉部分スキームとする。
このとき
\begin{enumerate}
\item アフィン形式代数空間 $X^\wedge_Z$ は弱 adic である。
\item $I$ が有限生成ならば，
$X^\wedge_Z = \text{Spf}(A^\wedge)$ である。ここで $A^\wedge$ は
$A$ の $I$-進完備化である。
\item $Z \to X$ が有限表示，すなわち $I$ が有限生成ならば，
$X^\wedge_Z$ は adic* である。
\item $X$ が Noether ならば，$X^\wedge_Z$ は Noether である。
\end{enumerate}
\end{lemma}

\begin{proof}
省略する。
\end{proof}

\begin{lemma}
\label{formal-spaces-lemma-formal-completion-subspace-types}
$S$ をスキームとする。$X$ を $S$ 上の代数空間とする。
$Z \subset X$ を閉部分空間とする。$X^\wedge_Z$ を
$Z$ に沿う $X$ の形式完備化とする。
\begin{enumerate}
\item 形式代数空間 $X^\wedge_Z$ は局所弱 adic である。
\item $Z \to X$ が有限表示ならば，$X^\wedge_Z$ は局所 adic* である。
% P10-E0381
\item $X$ が局所 Noether ならば，$X_Z$ は局所 Noether である。
\end{enumerate}
\end{lemma}

\begin{proof}
省略する。
\end{proof}

\begin{lemma}
\label{formal-spaces-lemma-when-isomorphism}
$S$ をスキームとする。$X$ を $S$ 上の代数空間とする。
$Z \subset X$ を閉部分空間とし，$T = |Z|$ と置く。標準射
$c : X^\wedge_Z \to X_{/T}$ は，$Z \to X$ が有限表示ならば同型であるが，
一般には同型ではない。
\end{lemma}

\begin{proof}
両方の構成はエタール局所化と可換なので，$X$ がアフィンの場合を
証明すれば十分である。$X = \Spec(A)$ と書き，$Z$ がイデアル
$I \subset A$ に対応するとする。$I$ が有限生成ならば，
$X^\wedge_Z$ と $X_{/T}$ はともに $\text{Spf}(A^\wedge)$ に等しい。
ここで $A^\wedge$ は $A$ の $I$-進完備化である。補題
\ref{formal-spaces-lemma-affine-formal-completion-types} および
\ref{formal-spaces-lemma-affine-formal-completion-subspace-types} を参照せよ。
$A = \mathbf{Z}[x_1, x_2, \ldots]$ かつ $I = (x_1, x_2, \ldots)$ ならば，
$\Spec(A/(x_n^n, n \geq 1)) \to X$ は $X_{/T}$ を経由するが
$X^\wedge_Z$ を経由しない。従って $c$ は同型ではない。
\end{proof}
